\documentclass[11pt]{amsart}

\usepackage[utf8]{inputenc}
\usepackage[T1]{fontenc}
\usepackage{lmodern}
\usepackage[protrusion=true,expansion=false]{microtype}
\usepackage{amsmath,amssymb,amsthm}
\usepackage{mathtools}
\usepackage{mathrsfs}
\usepackage{stackrel}
\usepackage{graphicx}
\usepackage{booktabs}
\usepackage{tabularray}

\hfuzz=2pt
\hbadness=5000
\vbadness=5000
\emergencystretch=2em
\raggedbottom
\allowdisplaybreaks[2]

\newtheorem{example}{Example}[section]

\DeclareMathOperator{\Sing}{Sing}
\DeclareMathOperator{\CritVal}{CritVal}
\DeclareMathOperator{\CritLoc}{CritLoc}

\newcommand{\MorseCat}[1]{\mathsf{Q}_{#1}}
\newcommand{\MBcat}[1]{\mathsf{MB}(#1)}
\newcommand{\MBtwcat}[2]{\mathsf{MB}(#1;[#2])}
\newcommand{\Aper}[2]{\mathsf{A}^{\mathrm{per}}_{#2}(#1)}
\newcommand{\QAper}[2]{\mathsf{QA}^{\mathrm{per}}_{#2}(#1)}
\newcommand{\QAtwist}[3]{\mathsf{QA}^{\mathrm{per}}_{#2}(#1;[#3])}
\newcommand{\Dfourcat}[2]{\mathsf{D}^{\mathrm{per}}_{4}(#1;[#2])}
\newcommand{\fitdisplay}[1]{%
  \begin{center}%
  \resizebox{\textwidth}{!}{$\displaystyle #1$}%
  \end{center}%
}

% Numbered references below point to the corresponding results in the main
% article. They are declared locally so that this companion appendix compiles
% without requiring the main article's auxiliary file.
\makeatletter
\@namedef{r@thm:general_formal}{{3.10}{1}}
\@namedef{r@prop:Ade_category_model}{{4.12}{1}}
\@namedef{r@prop:morse_bott_clifford}{{4.18}{1}}
\@namedef{r@sec:mather_yau}{{4}{1}}
\@namedef{r@sec:fano_threefold_decompositions}{{6}{1}}
\makeatother

\begin{document}

\title{Detailed Computations for the Fano Threefold Models}
\author{Victor Ibrahim Santos El Adji}
\date{}
\maketitle

\section{Conventions and Computational Scope}
\label{sec:fano_threefold_appendix}
This appendix supplies the detailed calculations supporting the inventory in
Section~\ref{sec:fano_threefold_decompositions}. We use the
Minkowski-polynomial, Graded Ring Database, and Iskovskikh--Mori--Mukai
classifications
\cite{akhtar2012minkowski,Coates_2019,Fano3,Iskovskikh1977FanoI,Iskovskikh1978FanoII,MoriMukai1981FanoB2,MoriMukai1983FanoB2,MoriMukai2003Erratum}.
Every decomposition below is interpreted in the ungraded two-periodic
convention of Section~\ref{sec:mather_yau}. Isolated Morse points contribute
$\MorseCat{3}$. A smooth critical component is written as $\MBcat{Z}$ only
when a split relative quadratic formal model has been exhibited. Likewise,
$\Aper{Z}{m}$ is used only after verifying a globally trivial formal normal
coordinate and a global contact equivalence to $u^{m+1}$. Direct sums are
asserted only for critical values or closed critical components that are
genuinely disjoint. The calculations concern the affine torus; possible
vanishing cycles at infinity require a compactification and are not included.
When the discriminant of a verified quadratic bundle is not a square in
$\mathcal O_Z^*$, the notation $\MBtwcat{Z}{\delta}$ records its nontrivial
square class explicitly. Thus $\MBcat{Z}$ is reserved for the untwisted case.
Following the notation of Sections~5 and~6 of the reference article, a mixed
formal normal form $v^2+u^{m+1}$ contributes $\QAper{Z}{m}$, whereas
$v^2+\delta u^{m+1}$ contributes $\QAtwist{Z}{m}{\delta}$ when the square
class $[\delta]\in\mathcal O_Z^*/\mathcal O_Z^{*2}$ is nontrivial. The
twisted residual $D_4$ model is denoted by $\Dfourcat{Z}{\delta}$.

The Mori--Mukai examples in this section are accompanied by the ancillary certificate file \texttt{supplementary-morimukai-certificates.tex}. It is regenerated, and its underlying symbolic checks are rerun, by the commands
\begin{center}
\begin{tabular}{l}
\texttt{python computations/build\_morimukai\_certificates.py}\\
\texttt{python computations/build\_morimukai\_correspondence.py}\\
\texttt{python computations/mori\_mukai\_validation.py --all}\\
\texttt{python computations/rank\_one\_fano\_validation.py --all}\\
\texttt{python computations/rank\_one\_five\_chart\_gluing.py}\\
\texttt{python computations/morse\_bott\_audit.py}.
\end{tabular}
\end{center}
The Mori--Mukai validator covers the specified models for the eighty-eight
Mori--Mukai families, and the rank-one validator covers the specified models
for the seventeen Picard-rank-one families. Together they give executable
checks for one specified Laurent-polynomial model for each of the $105$
deformation families in the summary table. The Morse--Bott audit independently checks the retained smooth
components and crossing strata in selected models.  The generated certificate
supplement records, label by label for the eighty-eight Mori--Mukai cases, the
cleared logarithmic critical equations, saturations by the torus monomial,
point and Hessian checks, algebraic branches, and the normal-form or
formal-expansion checks used in the decompositions.  The rank-one validator
performs the corresponding exact checks directly from the compact and expanded
potentials displayed below.  The five-chart validator checks the projective
chart calculations and the overlap identifications for the compactified
$B_4$ and $V_{18}$ pencils discussed below.

\begin{example}
    Consider the LG model for $\mathbb{P}^3$ given by
    \[
    \left((\mathbb{C}^*)^3,x+y+z+\frac{1}{xyz}=W_{\mathbb{P}^3}\right).
    \]
    The critical locus is
    \[
    \CritLoc(W_{\mathbb{P}^3})=\{(\omega_4^i,\omega_4^i,\omega_4^i)\mid i=0,1,2,3\},
    \]
    where $\omega_4=e^{2\pi i/4}$, and the corresponding critical values are
    \[
    \CritVal(W_{\mathbb{P}^3})=\{4\omega_4^i\mid i=0,1,2,3\}.
    \]
    Let $Z=V(x-y,y-z)$. Observe that $\CritLoc(W_{\mathbb{P}^3})\subset Z$. The rank-two normal bundle is trivial and has nondegenerate quadratic part, so Theorem~\ref{thm:general_formal}, together with the local quadratic normal form, gives
    \[
    \overline{MF((\mathbb{C}^*)^3,W_{\mathbb{P}^3})}\cong \overline{MF^Z((\mathbb{C}^*)^3,W_{\mathbb{P}^3})}\cong\overline{MF\left((\mathbb{C}^*),3x+\frac{1}{x^3}\right)}
    \]
    The contact equivalence $x+\frac{1}{x}-2\overset{c}{\sim} 3x+\frac{1}{x^3}-4\omega_{4}^i$ implies
    \[
    \overline{MF_{4\omega_4^i}\left((\mathbb{C}^*),3x+\frac{1}{x^3}\right)}\cong \overline{MF_2\left((\mathbb{C}^*),x+\frac{1}{x}\right)}
    \]

The one-dimensional pair
$\left((\mathbb{C}^*),x+\frac{1}{x}\right)$ is the standard LG model for
$\mathbb{P}^1$, and its contributions at $2$ and $-2$ are equivalent. Hence
\[
\overline{MF((\mathbb{C}^*)^3,W_{\mathbb{P}^3})}
\cong \bigoplus_{n=1}^{4}\MorseCat{3}.
\]

\end{example}
\begin{example}
    Consider the LG model for $Q\subset\mathbb{P}^4$ given by
    \[
    \left((\mathbb{C}^*)^3,x+y+z+\frac{1}{xy}+\frac{1}{xz}=W_{Q}\right).
    \]
    The critical locus is
    \[
    \CritLoc(W_Q)=\left\{\left(\sqrt[3]{4}\,\omega_3^i,\frac{\omega_3^i}{\sqrt[3]{2}},\frac{\omega_3^i}{\sqrt[3]{2}}\right)\ \middle|\ i=0,1,2\right\},
    \]
    where $\omega_3=e^{2\pi i/3}$, and the corresponding critical values are
    \[
    \CritVal(W_Q)=\{3\sqrt[3]{4}\,\omega_3^i\mid i=0,1,2\}.
    \]
    Let $Z=V(y-z)$. Observe that $\CritLoc(W_Q)\subset Z$. The normal line to $Z$ is trivial and the transverse quadratic part is nondegenerate, so Theorem~\ref{thm:general_formal}, together with the local quadratic normal form, gives
    \[
    \overline{MF((\mathbb{C}^*)^3,W_{Q})}\cong \overline{MF^Z((\mathbb{C}^*)^3,W_{Q})}\cong\overline{MF\left((\mathbb{C}^*)^2,x+2y+\frac{2}{xy}\right)}[1]
    \]
    The contact equivalence $x+y+\frac{1}{xy}-3\omega_3^i\overset{c}{\sim}x+2y+\frac{2}{xy}-3\sqrt[3]{4}\omega_3^i$ implies
    \[
    \overline{MF_{3\sqrt[3]{4}\omega_3^i}\left((\mathbb{C}^*)^3,W_{Q}\right)}\cong \overline{MF_{3\omega_3^i}\left((\mathbb{C}^*)^2,x+y+\frac{1}{xy}\right)}[1]
    \]

The two-dimensional pair
$\left((\mathbb{C}^*)^2,x+y+\frac{1}{xy}\right)$ is the standard LG model
for $\mathbb{P}^2$, and
\[
\overline{MF_{3\omega_3^{i}}\left((\mathbb{C}^*)^2,x+y+\frac{1}{xy}\right)}
\cong
\overline{MF_{3}\left((\mathbb{C}^*)^2,x+y+\frac{1}{xy}\right)}.
\]
Hence
\[
\overline{MF((\mathbb{C}^*)^3,W_{Q})}
\cong \bigoplus_{n=1}^{3}\MorseCat{3}.
\]

\end{example}

\begin{example}[The Landau--Ginzburg model for $B_5$]
Appendix~A of \cite{CoatesCortiGalkinKasprzyk2016} gives the $7$-term
Laurent polynomial
\[
W_{B_5}=x+y+z+\frac1x+\frac1y+\frac1z+\frac1{xyz}.
\]
Its classical period begins
\[
\sum_{d\geq0}\mathrm{CT}(W_{B_5}^d)t^d
=1+6t^2+114t^4+2940t^6+87570t^8+\cdots,
\]
and agrees through order $t^{10}$ with the regularized quantum period in
Section~H.7 of the same article.  The saturated logarithmic critical ideal is
\[
(x-y,\ y-z,\ z^4-z^2-1).
\]
Hence
\[
\CritLoc(W_{B_5})
=\{(a,a,a):a^4-a^2-1=0\}.
\]
At such a point,
\[
W_{B_5}(a,a,a)=3a+\frac3a+\frac1{a^3},
\]
and
\[
\det(\theta_i\theta_jW_{B_5})(a,a,a)
=\frac{(a^2+1)^2(a^4+a^2+3)}{a^5}.
\]
The numerator is coprime to $a^4-a^2-1$, so all four points are Morse.
Consequently,
\[
\overline{MF((\mathbb C^*)^3,W_{B_5})}
\simeq\bigoplus_{a:\,a^4-a^2-1=0}\MorseCat{3}.
\]
\end{example}

\begin{example}[The Landau--Ginzburg model for $B_4$]
The source polynomial in Appendix~A of
\cite{CoatesCortiGalkinKasprzyk2016} is
\[
W_{B_4}
=x+yz^2+2yz+y+\frac2{yz}+\frac1{xy^2z^2}.
\]
Its classical period begins
\[
\sum_{d\geq0}\mathrm{CT}(W_{B_4}^d)t^d
=1+8t^2+216t^4+8000t^6+343000t^8+\cdots,
\]
and agrees through order $t^{10}$ with the regularized quantum period in
Section~H.6 of the same article.  Saturation of the logarithmic critical ideal
gives
\[
I_{\mathrm{crit}}
=(xy-1,z+1)\cap(x-y,y^2-1,z-1).
\]
Thus the reduced critical locus is the disjoint union of the smooth curve
\[
\Gamma=\{(t,t^{-1},-1):t\in\mathbb C^*\}
\]
and the two points
\[
P_+=(1,1,1),
\qquad
P_-=(-1,-1,1).
\]
The critical values are $0,8,-8$, respectively, and
\[
\det(\theta_i\theta_jW_{B_4})(P_+)=32,
\qquad
\det(\theta_i\theta_jW_{B_4})(P_-)=-32.
\]
Hence both isolated points are Morse.

Along $\Gamma$, put
\[
t=x,
\qquad
U=xy-1,
\qquad
V=z+1.
\]
Then $y=(1+U)/t$, $z=V-1$, and the transverse quadratic part is
\[
W_{B_4}
\equiv
t(U-V)^2+\frac1tV^2
\pmod{(U,V)^3}.
\]
It splits globally as
\[
t\left(U-\left(1-\frac{\mathrm i}{t}\right)V\right)
 \left(U-\left(1+\frac{\mathrm i}{t}\right)V\right),
\]
and the corresponding linear change has determinant $-2\mathrm i/t$, a unit
on $\Gamma$.  The relative formal Morse lemma therefore gives an untwisted
Morse--Bott model along the whole curve. Consequently,
\[
\overline{MF((\mathbb C^*)^3,W_{B_4})}
\simeq
\MBcat{\Gamma}\oplus\MorseCat{3}\oplus\MorseCat{3}.
\]
\end{example}

\begin{example}[The Landau--Ginzburg model for $B_3$]
The $9$-term source polynomial is
\[
W_{B_3}=x+\frac{x}{yz}+y+z+\frac2z+\frac2y
+\frac{y}{xz}+\frac2x+\frac{z}{xy}
\]
\cite{CoatesCortiGalkinKasprzyk2016}. Its classical period begins
\[
\sum_{d\geq0}\mathrm{CT}(W_{B_3}^d)t^d
=1+12t^2+540t^4+33600t^6+2425500t^8+\cdots,
\]
and agrees through order $t^{10}$ with the regularized quantum period in
Section~H.5 of the same article. The saturated logarithmic critical ideal is
\[
(x-y,\ y-z,\ z^2-3).
\]
Thus the critical points are $(a,a,a)$ with $a^2=3$. At either point,
\[
W_{B_3}(a,a,a)=6a,
\qquad
\det(\theta_i\theta_jW_{B_3})(a,a,a)=54a.
\]
Both points are Morse, and hence
\[
\overline{MF((\mathbb C^*)^3,W_{B_3})}
\simeq\bigoplus_{a^2=3}\MorseCat{3}.
\]
\end{example}

\begin{example}[The Landau--Ginzburg model for $B_2$]
The $10$-term Laurent polynomial of Appendix~A of
\cite{CoatesCortiGalkinKasprzyk2016} can be written compactly as
\[
W_{B_2}
=x(y+z+1)^2+\frac{2(y+z+1)}{xyz}+\frac1{x^3y^2z^2}.
\]
Its classical period begins
\[
\sum_{d\geq0}\mathrm{CT}(W_{B_2}^d)t^d
=1+24t^2+2520t^4+369600t^6+63063000t^8+\cdots,
\]
and agrees through order $t^{10}$ with the regularized quantum period in
Section~H.4 of the same article. Set
\[
U=y+z+1+\frac1{x^2yz}.
\]
Then the exact identity
\[
W_{B_2}=xU^2
\]
shows that the smooth surface $S=V(U)$ is critical with value $0$. It is
smooth because
\[
\frac{\partial U}{\partial x}=-\frac2{x^3yz}
\]
is a unit on the torus. Away from $S$, saturation gives the complementary
critical ideal
\[
(x^2-1,\ y-1,\ z-1).
\]
Thus the remaining critical points are $(1,1,1)$ and $(-1,1,1)$, with
critical values $16$ and $-16$ and logarithmic Hessian determinants $512$ and
$-512$, respectively. Both are Morse.

The function $U$ is a global formal normal coordinate along $S$, but the
coefficient $x$ has nontrivial square class in $\mathcal O_S^*$. Indeed, the
curve
\[
C=\{(t,-1,t^{-1}):t\in\mathbb C^*\}\subset S
\]
pulls $x$ back to $t$, which is not a square in $\mathbb C[t^{\pm1}]$.
Therefore
\[
\overline{MF_0^S((\mathbb C^*)^3,W_{B_2})}
\simeq\MBtwcat{S}{x}.
\]
Since the three critical components are disjoint, we obtain
\[
\overline{MF((\mathbb C^*)^3,W_{B_2})}
\simeq
\MBtwcat{S}{x}\oplus\MorseCat{3}\oplus\MorseCat{3}.
\]
\end{example}

\begin{example}[The Landau--Ginzburg model for $B_1$]
The $18$-term weighted-projective complete-intersection mirror listed in
Appendix~A of \cite{CoatesCortiGalkinKasprzyk2016} has the compact form
\[
W_{B_1}
=(x+y)\left((z+1)^2+\frac{1}{xyz^4}\right)^2.
\]
Equivalently, its expansion is
\[
\begin{aligned}
W_{B_1}={}&xz^4+4xz^3+6xz^2+4xz+x
 +yz^4+4yz^3+6yz^2+4yz+y\\
&+\frac{2}{yz^2}+\frac{4}{yz^3}+\frac{2}{yz^4}
 +\frac{2}{xz^2}+\frac{4}{xz^3}+\frac{2}{xz^4}
 +\frac{1}{xy^2z^8}+\frac{1}{x^2yz^8}.
\end{aligned}
\]
Its classical period begins
\[
\sum_{d\geq0}\mathrm{CT}(W_{B_1}^d)t^d
=1+120t^2+83160t^4+81681600t^6+93699005400t^8+\cdots,
\]
and agrees through order $t^{10}$ with the regularized quantum period in
Section~H.3 of the same article.

Set
\[
A=x+y,
\qquad
G=(z+1)^2+\frac{1}{xyz^4},
\qquad
q=\frac{1}{xyz^4}.
\]
Thus $W_{B_1}=AG^2$, and
\[
\theta_xW_{B_1}=G(xG-2Aq),
\qquad
\theta_yW_{B_1}=G(yG-2Aq),
\]
\[
\theta_zW_{B_1}=4AG\bigl(z(z+1)-2q\bigr).
\]
The surface
\[
S=V(G)
\]
is therefore critical with value $0$.  It is smooth because
$\partial_xG=-1/(x^2yz^4)$ is a unit.  Outside $S$, subtracting the first two
equations gives $x=y$, and the remaining equations reduce to
\[
G=4q,
\qquad
z(z+1)=2q.
\]
They give $z=2$ and $48x^2=1$.  Thus the isolated critical points are
\[
P_a=(a,a,2),
\qquad
48a^2-1=0,
\]
with
\[
W_{B_1}(P_a)=288a,
\qquad
\det(\theta_i\theta_jW_{B_1})(P_a)=124416a.
\]
Both are Morse.

The quadratic surface contribution changes type along
\[
K=V(A,G)\subset S.
\]
Equivalently, $K$ is the disjoint union of the two smooth curves
\[
xz^2(z+1)=1,\quad y=-x,
\qquad\text{and}\qquad
xz^2(z+1)=-1,\quad y=-x.
\]
On $S^\circ=S\setminus K$, the coefficient $A$ is a unit, but its square class
is nontrivial. To see this, restrict to $z=1$, where
\[
y=-\frac1{4x},
\qquad
A=\frac{4x^2-1}{4x}
\]
in $\mathbb C[x^{\pm1},(4x^2-1)^{-1}]$; the factors $x$, $2x-1$, and
$2x+1$ occur with odd valuation. Thus the exact factorization $W_{B_1}=AG^2$
gives the twisted contribution $\MBtwcat{S^\circ}{A}$. Along $K$, the
functions $(A,G)$ are global formal normal coordinates and the transverse
model is the nonsplit crossing $uv^2$.
Define
\[
\mathcal C_{B_1,0}:=
\overline{MF_0((\mathbb C^*)^3,W_{B_1})}
\]
and
\[
\mathcal P_{B_1}:=
\overline{MF^t_0\!\left(
\operatorname{Spf}(\mathcal O_K[[A,G]]),AG^2
\right)}.
\]
Supported localization gives the Verdier sequence
\[
\mathcal P_{B_1}\longrightarrow
\mathcal C_{B_1,0}\longrightarrow
\MBtwcat{S^\circ}{A},
\]
which is not asserted to split.  Therefore
\[
\overline{MF((\mathbb C^*)^3,W_{B_1})}
\cong
\mathcal C_{B_1,0}
\oplus
\bigoplus_{a:\,48a^2-1=0}\MorseCat{3}.
\]
\end{example}


\medskip
\noindent\textbf{Why there is no $V_{20}$.}
For a prime Fano threefold of Picard rank one and index one, the label $V_d$
records the anticanonical degree $d=(-K_X)^3=2g-2$.  Thus $V_{20}$ would
have genus $11$, which does not occur in the Iskovskikh classification.
Accordingly, the rank-one list jumps directly from $V_{18}$ to $V_{22}$
\cite{Iskovskikh1977FanoI,Iskovskikh1978FanoII,
CoatesCortiGalkinKasprzyk2016}.

\begin{example}[The Landau--Ginzburg model for $V_{22}$]
Appendix~A of \cite{CoatesCortiGalkinKasprzyk2016} gives the following
$11$-term Laurent polynomial:
\[
W_{V_{22}}=xy+\frac{xy}{z}+x+y+\frac{2y}{z}+z
+\frac{2z}{y}+\frac1y+\frac{z}{y^2}+\frac{y}{xz}+\frac1x.
\]
Its classical period begins
\[
\sum_{d\geq0}\mathrm{CT}(W_{V_{22}}^d)t^d
=1+12t^2+60t^3+636t^4+5760t^5+\cdots,
\]
and agrees through order $t^9$ with the regularized quantum period computed
for $V_{22}$ in Section~H.17 of the same article.
Clearing the logarithmic derivatives and saturating by $xyz$ gives the reduced
critical ideal
\[
I_{\mathrm{crit}}=(x-z,\ y-z,\ z^3+2z^2-2).
\]
Consequently,
\[
\CritLoc(W_{V_{22}})=
\{(a,a,a):a^3+2a^2-2=0\}.
\]
At such a point the critical value is
\[
c_a=\frac{a^3+4a^2+4a+4}{a},
\]
and the logarithmic Hessian determinant is
\[
\det(\theta_i\theta_jW_{V_{22}})(a,a,a)
=\frac{4(a+1)^4(a^3+a^2+1)}{a^3}.
\]
Its numerator is coprime to $a^3+2a^2-2$, so all three points are isolated
Morse singularities. Therefore
\[
\overline{MF((\mathbb{C}^*)^3,W_{V_{22}})}
\cong \bigoplus_{a:\,a^3+2a^2-2=0}\MorseCat{3}.
\]
\end{example}

\begin{example}[The Landau--Ginzburg model for $V_{18}$]
Appendix~A of \cite{CoatesCortiGalkinKasprzyk2016} gives the following
$13$-term Laurent polynomial for the codimension-$2$ linear section of the
$G_2$-Grassmannian:
\[
\begin{aligned}
W_{V_{18}}={}&xy^2+2xy+x+2y+z+\frac1z+\frac2y+\frac1{yz}
+\frac1x\\
&+\frac{2z}{xy}+\frac1{xy}+\frac1{xy^2}+\frac{z}{x^2y^2}.
\end{aligned}
\]
Its classical period begins
\[
\sum_{d\geq0}\mathrm{CT}(W_{V_{18}}^d)t^d
=1+18t^2+120t^3+1566t^4+18360t^5+\cdots,
\]
and agrees through order $t^9$ with the regularized quantum period computed
for $V_{18}$ in Section~H.16 of the same article.

The saturated logarithmic critical ideal is
\[
(2x+12z-11,\ 3y-6z+1,\ 36z^2-12z-11).
\]
Thus its two points are indexed by the roots of the last polynomial:
\[
P_a=\left(\frac{11-12a}{2},\frac{6a-1}{3},a\right),
\qquad
36a^2-12a-11=0.
\]
At these points,
\[
W_{V_{18}}(P_a)=18a+3,
\qquad
\det(\theta_i\theta_jW_{V_{18}})(P_a)=9(42a+17).
\]
The determinant is relatively prime to $36a^2-12a-11$, so both points are
Morse. Their critical values are $6\pm6\sqrt3$. Consequently,
\[
\overline{MF((\mathbb C^*)^3,W_{V_{18}})}
\simeq
\bigoplus_{a:\,36a^2-12a-11=0}\MorseCat{3}.
\]
\end{example}


\begin{example}[The Landau--Ginzburg model for $V_{16}$]
Appendix~A of \cite{CoatesCortiGalkinKasprzyk2016} gives the following
$14$-term Laurent polynomial:
\[
\begin{aligned}
W_{V_{16}}={}&x+\frac{2x}{yz}+\frac{x}{y^2z^2}
+yz^2+2yz+y+2z+\frac2z+\frac1y\\
&+\frac{2}{yz}+\frac{1}{yz^2}+\frac{z}{x}+\frac2x+\frac{1}{xz}.
\end{aligned}
\]
Its classical period begins
\[
\sum_{d\geq0}\mathrm{CT}(W_{V_{16}}^d)t^d
=1+24t^2+192t^3+2904t^4+40320t^5+\cdots,
\]
and agrees through order $t^9$ with the regularized quantum period computed
for $V_{16}$ in Section~H.15 of the same article.
The saturated logarithmic critical ideal is generated by
\[
xy-x-6y+z+7,
\qquad
xz+x+4y-2z-6,
\]
\[
2y^2-z-3,
\qquad
(y-1)(z-1),
\qquad
(z-1)(z+1).
\]
It follows that the critical locus is the disjoint union of the smooth curve
\[
\Gamma=\{y=1,\ z=-1\}
\]
and the two points
\[
P_a=(4-2a,a,1),
\qquad
a^2=2.
\]
Along the curve, $W_{V_{16}}=-4$. Put $t=x$, $U=y-1$, and $V=z+1$.
Then
\[
\widehat X_\Gamma
\simeq
\operatorname{Spf}\!\left(\mathbb C[t^{\pm1}][[U,V]]\right)
\]
and
\[
W_{V_{16}}+4
\equiv
tU^2-2tUV+\frac{t^2-1}{t}V^2
\pmod{(U,V)^3}.
\]
The quadratic form splits globally as
\[
t\left(U-\left(1+\frac1t\right)V\right)
 \left(U-\left(1-\frac1t\right)V\right).
\]
The corresponding linear change has determinant $2/t$, a unit on $\Gamma$.
The relative formal Morse lemma therefore gives a globally split, untwisted
Morse--Bott model on $\Gamma$. At $P_a$ one has
\[
W_{V_{16}}(P_a)=8a+8,
\qquad
\det(\theta_i\theta_jW_{V_{16}})(P_a)=16(17a+24),
\]
which is nonzero modulo $a^2-2$. Therefore
\[
\overline{MF((\mathbb{C}^*)^3,W_{V_{16}})}
\cong
\MBcat{\Gamma}
\oplus
\bigoplus_{a^2=2}\MorseCat{3}.
\]
\end{example}




\begin{example}[The Landau--Ginzburg model for $V_{14}$]
Appendix~A of \cite{CoatesCortiGalkinKasprzyk2016} gives the following
$15$-term Laurent polynomial for the codimension-$5$ linear section of
$\operatorname{Gr}(2,6)$:
\[
\begin{aligned}
W_{V_{14}}={}&xz+x+\frac{x}{yz}+yz^3+3yz^2+3yz+y+z+\frac3z\\
&+\frac1{yz}+\frac3{yz^2}+\frac1{y^2z^3}
+\frac zx+\frac1x+\frac1{xyz}.
\end{aligned}
\]
Its classical period begins
\[
\sum_{d\geq0}\mathrm{CT}(W_{V_{14}}^d)t^d
=1+32t^2+312t^3+5520t^4+91680t^5+\cdots,
\]
and agrees through order $t^9$ with the regularized quantum period computed
for $V_{14}$ in Section~H.14 of the same article.

Set
\[
A=yz^2+yz+1
\]
and
\[
B=x^2yz^2+xy^2z^4+2xy^2z^3+xy^2z^2+2xyz+x+yz^2.
\]
Then the fiber at $-4$ has the exact factorization
\[
W_{V_{14}}+4=\frac{AB}{xy^2z^3}.
\]
The hypersurface $V(A)$ is smooth because
$\partial A/\partial y=z(z+1)$ is a unit there. On $A=0$ one has
\[
y=-\frac1{z(z+1)},
\qquad
B=-\frac{z(x-1)^2}{z+1}.
\]
Consequently, the positive-dimensional critical component is the smooth curve
\[
\Gamma=V(A,B)
=\left\{
\left(1,-\frac1{t(t+1)},t\right)
\;\middle|\;
t\in\mathbb C^*\setminus\{-1\}
\right\},
\]
and $W_{V_{14}}|_\Gamma=-4$.

Away from $V(A)$, saturation of the logarithmic critical ideal gives
\[
(2x-3z+1,\ y+6z-7,\ (z-1)(3z+1)).
\]
Thus there are two isolated critical points,
\[
P_+=(1,1,1),
\qquad
P_-=\left(-1,9,-\frac13\right).
\]
Their critical values and logarithmic Hessian determinants are
\[
\begin{array}{c|cc}
 & W_{V_{14}} & \det(\theta_i\theta_jW_{V_{14}})\\ \hline
P_+ & 23 & 1458\\
P_- & -5 & -2.
\end{array}
\]
Hence both isolated points are Morse.

It remains to identify the formal category along $\Gamma$. Let
\[
R=\mathbb C[t^{\pm1},(t+1)^{-1}],
\qquad
U=A,
\qquad
V=x-1.
\]
Since $y=(U-1)/(t(t+1))$, the pair $(U,V)$ gives global formal normal
coordinates and
\[
\widehat X_\Gamma\simeq\operatorname{Spf}(R[[U,V]]).
\]
In these coordinates the exact shifted potential is
\[
W_{V_{14}}+4
=
\frac{U(t+1)}{t(U-1)^2(V+1)}
\left((t+1)(V+1)U^2+t(U-1)V^2\right).
\]
After multiplication by a unit, this becomes
\[
q+U^3V+cU^2V^2,
\qquad
q=U(U^2-cV^2),
\qquad
c=\frac{t}{t+1}\in R^*.
\]
The relative derivatives of $q$ are
\[
q_U=3U^2-cV^2,
\qquad
q_V=-2cUV.
\]
The identities
\[
\begin{aligned}
U^3&=\frac U3q_U-\frac V6q_V,
&U^2V&=-\frac{U}{2c}q_V,\\
UV^2&=-\frac{V}{2c}q_V,
&V^3&=-\frac Vc q_U-\frac{3U}{2c^2}q_V
\end{aligned}
\]
give the usual relative finite-determinacy argument: successive
$R$-linear formal coordinate changes remove all terms of order at least four.
Thus the completed potential is globally contact-equivalent to
\[
U\left(U^2-\frac{t}{t+1}V^2\right).
\]
Its discriminant class is
\[
\left[\frac{t}{t+1}\right]=[t(t+1)]
\quad\text{in }R^*/R^{*2},
\]
and is nontrivial. The formal singular locus of
$U(U^2-cV^2)$ is the zero section $V(U,V)$, so the formal comparison theorem
applies. Hence the curve contributes the twisted periodic
$D_4$ category $\Dfourcat{\Gamma}{t(t+1)}$. Therefore
\[
\overline{MF((\mathbb C^*)^3,W_{V_{14}})}
\simeq
\Dfourcat{\Gamma}{t(t+1)}
\oplus\MorseCat{3}\oplus\MorseCat{3}.
\]
\end{example}

\begin{example}[The Landau--Ginzburg model for $V_{12}$]
Appendix~A of \cite{CoatesCortiGalkinKasprzyk2016} gives the following
$18$-term Laurent polynomial:
\[
\begin{aligned}
W_{V_{12}}={}&x^2y^3z+x^2y^2z+2xy^2z+xy^2+2xyz+2xy+x+yz+3y+z\\
&+\frac2y+\frac1x+\frac1{xz}+\frac2{xy}+\frac3{xyz}
+\frac1{xy^2}+\frac3{xy^2z}+\frac1{xy^3z}.
\end{aligned}
\]
Its classical period begins
\[
\sum_{d\geq0}\mathrm{CT}(W_{V_{12}}^d)t^d
=1+48t^2+600t^3+13176t^4+276480t^5+\cdots,
\]
and agrees through order $t^9$ with the regularized quantum period computed
for $V_{12}$ in Section~H.13 of the same article.

Set
\[
\begin{aligned}
H={}&x^3y^5z^2+2x^2y^4z^2+x^2y^4z+x^2y^3z
+xy^3z^2+3xy^3z\\
&+2xy^2z+y^2z+y^2+yz+2y+1.
\end{aligned}
\]
The critical fiber at $-5$ admits the exact factorization
\[
W_{V_{12}}+5=\frac{(y+1)H}{xy^3z}.
\]
Its saturated logarithmic critical ideal is generated by
\[
16x^2z-32xz+222yz-237y+238z-221,
\]
\[
(y+1)(x-2z+1),
\qquad
(y+1)(9y-16z+8),
\qquad
(y+1)(32z^2-32z-1).
\]
Hence the critical locus is the disjoint union of the smooth curve
\[
\Gamma=\{y=-1,\ z(x-1)^2+1=0\}
\]
and the two points
\[
P_a=\left(2a-1,\frac{16a-8}{9},a\right),
\qquad
32a^2-32a-1=0.
\]
The curve is parametrized by
\[
(x,y,z)=\left(t,-1,-\frac1{(t-1)^2}\right),
\qquad
t\in\mathbb C^*,\quad t\neq1,
\]
and $W_{V_{12}}|_\Gamma=-5$. Put
\[
U=y+1,
\qquad
V=z(x-1)^2+1.
\]
The factor above restricts on $U=0$ to
\[
H|_{y=-1}=-xzV.
\]
Since $x$ and $z$ are units on the torus, the two factors of the fiber meet
precisely along $\Gamma$.
The Jacobian minor in the $(y,z)$ variables is $(x-1)^2$, a unit on
$\Gamma$, and the inverse substitution
\[
x=t,
\qquad
y=-1+U,
\qquad
z=\frac{V-1}{(t-1)^2}
\]
gives the global formal splitting
\[
\widehat X_\Gamma
\simeq
\operatorname{Spf}\!\left(
\mathbb C[t^{\pm1},(t-1)^{-1}][[U,V]]
\right).
\]
Modulo $(U,V)^3$, the shifted potential is
\[
W_{V_{12}}+5
\equiv
\frac{t^3-t^2+3t-1}{t(t-1)}U^2+UV.
\]
The determinant of the transverse Hessian is $-1$. More strongly, the exact
factorization by $U=y+1$ makes
\[
\widetilde V:=\frac{W_{V_{12}}+5}{U}
\]
a regular element of the completed coordinate ring, and
\[
\widetilde V
\equiv
V+\frac{t^3-t^2+3t-1}{t(t-1)}U
\pmod{(U,V)^2}.
\]
Hence $(U,\widetilde V)$ is a global formal coordinate system and the shifted
potential is exactly $U\widetilde V$. The quadratic bundle is therefore
globally hyperbolic and untwisted. At the isolated points,
\[
W_{V_{12}}(P_a)=32a-4,
\qquad
\det(\theta_i\theta_jW_{V_{12}})(P_a)
=\frac{8(1154a+35)}{3},
\]
and the determinant is nonzero modulo $32a^2-32a-1$. Therefore
\[
\overline{MF((\mathbb{C}^*)^3,W_{V_{12}})}
\cong
\MBcat{\Gamma}
\oplus
\bigoplus_{a:\,32a^2-32a-1=0}\MorseCat{3}.
\]
\end{example}


\begin{example}[The Landau--Ginzburg model for $V_{10}$]
Appendix~A of \cite{CoatesCortiGalkinKasprzyk2016} gives the following
$21$-term Laurent polynomial for the complete intersection of two
hyperplanes and a quadric in $\operatorname{Gr}(2,5)$:
\begin{align*}
W_{V_{10}}={}&xyz^3+3xyz^2+3xyz+xy+xz^2+2xz+x
+yz^2+2yz+y \\
&+3z+\frac3z+\frac2y+\frac2{yz}+\frac2x+\frac2{xz}
+\frac3{xyz}+\frac3{xyz^2} \\
&+\frac1{xy^2z^2}+\frac1{x^2yz^2}+\frac1{x^2y^2z^3}.
\end{align*}
Set
\[
F=xyz^2+xyz+1,
\qquad
G=xyz^2+xyz+xz+yz+1.
\]
Grouping the displayed terms gives the exact identity
\[
W_{V_{10}}+6=\frac{F^2G}{x^2y^2z^3}.
\]
The product before the shift has constant term $6$, and its other $21$
monomials agree term by term with the polynomial in the cited table. Its
classical period begins
\[
\sum_{d\geq0}\mathrm{CT}(W_{V_{10}}^d)t^d
=1+78t^2+1320t^3+37746t^4+1051920t^5+\cdots,
\]
and agrees through order $t^9$ with the regularized quantum period computed
for $V_{10}$ in Section~H.12 of the same article.

Every logarithmic derivative of $W_{V_{10}}$ is divisible by $F$. In
particular,
\[
\theta_zW_{V_{10}}
=\frac{(xyz^2-1)F(3xyz^2+3xyz+2xz+2yz+3)}{x^2y^2z^3}.
\]
Thus the smooth surface
\[
S=V(F)
\]
is contained in the critical locus over the value $-6$. It is smooth because
$\partial F/\partial x=yz(z+1)$ is a unit on $S$. Saturating the logarithmic
critical ideal by $F$ gives
\[
(\theta_xW_{V_{10}},\theta_yW_{V_{10}},\theta_zW_{V_{10}}):F^\infty
=
(x-4z+2,\ y-4z+2,\ 4z^2-2z-1).
\]
Hence the only critical points outside $S$ are
\[
P_a=(4a-2,4a-2,a),
\qquad
4a^2-2a-1=0.
\]
At these points,
\[
W_{V_{10}}(P_a)=40a+6,
\qquad
\det(\theta_i\theta_jW_{V_{10}})(P_a)=64(123a+38).
\]
The determinant is relatively prime to $4a^2-2a-1$, so both points are
Morse. Therefore
\[
\CritLoc(W_{V_{10}})_{\mathrm{red}}=S\sqcup\{P_a:4a^2-2a-1=0\}.
\]

The coefficient $G$ vanishes on the smooth curve
\[
K=V(F,G)
=\{y=-x,\ x^2z(z+1)=1\}\subset S.
\]
Indeed, $G-F=z(x+y)$, and the Jacobian minor
\[
\det\!\left(\frac{\partial(F,G)}{\partial(x,y)}\right)\bigg|_K
=-2xz^2(z+1)
\]
is a unit on $K$. Thus $(F,G)$ is a global system of formal normal
coordinates along $K$. On the open stratum $S^\circ=S\setminus K$, the factor
$G$ is a unit and the shifted potential has transverse contact type $F^2$,
giving $\Aper{S^\circ}{1}$. Along $K$, the exact factorization above gives the
untwisted crossing form $u^2v$ up to multiplication by a unit.

Let
\[
R=
\frac{\mathbb C[x^{\pm1},z^{\pm1},(z+1)^{-1}]}
{(x^2z(z+1)-1)}
\]
be the coordinate ring of $K$, and put
\[
\mathcal P_{V_{10}}
:=
\overline{MF^t_0\!\left(
\operatorname{Spf}(R[[u,v]]),u^2v
\right)}.
\]
Define the connected contribution of the fiber at $-6$ by
\[
\mathcal C_{V_{10},-6}
:=
\overline{MF^t_{-6}(\widehat X_S,\widehat W_{V_{10}})}.
\]
Supported localization gives the Verdier sequence
\[
\mathcal P_{V_{10}}
\longrightarrow
\mathcal C_{V_{10},-6}
\longrightarrow
\Aper{S^\circ}{1}.
\]
No splitting is asserted. Consequently,
\[
\overline{MF((\mathbb C^*)^3,W_{V_{10}})}
\simeq
\mathcal C_{V_{10},-6}
\oplus
\bigoplus_{a:\,4a^2-2a-1=0}\MorseCat{3}.
\]
\end{example}

\begin{example}[The Landau--Ginzburg model for $V_8$]
Appendix~A of \cite{CoatesCortiGalkinKasprzyk2016} gives the following
$18$-term Laurent polynomial for the complete intersection of three quadrics
in $\mathbb P^6$:
\begin{align*}
W_{V_{8}} &= xy^2 + 2xyz^2 + 4xyz + 2xy + xz^4 + 4xz^3 + 6xz^2 + 4xz + x \\
&\quad +\frac{4y}{z} + 4z + \frac{4}{z} + \frac{6}{xz^2} + \frac{2}{xy} + \frac{4}{xyz} + \frac{2}{xyz^2} + \frac{4}{x^2yz^3} + \frac{1}{x^3y^2z^4}
\end{align*}

Set
\[
A=1+\frac1{xyz},
\qquad
B=z+1,
\qquad
U=yA^2+B^2.
\]
Grouping the displayed terms gives the exact identity
\[
U=y+(z+1)^2+\frac{2}{xz}+\frac{1}{x^2yz^2},
\qquad
W_{V_8}+8=xU^2.
\]
The product $xU^2$ has constant term $8$, and its other $18$ monomials agree
term by term with the polynomial in the cited table. Its classical period
begins
\[
\sum_{d\geq0}\mathrm{CT}(W_{V_8}^d)t^d
=1+152t^2+3840t^3+157656t^4+\cdots,
\]
where $\mathrm{CT}$ denotes the constant term. This agrees through order
$t^9$ with the regularized quantum period computed for $V_8$ in
Section~H.11 of the same article.

Put $q=(xyz)^{-1}$. The logarithmic critical equations factor as
\[
\begin{aligned}
\theta_xW_{V_8}&=xU\bigl(U-4yqA\bigr),\\
\theta_yW_{V_8}&=2xyU(1-q^2),\\
\theta_zW_{V_8}&=4xU\bigl(zB-yqA\bigr).
\end{aligned}
\]
Thus $S=V(U)$ is contained in the critical locus over $-8$. Unlike the
smooth square loci in the $B_2$ and $MM_{2-8}$ models, however, $S$ is not
smooth. Indeed,
\[
\Sing(S)=K=V(A,B)
=\{z=-1,\ xy=1\}.
\]
On $S^\circ=S\setminus K$, the global equation $U$ is a regular normal
coordinate and $W_{V_8}+8=xU^2$ gives the untwisted $A_1$ contribution
$\Aper{S^\circ}{1}$.

On the locus where $U\neq0$, the second equation gives $q=\pm1$. The choice
$q=-1$ would give $A=0$ and contradict the first equation. Hence $q=1$, and
the remaining equations yield the unique point
\[
P=(1,1,1),
\qquad
W_{V_8}(P)=56.
\]
Consequently,
\[
\CritLoc(W_{V_8})_{\mathrm{red}}=S\sqcup\{P\},
\qquad
\CritVal(W_{V_8})=\{-8,56\}.
\]
The logarithmic Hessian at $P$ is
\[
(\theta_i\theta_jW_{V_8})(P)=
\begin{pmatrix}
64&32&96\\
32&32&32\\
96&32&192
\end{pmatrix},
\qquad
\det(\theta_i\theta_jW_{V_8})(P)=2^{15}.
\]
Thus $P$ is Morse and contributes $\MorseCat{3}$.

The stratum $K$ carries additional crossing data and cannot be absorbed into a
Morse--Bott symbol. Parametrize it by
\[
(x,y,z)=\left(t,t^{-1},-1\right),
\qquad t\in\mathbb C^*.
\]
The functions $(A,B)$ give global formal normal coordinates along $K$: the
Jacobian minor in the $(y,z)$ variables restricts to $t$, hence is a unit, and
\[
\widehat X_K
\simeq
\operatorname{Spf}\!\left(\mathbb C[t^{\pm1}][[A,B]]\right).
\]
In these coordinates the exact potential is
\[
W_{V_8}+8
=x\bigl(yA^2+B^2\bigr)^2.
\]
Here $x=t$, $z=B-1$, and
\[
y=\frac{1}{t(B-1)(A-1)}
\]
in the completed coordinate ring. Its leading transverse form is
\[
\frac{(A^2+tB^2)^2}{t};
\]
the square class $[t]$ is nontrivial on $K$, so this is a genuinely twisted
formal crossing family rather than a global $u^2v^2$ product.

Put
\[
\mathcal C_{V_8,-8}
:=
\overline{MF^t_{-8}(\widehat X_S,\widehat W_{V_8})}.
\]
Also set
\fitdisplay{\mathcal P_{V_8}:=
\overline{MF^t_0\!\left(
\operatorname{Spf}\mathbb C[t^{\pm1}][[A,B]],
t\left(\frac{A^2}{t(B-1)(A-1)}+B^2\right)^2
\right)}}
Supported localization gives the Verdier sequence
\[
\mathcal P_{V_8}
\longrightarrow
\mathcal C_{V_8,-8}
\longrightarrow
\Aper{S^\circ}{1}.
\]
It is not asserted to split. The hypothesis
$\Sing(\operatorname{Spec}B)\subseteq V((U))$ of
Theorem~\ref{thm:general_formal} holds because the formal zero fiber is defined
by the nilpotent equation $U^2=0$. Consequently,
\[
\overline{MF((\mathbb{C}^*)^3,W_{V_8})}
\simeq
\mathcal C_{V_8,-8}\oplus\MorseCat{3}.
\]

\end{example}

\begin{example}[The Landau--Ginzburg model for $V_6$]
Appendix~A of \cite{CoatesCortiGalkinKasprzyk2016} lists a $29$-term
Laurent polynomial for the complete intersection of a quadric and a cubic in
$\mathbb P^5$. Grouping those terms gives the exact expression
\[
W_{V_6}
=x(y+1)^2\left(z+1+\frac1{xyz}\right)^3-12.
\]
The product before the shift has constant term $12$, and its other $29$
monomials agree term by term with the polynomial in the cited table. Its
classical period begins
\[
\sum_{d\geq0}\mathrm{CT}(W_{V_6}^d)t^d
=1+396t^2+17616t^3+1217052t^4+\cdots,
\]
and agrees through order $t^9$ with the regularized quantum period computed
for $V_6$ in Section~H.10 of the same article.

Put
\[
A=y+1,
\qquad
U=z+1+\frac1{xyz},
\qquad
q=\frac1{xyz}.
\]
Then $W_{V_6}+12=xA^2U^3$, and the logarithmic critical equations are
\[
\begin{aligned}
\theta_xW_{V_6}&=xA^2U^2(U-3q),\\
\theta_yW_{V_6}&=xAU^2\bigl(2yU-3Aq\bigr),\\
\theta_zW_{V_6}&=3xA^2U^2(z-q).
\end{aligned}
\]
Hence the critical fiber at $-12$ contains the two smooth surfaces
\[
S_A=V(A),
\qquad
S_U=V(U).
\]
They meet transversely along the smooth curve
\[
K=V(A,U)
=\left\{
\left(\frac1{t(t+1)},-1,t\right)
\;\middle|\;
t\in\mathbb C^*\setminus\{-1\}
\right\}.
\]
On the complement of $S_A\cup S_U$, the critical equations give
$z=q$, $U=3q$, and $2yU=3Aq$. It follows that $y=z=q=1$ and therefore
the only remaining critical point is
\[
P=(1,1,1),
\qquad
W_{V_6}(P)=96.
\]
Thus
\[
\CritLoc(W_{V_6})_{\mathrm{red}}
=(S_A\cup S_U)\sqcup\{P\},
\qquad
\CritVal(W_{V_6})=\{-12,96\}.
\]
The logarithmic Hessian at $P$ is
\[
(\theta_i\theta_jW_{V_6})(P)=
\begin{pmatrix}
72&72&108\\
72&126&108\\
108&108&216
\end{pmatrix},
\qquad
\det(\theta_i\theta_jW_{V_6})(P)=2^5 3^8,
\]
so $P$ is Morse.

The fiber at $-12$ is not a single residual $A_2$ family. On
$S_A^\circ=S_A\setminus K$ it has transverse form $A^2$, while on
$S_U^\circ=S_U\setminus K$ it has transverse form $U^3$. Along $K$, the
functions $(A,U)$ are formal normal coordinates and the shifted potential has
contact normal form $a^2u^3$. Define
\[
\mathcal C_{V_6,-12}
:=\overline{MF_{-12}((\mathbb C^*)^3,W_{V_6})}
\]
and let $\mathcal P_{V_6}$ be the category supported on $K$ for this
$a^2u^3$ crossing family. Supported localization gives the Verdier sequence
\[
\mathcal P_{V_6}
\longrightarrow
\mathcal C_{V_6,-12}
\longrightarrow
\Aper{S_A^\circ}{1}\oplus\Aper{S_U^\circ}{2}.
\]
No splitting of this sequence is asserted. The complete open-torus category
is therefore
\[
\overline{MF((\mathbb C^*)^3,W_{V_6})}
\simeq
\mathcal C_{V_6,-12}\oplus\MorseCat{3}.
\]
\end{example}

\begin{example}[The Landau--Ginzburg model for $V_4$]
Appendix~A of \cite{CoatesCortiGalkinKasprzyk2016} gives the following
$34$-term Laurent polynomial for the quartic threefold
$V_4\subset\mathbb P^4$:
\fitdisplay{\begin{aligned}
W_{V_4} &= xy^4 + 4xy^3z + 4xy^3 + 6xy^2z^2 + 12xy^2z + 6xy^2 + 4xyz^3 + 12xyz^2 + 12xyz \\
&\quad + 4xy + xz^4 + 4xz^3 + 6xz^2 + 4xz + x + \frac{4y^2}{z} + 12y + \frac{12y}{z} + 12z \\
&\quad + \frac{12}{z} + \frac{4z^2}{y} + \frac{12z}{y} + \frac{12}{y} + \frac{4}{yz} + \frac{6}{xz^2} + \frac{12}{xyz} + \frac{12}{xyz^2} + \frac{6}{xy^2} \\
&\quad + \frac{12}{xy^2z} + \frac{6}{xy^2z^2} + \frac{4}{x^2y^2z^3} + \frac{4}{x^2y^3z^2} + \frac{4}{x^2y^3z^3} + \frac{1}{x^3y^4z^4}
\end{aligned}}

Set
\[
U=y+z+1+\frac{1}{xyz}.
\]
Grouping the displayed terms by their multinomial coefficients gives the
exact identity
\[
W_{V_4}=xU^4-24.
\]
Indeed, $xU^4$ has constant term $24$, and its other $34$ monomials agree
term by term with the polynomial in the cited table. Its classical period
begins
\[
\sum_{d\geq0}\mathrm{CT}(W_{V_4}^d)t^d
=1+1944t^2+215808t^3+35295192t^4+\cdots,
\]
where $\mathrm{CT}$ denotes the constant term. This agrees through order
$t^9$ with the regularized quantum period computed for $V_4$ in
Section~H.9 of the same article.

Put
\[
q=\frac1{xyz}.
\]
The logarithmic critical equations factor as
\[
\theta_xW_{V_4}=xU^3(U-4q),
\qquad
\theta_yW_{V_4}=4xU^3(y-q),
\qquad
\theta_zW_{V_4}=4xU^3(z-q).
\]
Thus one critical component is the surface
\[
S=V(U),
\qquad
W_{V_4}|_S=-24.
\]
It is smooth because
$\partial_xU=-1/(x^2yz)$ is a unit on the torus. On the complement of $S$,
the equations give $y=q$, $z=q$, and $U=4q$. Since
$U=y+z+1+q$, these equalities imply $q=1$ and yield the single point
\[
P=(1,1,1),
\qquad
W_{V_4}(P)=232.
\]
Consequently,
\[
\CritLoc(W_{V_4})_{\mathrm{red}}=S\sqcup\{P\},
\qquad
\CritVal(W_{V_4})=\{-24,232\}.
\]
The logarithmic Hessian at $P$ is
\[
(\theta_i\theta_jW_{V_4})(P)=
\begin{pmatrix}
192&256&256\\
256&512&256\\
256&256&512
\end{pmatrix},
\qquad
\det(\theta_i\theta_jW_{V_4})(P)=2^{22}.
\]
Hence $P$ is Morse. Along $S$, the identity $W_{V_4}+24=xU^4$ and the fact
that $x$ is a unit give a global residual $A_3$ normal form. Therefore
Proposition~\ref{prop:Ade_category_model} and
Theorem~\ref{thm:general_formal} give
\[
\overline{MF((\mathbb{C}^*)^3, W_{V_{4}})}
\cong \Aper{S}{3}\oplus \MorseCat{3}.
\]
\end{example}
\begin{example}[The Landau--Ginzburg model for $V_2$]
Appendix~A of \cite{CoatesCortiGalkinKasprzyk2016} lists an expanded
Laurent polynomial for the sextic hypersurface
$V_2\subset\mathbb P(1,1,1,1,3)$. Grouping its terms by their multinomial
coefficients gives the compact expression
\[
W_{V_2}
=x\left(y+z+1+\frac{1}{xy^3z}\right)^6-120.
\]
Indeed, $x(y+z+1+1/(xy^3z))^6$ has constant term $120$; after this term is
removed, its remaining $83$ monomials agree term by term with the polynomial
in the cited table. As an independent check, its classical period begins
\[
\sum_{d\geq0}\mathrm{CT}(W_{V_2}^d)t^d
=1+68760t^2+55200000t^3+61054781400t^4+\cdots,
\]
where $\mathrm{CT}$ denotes the constant term. This series agrees through
order $t^9$ with the regularized quantum period computed for $V_2$ in
Section~H.8 of the same article.

Put
\[
F=y+z+1+\frac{1}{xy^3z},
\qquad
q=\frac{1}{xy^3z}.
\]
The logarithmic critical equations factor as
\[
\theta_xW_{V_2}=xF^5(F-6q),
\qquad
\theta_yW_{V_2}=6xF^5(y-3q),
\qquad
\theta_zW_{V_2}=6xF^5(z-q).
\]
Hence one critical component is the smooth surface
\[
S=V(F),
\qquad
W_{V_2}|_S=-120.
\]
It is smooth because
$\partial_xF=-1/(x^2y^3z)$ is a unit on the torus.  On the complement of
$S$, the equations give $y=3q$, $z=q$, and $F=6q$. Since
$F=y+z+1+q$, these equalities imply $q=1$ and yield the single point
\[
P=\left(\frac1{27},3,1\right),
\qquad
W_{V_2}(P)=1608.
\]
Thus
\[
\CritLoc(W_{V_2})_{\mathrm{red}}=S\sqcup\{P\},
\qquad
\CritVal(W_{V_2})=\{-120,1608\}.
\]
The logarithmic Hessian at this point is
\[
(\theta_i\theta_jW_{V_2})(P)=
\begin{pmatrix}
1440&5184&1728\\
5184&20736&5184\\
1728&5184&3456
\end{pmatrix},
\]
with determinant $2^{17}3^9$, so $P$ is Morse. Along $S$ the exact identity
$W_{V_2}+120=xF^6$, together with the fact that $x$ is a unit on the torus,
gives a global residual $A_5$ normal form. Consequently,
\[
\overline{MF((\mathbb C^*)^3,W_{V_2})}
\cong
\Aper{S}{5}\oplus\MorseCat{3}.
\]
\end{example}

\begin{example}[The Laurent mirror for $MM_{2-1}$]
For the degree-$4$ hypersurface-in-a-product model, Appendix~A of
\cite{CoatesCortiGalkinKasprzyk2016} gives
\fitdisplay{\begin{aligned}
W_{MM_{2-1}}={}&x^7y^7z^{18}+6x^6y^6z^{15}
 +6x^5y^5z^{13}+15x^5y^5z^{12}
 +30x^4y^4z^{10}+20x^4y^4z^9\\
&+x^4y^3z^9+x^3y^4z^9+15x^3y^3z^8
 +60x^3y^3z^7+15x^3y^3z^6+3x^3y^2z^6+3x^2y^3z^6\\
&+60x^2y^2z^5+60x^2y^2z^4+6x^2y^2z^3
 +3x^2yz^4+3x^2yz^3+3xy^2z^4+3xy^2z^3\\
&+20xyz^3+90xyz^2+30xyz+xy+6xz+x+6yz+y
 +\frac{60}{z}+\frac{6}{z^2}+\frac{3}{yz}+\frac{3}{yz^2}\\
&+\frac{3}{xz}+\frac{3}{xz^2}
 +\frac{15}{xyz^2}+\frac{60}{xyz^3}+\frac{15}{xyz^4}
 +\frac{3}{xy^2z^4}+\frac{3}{x^2yz^4}\\
&+\frac{30}{x^2y^2z^5}+\frac{20}{x^2y^2z^6}
 +\frac{1}{x^2y^3z^6}+\frac{1}{x^3y^2z^6}
 +\frac{6}{x^3y^3z^7}+\frac{15}{x^3y^3z^8}\\
&+\frac{6}{x^4y^4z^{10}}+\frac{1}{x^5y^5z^{12}}.
\end{aligned}}

The large expression admits a compact form after setting
\[
p=xy,
\qquad
F=p^2z^5+pz^2+1,
\qquad
R=\frac{F^3}{p^2z^6}.
\]
Indeed, direct expansion gives the exact identity
\[
W_{MM_{2-1}}+60=\frac{R(R+x+y)}{p}.
\]
All three cleared logarithmic derivatives contain $F^2$.  Hence
\[
S=V(F),
\qquad
W_{MM_{2-1}}|_S=-60,
\]
is a critical surface.  It is smooth: in the coordinates $(p,z)$, the
simultaneous vanishing of $\partial_pF$ and $F$ forces
$p=-1/(2z^3)$, while then $\partial_zF=1/(4z^2)\neq0$.

On $S$ the coefficient of $F^3$ is
\[
\left.\frac{R+x+y}{p}\frac{R}{F^3}\right|_S
=\frac{x+y}{p^3z^6}.
\]
It is a unit away from the smooth curve
\[
K=V(F,x+y)\subset S.
\]
Thus $S^\circ=S\setminus K$ carries a global residual $A_2$ model.  Along
$K$, put $V=R+x+y$.  The change from $(F,x+y)$ to $(F,V)$ is formally
triangular, and
\[
W_{MM_{2-1}}+60=\frac{R}{p}V
\]
has the exact transverse form $u^3v$ up to a unit.  Define
\[
\mathcal C_{MM_{2-1},-60}:=
\overline{MF_{-60}((\mathbb C^*)^3,W_{MM_{2-1}})}
\]
and let $\mathcal P_{MM_{2-1}}$ denote the torsion matrix-factorization
category of this $u^3v$ formal model along $K$.  Supported localization gives
\[
\mathcal P_{MM_{2-1}}\longrightarrow
\mathcal C_{MM_{2-1},-60}\longrightarrow
\Aper{S^\circ}{2},
\]
and no splitting is asserted.

There is a second positive-dimensional critical component.  The identity
\[
W_{MM_{2-1}}+61
=\frac{(R+x)(R+y)}{xy}
\]
shows that
\[
\Gamma=V(R+x,R+y)
\]
is a quadratic component with critical value $-61$.  Equivalently,
\[
\Gamma=\left\{x=y=a,\,
E(a,z)=0\right\},
\]
where
\[
E(a,z)=\bigl(a^4z^5+a^2z^2+1\bigr)^3+a^5z^6.
\]
A Gr\"obner-basis calculation gives
$(E,\partial_aE,\partial_zE)=(1)$, so $\Gamma$ is smooth.  The functions
$R+x$ and $R+y$ generate its conormal ideal, and the displayed exact product
gives a global hyperbolic quadratic model.  Hence $\Gamma$ contributes
$\MBcat{\Gamma}$.

After removing $S$ and $\Gamma$, the residual critical ideal is
\[
z-6,
\qquad
x-y,
\qquad
108x^2-1.
\]
Thus there are two isolated points
\[
P_a=(a,a,6),
\qquad
108a^2-1=0.
\]
Their critical values and logarithmic Hessian determinants, reduced modulo
$108a^2-1$, are
\[
W_{MM_{2-1}}(P_a)=372+432a,
\]
\[
\det(\theta_i\theta_jW_{MM_{2-1}})(P_a)
=186624(433a+4).
\]
The determinant is coprime to $108a^2-1$, so both points are Morse.  We
conclude that
\[
\overline{MF((\mathbb C^*)^3,W_{MM_{2-1}})}
\cong
\mathcal C_{MM_{2-1},-60}
\oplus\MBcat{\Gamma}
\oplus
\bigoplus_{a:\,108a^2-1=0}\MorseCat{3}.
\]
\end{example}

\begin{example}[The Laurent mirror for $MM_{2-2}$]
The Quantum Lefschetz construction with mirror map gives
\[
\begin{aligned}
W_{MM_{2-2}}={}&xy^2+2xyz+2xy+xz^2+2xz+x
 +\frac{y^2}{z}+4y+\frac{4y}{z}+6z+\frac{6}{z}\\
&+\frac{4z^2}{y}+\frac{14z}{y}+\frac{14}{y}+\frac{4}{yz}
 +\frac{z^3}{y^2}+\frac{4z^2}{y^2}+\frac{6z}{y^2}
 +\frac{4}{y^2}+\frac{1}{y^2z}\\
&+\frac{4}{xz}+\frac{12}{xy}+\frac{12}{xyz}
 +\frac{12z}{xy^2}+\frac{25}{xy^2}+\frac{12}{xy^2z}\\
&+\frac{4z^2}{xy^3}+\frac{12z}{xy^3}+\frac{12}{xy^3}
 +\frac{4}{xy^3z}+\frac{6}{x^2y^2z}
 +\frac{12}{x^2y^3}+\frac{12}{x^2y^3z}\\
&+\frac{6z}{x^2y^4}+\frac{12}{x^2y^4}+\frac{6}{x^2y^4z}
 +\frac{4}{x^3y^4z}+\frac{4}{x^3y^5}
 +\frac{4}{x^3y^5z}+\frac{1}{x^4y^6z}.
\end{aligned}
\]

Put
\[
F=xy(y+z+1)+1.
\]
The logarithmic critical equations have the common factor $F$, and direct
division gives the exact identity
\[
W_{MM_{2-2}}+14=F^2Q,
\]
where
\fitdisplay{Q=
\frac{
x^3y^4z+x^2y^4+2x^2y^3z+2x^2y^3+x^2y^2z^2
+2x^2y^2z+x^2y^2+2xy^2+2xyz+2xy+1
}{x^4y^6z}}
Thus
\[
S=V(F),
\qquad
W_{MM_{2-2}}|_S=-14,
\]
is a smooth critical surface.  Indeed,
\[
\left.\partial_xF\right|_S
=y(y+z+1)=-\frac1x
\]
is a unit, and
\[
Q|_S=-\frac{y+z+1}{y}=\frac1{xy^2}
\]
is also a unit.  Therefore the exact factorization gives a global
nondegenerate quadratic normal model and the surface contributes
$\MBcat{S}$.

After removing the factor $F$ and saturating by $xyz$, the residual critical
ideal is
\[
z-1,
\qquad
x-4y+8,
\qquad
4y^3-16y^2+16y-1.
\]
For every root $t$ of
\[
q(t)=4t^3-16t^2+16t-1,
\]
there is therefore a critical point
\[
P_t=(4t-8,t,1).
\]
Modulo $q(t)$, its critical value and logarithmic Hessian determinant are
\[
W_{MM_{2-2}}(P_t)=2(24t^2-32t-5),
\]
\[
\det(\theta_i\theta_jW_{MM_{2-2}})(P_t)
=128(1264t^2-2237t+144).
\]
The last factor is coprime to $q$, so all three isolated points are Morse.
Consequently,
\[
\overline{MF((\mathbb C^*)^3,W_{MM_{2-2}})}
\cong
\MBcat{S}
\oplus
\bigoplus_{t:\,q(t)=0}\MorseCat{3}.
\]
\end{example}

\begin{example}[The Laurent mirror for $MM_{2-3}$]
For the degree-$8$ hypersurface-in-a-product model, the Laurent potential is
\fitdisplay{\begin{aligned}
W_{MM_{2-3}}={}&x^2y^5z^2+4x^2y^4z^2+6x^2y^3z^2
 +4x^2y^2z^2+x^2yz^2+xy^3z^2+4xy^3z\\
&+2xy^2z^2+12xy^2z+xy^2+xyz^2+12xyz+2xy+4xz+x\\
&+2yz+6y+2z+\frac{2}{z}+\frac{6}{y}+\frac{2}{yz}
 +\frac{1}{xy}+\frac{4}{xyz}+\frac{4}{xy^2z}
 +\frac{1}{xy^2z^2}+\frac{1}{x^2y^3z^2}.
\end{aligned}}

Set
\[
F=xyz(y+1)+1.
\]
The three logarithmic derivatives have the common factor $F$, and the
potential satisfies
\[
W_{MM_{2-3}}+12=F^2Q,
\]
where
\[
Q=
\frac{x^2y^4z^2+2x^2y^3z^2+x^2y^2z^2+xy^2z^2
+2xy^2z+2xyz+xy+1}{x^2y^3z^2}.
\]
Hence
\[
S=V(F),
\qquad
W_{MM_{2-3}}|_S=-12,
\]
is a smooth critical surface, since
$\partial_xF=yz(y+1)=-1/x$ on $S$.  The restriction of the quadratic
coefficient is
\[
Q|_S=-\frac{(y+1)(yz^2+1)}{yz}.
\]
It is a unit on
\[
S^\circ=S\setminus K,
\qquad
K=V(F,yz^2+1).
\]
The curve $K$ is smooth and has the parametrization
\[
(x,y,z)=\left(\frac{t^3}{t^2-1},-\frac1{t^2},t\right),
\qquad
t\in\mathbb C^*,\quad t^2\neq1.
\]
Thus $S^\circ$ carries the quadratic contribution $\MBcat{S^\circ}$.  Along
$K$, the pair $(F,Q)$ gives formal normal coordinates and the exact identity
$W_{MM_{2-3}}+12=F^2Q$ has transverse model $u^2v$.  Put
\[
\mathcal C_{MM_{2-3},-12}:=
\overline{MF_{-12}((\mathbb C^*)^3,W_{MM_{2-3}})}
\]
and let $\mathcal P_{MM_{2-3}}$ be the torsion category of the $u^2v$ formal
model along $K$.  Then
\[
\mathcal P_{MM_{2-3}}\longrightarrow
\mathcal C_{MM_{2-3},-12}\longrightarrow
\MBcat{S^\circ}
\]
is the supported-localization sequence; it is not asserted to split.

The remaining non-isolated component lies in the critical fiber over $-13$.
Set
\[
A=yz^2-1
\]
and
\[
H=x^2(z^2+1)^2+xz^3(2z^2+z+2)+z^6.
\]
Then
\[
\Gamma=V(A,H)
\]
is a smooth curve, and on the hypersurface $A=0$ one has the exact identity
\[
W_{MM_{2-3}}+13=\frac{H^2}{x^2z^8}.
\]
The quadratic Taylor form in the global conormal generators $(A,H)$ is
nondegenerate.  If $\Delta$ denotes its discriminant, direct reduction modulo
$H$ gives
\[
\Delta\equiv-\frac1{x^2z^8}
=-\left(\frac1{xz^4}\right)^2
\quad\text{in }\mathcal O_\Gamma^*.
\]
Thus the discriminant square class is trivial over $\mathbb C$, and the
relative formal Morse lemma gives the untwisted contribution
$\MBcat{\Gamma}$.

After removing $S$ and $\Gamma$, the two remaining points are
\[
P_+=\left(\frac12,1,1\right),
\qquad
P_-=\left(-\frac12,1,-1\right).
\]
Their critical values are
\[
W_{MM_{2-3}}(P_+)=68,
\qquad
W_{MM_{2-3}}(P_-)=36,
\]
and their logarithmic Hessian determinants are respectively
$41472$ and $-25088$.  Hence both are Morse.  Consequently,
\[
\overline{MF((\mathbb C^*)^3,W_{MM_{2-3}})}
\cong
\mathcal C_{MM_{2-3},-12}
\oplus\MBcat{\Gamma}
\oplus\MorseCat{3}\oplus\MorseCat{3}.
\]
\end{example}


\begin{example}
Consider the Landau--Ginzburg model of the blowup of $\mathbb{P}^3$ in the intersection of two cubics defined on the algebraic torus $(\mathbb{C}^*)^3$ by the following 21-term potential:
$$\begin{aligned}
W_{MM_{2-4}} = {} & xy^2 + 2xy + x + \frac{2x}{z} + \frac{2x}{yz} + \frac{x}{y^2z^2} \\
& + y^3z + 3y^2z + 3yz + 4y + z + \frac{3}{y} \\
& + \frac{5}{yz} + \frac{4}{y^2z} + \frac{2}{y^3z^2}+ \frac{1}{x} + \frac{2}{xy}  \\
& + \frac{1}{xy^2} + \frac{2}{xy^2z} + \frac{2}{xy^3z} + \frac{1}{xy^4z^2}.
\end{aligned}$$
Set
\[
v=y+1+\frac{1}{yz}.
\]
A direct regrouping of the Laurent monomials gives the compact form
\[
W_{MM_{2-4}}=xv^2+\frac{v^2}{xy^2}+zv^3+\frac{v^2}{y}-\frac{v}{y}-7.
\]
The critical locus in the torus is obtained from the logarithmic derivatives. The $x$-equation is
\[
x\frac{\partial W_{MM_{2-4}}}{\partial x}=xv^2\left(1-\frac{1}{x^2y^2}\right),
\]
and the branch $v=0$ does not give a simultaneous critical point, since then the $z$-derivative of the term $-v/y$ is nonzero. Hence the critical points lie on the two branches $xy=\pm1$.

On the branch $x=-1/y$, the remaining logarithmic equations reduce to
\[
z=y,\qquad y^2+y-1=0.
\]
Thus this branch gives the two points
\[
\left(-\frac{1}{y},y,y\right),
\qquad
 y=\frac{-1\pm\sqrt{5}}{2}.
\]
Evaluating the potential gives the critical values
\[
\frac{-3\pm5\sqrt{5}}{2}.
\]
On the branch $x=1/y$, the remaining equations reduce to
\[
9y^2+y-9=0,
\qquad
z=\frac{153y-64}{45}.
\]
Equivalently,
\[
y=\frac{-1\pm5\sqrt{13}}{18},
\qquad
x=\frac{1}{y},
\qquad
z=\frac{153y-64}{45}.
\]
These two points have critical values
\[
\frac{33\pm13\sqrt{13}}{2}.
\]
A direct computation of the Hessian at each of the four critical points gives $\det H_W\neq0$. Therefore all singularities are isolated Morse singularities, hence of type $A_1$. Consequently,
\[
\overline{MF((\mathbb{C}^*)^3,W_{MM_{2-4}})}
\cong \bigoplus_{i=1}^{4}\MorseCat{3}.
\]
\end{example}

\begin{example}
Consider the Landau--Ginzburg model of the blowup of a cubic hypersurface in $\mathbb{P}^4$ along a plane cubic, defined on the algebraic torus $(\mathbb{C}^*)^3$ by the following 18-term Laurent potential:
\fitdisplay{\begin{aligned}
W_{MM_{2-5}}={}&\frac{x^3y^3}{z^2}+\frac{x^2y^3}{z^2}+\frac{3x^2y^2}{z}+\frac{x^2y}{z}+\frac{2xy^2}{z}+3xy+\frac{4xy}{z}+x+y+\frac{2y}{z}+z \\
&\quad +\frac{1}{y}+\frac{2}{x}+\frac{3z}{xy}+\frac{3}{xy}+\frac{1}{x^2y}+\frac{3z}{x^2y^2}+\frac{z}{x^3y^3}.
\end{aligned}}
Set
\[
S=\frac{xy}{z}+1+\frac{1}{xy}.
\]
A direct regrouping of the Laurent monomials gives the compact form
\[
W_{MM_{2-5}}=yS^2+xS+zS^3+\frac{xy}{z}-6.
\]
To compute the critical locus in the torus, introduce the torus coordinates
\[
a=\frac{xy}{z},\qquad b=\frac{1}{xy},
\]
so that $S=a+b+1$. The logarithmic critical equations show that the branch $S=0$ has no valid torus solution. The remaining branch gives
\[
a=-\frac{b^2y^2+by^2-1}{by^2}.
\]
Indeed, on the branch $S=0$ the equations contain
\[
(b+1)^2(by+1)=0,
\qquad
b(b+1)=0,
\]
which forces $b=-1$ and then $a=0$, outside the torus. After substituting the
displayed expression for $a$ on the remaining branch, the other two equations
factor as
\[
G(b,y)H(b,y)=0,
\qquad
G(b,y)(3b^2y^2-1)=0,
\]
where
\[
G(b,y)=b^3y^3+b^2y^3-by-1
\]
and
\[
H(b,y)=b^3y^3+b^2y^3+3b^2y^2+3by^2-by-2.
\]
Thus the critical locus splits into a one-dimensional component and an
isolated component. The one-dimensional component is cut out by
\[
b^3y^3+b^2y^3-by-1=0,
\]
or, in the original coordinates, by the parametrization
\[
\Gamma=\left\{\left(t,\frac{t^3+t^2-1}{t},-\frac{(t^3+t^2-1)^2}{t^3}\right):
 t\in\mathbb{C}^*,\ t^3+t^2-1\neq0\right\}.
\]
Along this component the potential is constant:
\[
W_{MM_{2-5}}|_{\Gamma}=-6.
\]
This component has a nondegenerate transverse quadratic part. Indeed, set
\[
f_1=xy-x^3-x^2+1,
\qquad
f_2=xz+y^2,
\]
so that $I_{\Gamma}=(f_1,f_2)$ on the open set parametrized above. With
$t=x|_{\Gamma}$ and $D=t^3+t^2-1$, put
\[
A_{\Gamma}=\mathbb C[t^{\pm1},D^{-1}].
\]
The substitutions
\[
x=t,
\qquad
y=\frac{U+D}{t},
\qquad
z=\frac{V-(U+D)^2/t^2}{t}
\]
identify $U=f_1$ and $V=f_2$ and give a global splitting
\[
\widehat X_{\Gamma}\simeq\operatorname{Spf}(A_{\Gamma}[[U,V]]).
\]
In these completed normal directions,
\[
W_{MM_{2-5}}+6\equiv A(t)f_1^2+B(t)f_1f_2+C(t)f_2^2
\pmod{I_{\Gamma}^3},
\]
where
\[
\begin{aligned}
A(t)&=-\frac{2}{t(t^3+t^2-1)},\\
B(t)&=\frac{t^3(4t+3)}{(t^3+t^2-1)^3},\\
C(t)&=-\frac{t^7(t+1)(2t+1)}{(t^3+t^2-1)^5}.
\end{aligned}
\]
The discriminant of this binary quadratic normal form is
\[
4A(t)C(t)-B(t)^2=-\frac{t^6}{(t^3+t^2-1)^6},
\]
which is the negative of the square unit $(t^3/D^3)^2$ in $A_{\Gamma}$. Over
$\mathbb C$ the quadratic form is therefore globally hyperbolic. The relative
formal Morse lemma removes the higher-order terms and gives a split formal
quadratic model over all of $\Gamma$. Hence the curve contributes
$\MBcat{\Gamma}$. In particular, the two rational points $(1,1,-1)$ and
$(-1,1,1)$ lie on this non-isolated component; they are not isolated Morse
points.

On the isolated component, a lexicographic Gr\"obner basis is
\[
575b+234y-652,
\qquad
78y^2-18y-23.
\]
It yields exactly the two critical points
\[
p_{\pm}=\left(\pm\sqrt{3},\frac{9\pm25\sqrt{3}}{78},\frac{243\pm77\sqrt{3}}{234}\right),
\]
with critical values
\[
W_{MM_{2-5}}(p_{\pm})=21\pm6\sqrt{3}.
\]
A direct computation of the Hessian determinant at $p_{+}$ and $p_{-}$ gives
a nonzero value, so both isolated singularities are Morse. Since the three
critical values $-6$ and $21\pm6\sqrt3$ are distinct, the corresponding
critical-fiber categories decompose as
\[
\overline{MF((\mathbb{C}^*)^3,W_{MM_{2-5}})}
\cong \MBcat{\Gamma} \oplus \MorseCat{3}\oplus \MorseCat{3}.
\]
\end{example}

\begin{example}
Consider the Landau--Ginzburg model $MM_{2-6}$ associated with a
$(2,2)$-divisor in $\mathbb{P}^2\times\mathbb{P}^2$. On the algebraic torus
$(\mathbb{C}^*)^3$, set
\[
B(z)=\frac{(z+1)^2}{z},\qquad v=yz.
\]
The Laurent potential can be written in the form
\[
W_{MM_{2-6}}
=x\left(1+\frac{1}{yz}\right)^2+yzB(z)+\frac{B(z)^2}{x}
+2\bigl(B(z)-2\bigr)+\frac{2B(z)}{yz}.
\]
Equivalently,
\[
W_{MM_{2-6}}
=x\left(\frac{v+1}{v}\right)^2+
\frac{B(z)^2}{x}+B(z)\left(\frac{v^2+2v+2}{v}\right)-4.
\]
The logarithmic critical equations in the variables $(x,v,z)$ have, after
clearing the invertible denominators, the form
\begin{align*}
&\bigl(vxz-v(z+1)^2+xz\bigr)\bigl(vxz+v(z+1)^2+xz\bigr)=0,\\
&v^3(z+1)^2-2vxz-2v(z+1)^2-2xz=0,\\
&(z^2-1)\bigl(xz(v^2+2v+2)+2v(z+1)^2\bigr)=0.
\end{align*}
In the torus these equations split into one isolated point and two
one-dimensional components. The isolated point is
\[
P_1=\left(\frac{8}{3},2,1\right),
\qquad
W_{MM_{2-6}}(P_1)=28.
\]
A direct computation gives $\det H_W(P_1)=108\neq0$, so $P_1$ is an isolated
Morse singularity of type $A_1$ and contributes
\[
\overline{MF^{P_1}_{28}((\mathbb{C}^*)^3,W_{MM_{2-6}})}
\cong \MorseCat{3}.
\]

The non-isolated critical locus lies in the fiber $W_{MM_{2-6}}=-4$ and has two
components:
\[
\Gamma_1=\{(t,1,-1):t\in\mathbb{C}^*\},
\]
and
\[
\Gamma_2=\left\{\left(2\frac{(t+1)^2}{t},-\frac{2}{t},t\right):
t\in\mathbb{C}^*,\ t\neq -1\right\}.
\]
Both satisfy $W_{MM_{2-6}}|_{\Gamma_i}=-4$. Their intersection in the torus is
empty. Indeed, solving the parametrized equations would require simultaneously
$t=-1$ from the $z$-coordinate and $1=-2/t$, which is impossible. Equivalently,
in the Laurent coordinate ring,
\[
(y-1,z+1)+(yz+2,xz-2(z+1)^2)=(1).
\]
Thus there is no additional intersection-point contribution, and the two
curve supports are genuinely disjoint.

It remains to record the transverse singularity type along the two curves. For
$\Gamma_1$, put
\[
A_1=\mathbb C[t^{\pm1}],
\qquad
Y=y-1,
\qquad
Z=z+1,
\qquad
t=x.
\]
These coordinates give a global identification
\[
\widehat X_{\Gamma_1}\simeq\operatorname{Spf}(A_1[[Y,Z]]).
\]
Indeed, the inverse substitution is $x=t$, $y=1+Y$, and $z=-1+Z$, whose
constant terms in the normal variables are units on the torus. In these
coordinates,
\[
W_{MM_{2-6}}+4\equiv tY^2-2tYZ+(t+1)Z^2
\pmod{(Y,Z)^3}.
\]
The determinant of the relative Hessian is
\[
\det H_{Y,Z}=4t,
\]
which is a unit in $A_1$. Its square class is $[t]\in
A_1^*/A_1^{*2}$, and it is nontrivial because the exponent of $t$ is odd.
Thus the relative formal Morse lemma identifies the completed potential with
a nondegenerate quadratic bundle whose Clifford twist must be retained.

For $\Gamma_2$, put
\[
A_2=\mathbb C[t^{\pm1},(t+1)^{-1}]
\]
and use the global coordinates
\[
R=\frac{yz+2}{z},
\qquad
S=\frac{xz-2(z+1)^2}{z},
\qquad
t=z.
\]
Their inverse substitution is
\[
z=t,
\qquad
y=R-\frac{2}{t},
\qquad
x=S+2\frac{(t+1)^2}{t},
\]
so they give a global identification
\[
\widehat X_{\Gamma_2}\simeq\operatorname{Spf}(A_2[[R,S]]).
\]
Modulo $(R,S)^3$, one obtains
\[
W_{MM_{2-6}}+4\equiv
-\frac{t\bigl((t+1)^2R+S\bigr)\bigl(3(t+1)^2R-S\bigr)}{8(t+1)^2}
\pmod{(R,S)^3}.
\]
The determinant of the relative Hessian is
\[
\det H_{R,S}=-\frac{t^2}{4},
\]
which is a square unit over $\mathbb C$. Thus the relative formal Morse lemma
gives an untwisted global quadratic model along $\Gamma_2$. By
Proposition~\ref{prop:morse_bott_clifford},
\[
\begin{aligned}
\overline{MF^{\Gamma_1}_{-4}((\mathbb C^*)^3,W_{MM_{2-6}})}
&\simeq \MBtwcat{\Gamma_1}{t},\\
\overline{MF^{\Gamma_2}_{-4}((\mathbb C^*)^3,W_{MM_{2-6}})}
&\simeq \MBcat{\Gamma_2}.
\end{aligned}
\]
Since the two curves are disjoint and the isolated point lies in the distinct
critical fiber over $28$, the full decomposition is
\[
\overline{MF((\mathbb{C}^*)^3,W_{MM_{2-6}})}
\cong
\MorseCat{3}\oplus \MBtwcat{\Gamma_1}{t}\oplus \MBcat{\Gamma_2}.
\]
\end{example}

\begin{example}
Consider another Laurent potential for the Mori--Mukai model $MM_{2-7}$ on the
algebraic torus $(\mathbb{C}^*)^3$. Set
\[
K=x(y+1)^2,
\qquad
S=1+\frac{1}{xy}.
\]
The potential has the compact form
\[
\widetilde W_{MM_{2-7}}=K\left(S^2+\frac{1}{z}\right)+zS^3+S-5.
\]
Equivalently,
\[
\widetilde W_{MM_{2-7}}
=x(y+1)^2\left(S^2+\frac{1}{z}\right)+zS^3+S-5,
\qquad S=1+\frac{1}{xy}.
\]
The differential is
\[
d\widetilde W_{MM_{2-7}}
=\left(S^2+\frac{1}{z}\right)dK
+\left(2KS+3zS^2+1\right)dS
+\left(S^3-\frac{K}{z^2}\right)dz.
\]
Moreover,
\[
\det\frac{\partial(K,S)}{\partial(x,y)}=\frac{(y-1)(y+1)}{xy^2}.
\]
The branch $y=-1$ gives no torus critical point. On the branch $y=1$, the
isolated critical points are
\[
p_{\alpha}=\left(\alpha,1,-28\alpha^2+56\alpha-24\right),
\qquad
4\alpha^3-4\alpha^2-4\alpha+3=0.
\]
The cubic has three distinct roots. The remaining solution on the $y=1$ branch,
$(-3/4,1,-9)$, lies on the non-isolated component below and is not counted as
an isolated point. The critical values
$w_{\alpha}=\widetilde W_{MM_{2-7}}(p_{\alpha})$ are the three roots of
\[
w^3-20w^2-96w+148=0.
\]
This polynomial is square-free and takes the value $3$ at $w=-5$. Hence the
three isolated critical values are distinct from one another and from the
critical value of the curve below.
A direct Hessian computation gives $\det H_W(p_{\alpha})\neq0$ for all three
roots $\alpha$, so these are isolated Morse singularities of type $A_1$.

On the complementary branch, the equations
\[
S^2+\frac{1}{z}=0,
\qquad
2KS+3zS^2+1=0,
\qquad
S^3-\frac{K}{z^2}=0
\]
give
\[
z=-\frac{1}{S^2},
\qquad
K=\frac{1}{S}.
\]
In the original torus coordinates this critical curve is
\[
\Gamma=\left\{\left(-\frac{t^2+t+1}{t(t+1)^2},\,t,\,-\frac{(t^2+t+1)^2}{t^2}\right):
t\in\mathbb{C}^*,\ t\neq -1,\ t^2+t+1\neq0\right\}.
\]
Along this curve,
\[
\widetilde W_{MM_{2-7}}|_{\Gamma}=-5.
\]
The point $(-3/4,1,-9)$ corresponds to $t=1$ in this parametrization.

It remains to classify the transverse formal type along $\Gamma$. Put
\[
D(t)=t^2+t+1,
\qquad
A_{\Gamma}=\mathbb C[t^{\pm1},(t+1)^{-1},D(t)^{-1}],
\]
and use the normal generators
\[
F_1=xy(y+1)^2+y^2+y+1,
\qquad
F_2=y^2z+(y^2+y+1)^2.
\]
Thus $I_{\Gamma}=(F_1,F_2)$ on the displayed open curve. The substitutions
\[
y=t,
\qquad
x=\frac{F_1-D(t)}{t(t+1)^2},
\qquad
z=\frac{F_2-D(t)^2}{t^2}
\]
give explicit inverse formal coordinates and identify the completed
neighborhood globally as
\[
\widehat X_{\Gamma}\simeq
\operatorname{Spf}(A_{\Gamma}[[F_1,F_2]]).
\]
The constant terms of $x$ and $z$ in the normal variables are units because
$t$, $t+1$, and $D(t)$ are invertible in $A_{\Gamma}$. Modulo
$(F_1,F_2)^3$, one has
\[
\widetilde W_{MM_{2-7}}+5\equiv
\frac{\bigl(D(t)(t^2+3t+1)F_1-tF_2\bigr)
\bigl(2D(t)(t+1)^2F_1-tF_2\bigr)}{tD(t)^5}.
\]
The determinant of the relative Hessian of this quadratic form is
\[
-\frac{1}{D(t)^6},
\]
which is a unit in $A_{\Gamma}$. The relative formal Morse lemma therefore
identifies the completed potential with a nondegenerate quadratic bundle over
the whole curve. Any resulting Clifford or super-Brauer twist is included in
the notation $\MBcat{\Gamma}$.

By Proposition~\ref{prop:morse_bott_clifford}, the curve contribution is
\[
\overline{MF^{\Gamma}_{-5}((\mathbb{C}^*)^3,\widetilde W_{MM_{2-7}})}
\cong \MBcat{\Gamma}.
\]
The three isolated Morse points contribute three copies of $\MorseCat{3}$.
Consequently,
\[
\overline{MF((\mathbb{C}^*)^3,\widetilde W_{MM_{2-7}})}
\cong \MBcat{\Gamma}\oplus\bigoplus_{\alpha:\,4\alpha^3-4\alpha^2-4\alpha+3=0}\MorseCat{3}.
\]
\end{example}

\begin{example}
Consider the Landau--Ginzburg model $MM_{2-8}$ on the algebraic torus
$(\mathbb{C}^*)^3$. Set
\[
u=y+z+1,
\qquad
v=1+\frac{1}{yz}.
\]
The Laurent potential is
\[
W_{MM_{2-8}}=xu^2+2uv+\frac{v^2}{x}-2.
\]
Equivalently, if
\[
F=xu+v=x(y+z+1)+1+\frac{1}{yz},
\]
then the potential is the exact square
\[
W_{MM_{2-8}}=\frac{F^2}{x}-2.
\]
The ordinary derivatives factor as
\begin{align*}
\frac{\partial W_{MM_{2-8}}}{\partial x}
&=\frac{F}{x^2}\,(xu-v),\\
\frac{\partial W_{MM_{2-8}}}{\partial y}
&=2F\left(1-\frac{1}{xy^2z}\right),\\
\frac{\partial W_{MM_{2-8}}}{\partial z}
&=2F\left(1-\frac{1}{xyz^2}\right).
\end{align*}
Hence the critical locus has a non-isolated branch and an isolated branch.

First, the equation $F=0$ defines the critical surface
\[
\Sigma=\left\{(x,y,z)\in(\mathbb{C}^*)^3:
x(y+z+1)+1+\frac{1}{yz}=0\right\}.
\]
This surface lies entirely in the singular fiber
\[
W_{MM_{2-8}}|_{\Sigma}=-2.
\]
It is smooth: if $dF=0$, then $y+z+1=0$; on $F=0$ this also forces
$1+1/(yz)=0$, and the remaining equations $F_y=F_z=0$ would force $y=z$, which
is incompatible with these two relations in the torus. The formal normal form
along $\Sigma$ is exact:
\[
W_{MM_{2-8}}+2=\frac{F^2}{x}.
\]
Since $x$ is a unit on the torus, this is a transverse $A_1$ normal form along
the whole surface $\Sigma$. The global generator $F$ is a basis of
$I_\Sigma/I_\Sigma^2$, hence
$\widehat X_\Sigma\simeq\operatorname{Spf}(\mathcal O_\Sigma[[F]])$.
Moreover, the formal zero fiber is defined by $F^2$, so
$\Sing(\operatorname{Spec}B)\subseteq V((F))$ and
Theorem~\ref{thm:general_formal} applies to the unrestricted formal category.

On the complementary branch $F\neq0$, the equations are
\[
xy^2z=1,
\qquad
xyz^2=1,
\qquad
xu-v=0.
\]
The first two equations give $y=z$ and $x=y^{-3}$. Substitution into $xu=v$
gives
\[
y^3-y-1=0.
\]
Thus the isolated critical points are
\[
p_{\rho}=\left(\rho^{-3},\rho,\rho\right),
\qquad
\rho^3-\rho-1=0.
\]
The cubic has three distinct roots. At these points the Hessian determinant is
nonzero; explicitly, after reduction modulo $\rho^3-\rho-1$, one obtains
\[
\det H_W(p_{\rho})=\frac{32(51\rho^2+64\rho+33)}{\rho(\rho+1)},
\]
and the numerator is coprime to $\rho^3-\rho-1$. Hence these three points are
isolated Morse singularities of type $A_1$. Their critical values are the roots
of
\[
w^3-2w^2-292w-2168=0.
\]

By Proposition~\ref{prop:Ade_category_model}, the surface contribution is
\[
\overline{MF^{\Sigma}_{-2}((\mathbb{C}^*)^3,W_{MM_{2-8}})}
\cong \MBcat{\Sigma},
\]
and each isolated Morse point contributes one copy of $\MorseCat{3}$.
Consequently,
\[
\overline{MF((\mathbb{C}^*)^3,W_{MM_{2-8}})}
\cong \MBcat{\Sigma}\oplus \bigoplus_{\rho:\,\rho^3-\rho-1=0}\MorseCat{3}.
\]
\end{example}

\begin{example}
Consider the Mori--Mukai model $MM_{2-9}$ on the algebraic torus
$(\mathbb{C}^*)^3$. Use the torus coordinates
\[
u=yz,
\qquad
v=1+\frac{1}{u},
\qquad
A=1+\frac{1}{z}.
\]
The Laurent potential has the compact form
\[
W_{MM_{2-9}}
=xA+\frac{uv^3}{x}+zA^2uv^2-\frac{1}{u}-4.
\]
Equivalently, in the coordinates $(x,u,z)$,
\[
W_{MM_{2-9}}
=\frac{x(z+1)}{z}
+\frac{(u+1)^3}{u^2x}
+\frac{(z+1)^2(u+1)^2}{zu}
-\frac{1}{u}-4.
\]

Let $\theta_x=x\partial_x$, $\theta_u=u\partial_u$, and
$\theta_z=z\partial_z$. The torus critical locus is obtained from the ideal
generated by $\theta_xW_{MM_{2-9}}$, $\theta_uW_{MM_{2-9}}$, and
$\theta_zW_{MM_{2-9}}$, after clearing denominators and saturating by the
torus monomial $xuz$. This gives the reduced zero-dimensional ideal
\fitdisplay{\left(
19x+49z^3-115z^2-36z+88,
u-z^3+3z^2+z-5,
z^4-2z^3-2z^2+2z+2
\right)}
Thus, in the $(x,u,z)$ coordinates, the critical points are
\[
\widetilde p_{\zeta}=\left(
-\frac{49\zeta^3-115\zeta^2-36\zeta+88}{19},
\zeta^3-3\zeta^2-\zeta+5,
\zeta
\right),
\qquad
q(\zeta)=0.
\]
Returning to the original $(x,y,z)$ coordinates, where $y=u/z$, these points
are
\[
p_{\zeta}=\left(
-\frac{49\zeta^3-115\zeta^2-36\zeta+88}{19},
\frac{\zeta^3-3\zeta^2-\zeta+5}{\zeta},
\zeta
\right),
\qquad
q(\zeta)=0,
\]
where
\[
q(z)=z^4-2z^3-2z^2+2z+2.
\]
The polynomial $q$ is square-free, so this gives four torus critical points.
The denominator conditions are automatic on this locus: $q$ is coprime to
$z$, $z^3-3z^2-z+5$, and $49z^3-115z^2-36z+88$.

The critical values $w_{\zeta}=W_{MM_{2-9}}(p_{\zeta})$ are the four roots of
\[
w^4-4w^3-202w^2-1280w-2207=0.
\]
Finally, the logarithmic Hessian is nondegenerate at each critical point. After
reduction modulo $q(\zeta)$, consider the numerator of
\[
\det(\theta_i\theta_jW_{MM_{2-9}})(\widetilde p_{\zeta}).
\]
It is proportional to
\[
9042\zeta^3-24194\zeta^2-3547\zeta+25168,
\]
which is coprime to $q(\zeta)$. Therefore all four singularities are isolated
Morse singularities of type $A_1$. Consequently,
\[
\overline{MF((\mathbb{C}^*)^3,W_{MM_{2-9}})}
\cong \bigoplus_{\zeta:\,q(\zeta)=0}\MorseCat{3}.
\]
\end{example}

\begin{example}
Consider the Mori--Mukai model $MM_{2-10}$ on the algebraic torus
$(\mathbb{C}^*)^3$, defined by the Laurent potential
\begin{align*}
W_{MM_{2-10}}={}&x+yz^2+2yz+y+4z+\frac{2yz}{x}
+\frac{2y}{x}+\frac{4}{x}+\frac{y}{x^2} \\
&+\frac{6}{y}+\frac{2}{yz}+\frac{4}{y^2z}
+\frac{2}{xyz}+\frac{1}{y^3z^2}.
\end{align*}
This apparently irregular expression is an exact square completion. Put
\[
F=yz+y+\frac{y}{x}+2+\frac{1}{yz}.
\]
Then the potential has the compact global form
\[
W_{MM_{2-10}}=\frac{F^2}{y}+x-4.
\]

The ordinary derivatives are
\begin{align*}
\frac{\partial W_{MM_{2-10}}}{\partial x}
&=\frac{2F}{y}\frac{\partial F}{\partial x}+1,\\
\frac{\partial W_{MM_{2-10}}}{\partial y}
&=\frac{F}{y^2}\left(2y\frac{\partial F}{\partial y}-F\right),\\
\frac{\partial W_{MM_{2-10}}}{\partial z}
&=\frac{2F}{y}\frac{\partial F}{\partial z}.
\end{align*}
The first equation excludes the branch $F=0$, since then it would read
$1=0$. Hence any torus critical point satisfies
\[
\frac{\partial F}{\partial z}=y-\frac{1}{yz^2}=0,
\qquad\text{so}\qquad
yz=\pm1.
\]
The branch $yz=-1$ is empty: on it, $F=y(1+1/x)$, and the $y$-equation forces
$F=0$, contradicting the $x$-equation. On the branch $yz=1$, one has
\[
F=y\left(1+\frac{1}{x}\right)+4.
\]
The $y$-equation gives
\[
y\left(1+\frac{1}{x}\right)=4,
\qquad\text{hence}\qquad
F=8.
\]
Substituting this into the $x$-equation gives
\[
-\frac{16}{x^2}+1=0,
\qquad\text{so}\qquad
x=\pm4.
\]
Therefore the torus critical locus consists of precisely two points,
\[
P_1=\left(4,\frac{16}{5},\frac{5}{16}\right),
\qquad
P_2=\left(-4,\frac{16}{3},\frac{3}{16}\right).
\]
Their critical values are
\[
W_{MM_{2-10}}(P_1)=20,
\qquad
W_{MM_{2-10}}(P_2)=4.
\]

Equivalently, after clearing denominators and saturating by the torus monomial
$xyz$, the critical ideal is generated by
\[
x-64z+16,
\qquad
15y+256z-128,
\qquad
(16z-5)(16z-3),
\]
which confirms that there are no hidden positive-dimensional components in the
torus. The Hessian determinants at the two critical points are
\[
\det H_W(P_1)=50,
\qquad
\det H_W(P_2)=-18.
\]
Thus both singularities are isolated Morse singularities of type $A_1$.
Consequently,
\[
\overline{MF((\mathbb{C}^*)^3,W_{MM_{2-10}})}
\cong \MorseCat{3}\oplus \MorseCat{3}.
\]
\end{example}

\begin{example}
Consider the Mori--Mukai model $MM_{2-11}$ on the algebraic torus
$(\mathbb{C}^*)^3$, with Laurent potential
\begin{align*}
W_{MM_{2-11}}={}&x+\frac{x}{yz}+yz+y+z+\frac{2}{z}+\frac{2}{y}
+\frac{yz}{x}+\frac{2y}{x}+\frac{2z}{x} \\
&+\frac{2}{x}+\frac{y}{xz}+\frac{z}{xy}.
\end{align*}
Set
\[
S=y+z+yz.
\]
The potential admits the compact symmetric form
\[
W_{MM_{2-11}}
=x\left(1+\frac{1}{yz}\right)
+S\left(1+\frac{2}{yz}\right)
+\frac{S^2}{xyz}-2.
\]
Indeed,
\[
\frac{S^2}{xyz}
=\frac{1}{x}\left(yz+2y+2z+2+\frac{y}{z}+\frac{z}{y}\right),
\]
which accounts for all the $x^{-1}$ terms.

We now compute the torus critical locus. On the open branch $y\neq z$, put
\[
p=yz,
\qquad
q=y+z,
\qquad
S=p+q.
\]
Since the Jacobian determinant of $(y,z)\mapsto(p,q)$ is $z-y$, the equations
$\partial_yW=\partial_zW=0$ are equivalent on this branch to
$\partial_pW=\partial_qW=0$. Eliminating from
\[
\partial_xW=\partial_pW=\partial_qW=0
\]
gives the consequence $p^2S^2=0$. Since $p=yz$ is a torus unit, this forces
$S=0$. Then $\partial_qW=0$ gives $p=-2$, while
$\partial_xW=1+1/p=1/2$, a contradiction. Hence no critical point occurs on
the branch $y\neq z$.

It remains to impose the symmetric branch $y=z=t$. The two independent
critical equations, after clearing denominators, are
\[
-t^4-4t^3+t^2x^2-4t^2+x^2=0,
\]
and
\[
t^4x+t^4+t^3x+2t^3-2tx-x^2=0.
\]
Saturating by the torus monomial $xt$ gives the reduced ideal
\[
\left(x-t(t+1),\ t^4+2t^3+t^2-2t-3\right).
\]
Equivalently, the first equation may be written as
\[
x^2(t^2+1)=t^2(t+2)^2,
\]
and the remaining equation selects the branch $x=t(t+1)$ together with the
quartic relation
\[
q(t)=t^4+2t^3+t^2-2t-3=0.
\]
The polynomial $q$ is irreducible over $\mathbb{Q}$ and square-free
($\operatorname{disc}(q)=-1856$). Therefore the critical locus consists of the
four torus points
\[
P_{\tau}=\left(\tau(\tau+1),\tau,\tau\right),
\qquad
q(\tau)=0.
\]
The torus conditions are automatic, since $q$ is coprime to $t$ and $t+1$.

The corresponding critical values are
\[
W_{MM_{2-11}}(P_{\tau})
=\frac{8\tau^3+25\tau^2+13\tau-2}{\tau^3+\tau+1},
\]
or, equivalently, they are the four roots of
\[
w^4-2w^3-129w^2-1086w-1823=0.
\]
Finally, the Hessian determinant is nonzero at each of these points. More
precisely, after reduction modulo $q(\tau)$, one obtains
\[
\det H_W(P_{\tau})
=2\frac{324\tau^3+655\tau^2+1239\tau+1329}
{103\tau^3+160\tau^2-17\tau-162},
\]
and both the numerator and denominator are coprime to $q(\tau)$. Thus all four
singularities are isolated Morse singularities of type $A_1$. Consequently,
\[
\overline{MF((\mathbb{C}^*)^3,W_{MM_{2-11}})}
\cong \bigoplus_{\tau:\,q(\tau)=0}\MorseCat{3}.
\]
\end{example}

\begin{example}
Consider the Mori--Mukai model $MM_{2-12}$ on the algebraic torus
$(\mathbb{C}^*)^3$, with Laurent potential
\begin{align*}
W_{MM_{2-12}}={}&x+yz^2+2yz+y+2z+\frac{2}{x}
+\frac{2}{xz}+\frac{2}{z}+\frac{1}{xz^2} \\
&+\frac{1}{y}+\frac{2}{xyz}+\frac{1}{x^2yz^2}.
\end{align*}
Set
\[
A=z+1,
\qquad
B=1+\frac{1}{xz},
\qquad
F=yA+B.
\]
Then the potential has the exact square-completed form
\[
W_{MM_{2-12}}
=yA^2+2AB+\frac{B^2}{y}+xB^2-2
=\frac{F^2}{y}+xB^2-2.
\]

The $y$-derivative factors as
\[
\frac{\partial W_{MM_{2-12}}}{\partial y}
=\frac{F}{y^2}\left(2y\frac{\partial F}{\partial y}-F\right)
=\frac{F}{y^2}\left(y(z+1)-B\right).
\]
Thus the critical locus splits into two branches.

More directly, saturating the logarithmic critical ideal by the torus monomial
$xyz$ gives
\[
\left(x-z-2,\ (3y-2)(z+1),\ (z-1)(z+1)\right).
\]
Consequently, the full torus critical locus is the disjoint union of the curve
\[
\Gamma=\{(1,t,-1):t\in\mathbb{C}^*\}.
\]
and the isolated point
\[
P_1=\left(3,\frac{2}{3},1\right).
\]
There are no further critical components.

Along the curve, $F=B=0$, so
\[
W_{MM_{2-12}}|_{\Gamma}=-2.
\]
Put
\[
A_{\Gamma}=\mathbb C[t^{\pm1}],
\qquad
X=x-1,
\qquad
Z=z+1,
\qquad
t=y.
\]
These coordinates give a global identification
\[
\widehat X_{\Gamma}\simeq\operatorname{Spf}(A_{\Gamma}[[X,Z]]).
\]
Modulo $(X,Z)^3$, the shifted potential has quadratic part
\[
W_{MM_{2-12}}+2\equiv
\frac{t+1}{t}X^2-\frac{2}{t}XZ+
\frac{t^2-t+1}{t}Z^2.
\]
The determinant of its relative Hessian is
\[
4t,
\]
which is a unit in $A_{\Gamma}$ but has the nontrivial square class
$[t]\in A_{\Gamma}^*/A_{\Gamma}^{*2}$. Hence the relative formal Morse lemma
gives a global quadratic formal model with the corresponding Clifford twist.

At the isolated point, $B=4/3$ and $F=8/3$, and therefore
\[
W_{MM_{2-12}}(P_1)=14.
\]
The Hessian determinant at the isolated point is nonzero; explicitly,
\[
\det H_W(P_1)=64.
\]
Thus $P_1$ is an isolated Morse singularity of type $A_1$.

By Proposition~\ref{prop:morse_bott_clifford}, the curve contribution is
\[
\overline{MF^{\Gamma}_{-2}((\mathbb{C}^*)^3,W_{MM_{2-12}})}
\cong \MBtwcat{\Gamma}{t},
\]
and the isolated point contributes $\MorseCat{3}$. Consequently,
\[
\overline{MF((\mathbb{C}^*)^3,W_{MM_{2-12}})}
\cong \MBtwcat{\Gamma}{t}\oplus \MorseCat{3}.
\]
\end{example}

\begin{example}
Consider the Mori--Mukai model $MM_{2-13}$ on the algebraic torus
$(\mathbb{C}^*)^3$, with Laurent potential
\begin{align*}
W_{MM_{2-13}}={}&x+\frac{x}{y}+y+z+\frac{1}{z}+\frac{2}{y}
+\frac{2}{yz}+\frac{1}{y^2z} \\
&+\frac{2yz}{x}+\frac{2y}{x}+\frac{2}{x}+\frac{y^2z}{x^2}.
\end{align*}
Set
\[
A=1+\frac{y}{x},
\qquad
B=1+\frac{1}{y},
\qquad
F=zA+B.
\]
Then the potential is the square completion
\[
W_{MM_{2-13}}
=zA^2+2AB+\frac{B^2}{z}+xB+y-2
=\frac{F^2}{z}+xB+y-2.
\]

The $z$-derivative factors as
\[
\frac{\partial W_{MM_{2-13}}}{\partial z}
=\frac{F}{z^2}\left(2z\frac{\partial F}{\partial z}-F\right)
=\frac{F}{z^2}\left(zA-B\right).
\]
Thus the critical locus has two branches.

First suppose $F=0$. Then the $x$-derivative reduces to $B=0$, hence
$y=-1$. The equation $F=0$ becomes $zA=0$, and since $z$ is a torus unit, one
gets $A=0$, hence $x=1$. Therefore this branch is the critical curve
\[
\Gamma=\{(1,-1,t):t\in\mathbb{C}^*\}.
\]
Along this curve, $F=B=0$, and therefore
\[
W_{MM_{2-13}}|_{\Gamma}=-3.
\]
Put
\[
A_{\Gamma}=\mathbb C[t^{\pm1}],
\qquad
X=x-1,
\qquad
Y=y+1,
\qquad
t=z.
\]
These coordinates give a global identification
\[
\widehat X_{\Gamma}\simeq\operatorname{Spf}(A_{\Gamma}[[X,Y]]).
\]
Modulo $(X,Y)^3$, the shifted potential has quadratic part
\[
W_{MM_{2-13}}+3\equiv
tX^2+(2t-3)XY+\frac{t^2-3t+1}{t}Y^2.
\]
Its relative Hessian is
\[
H_W^{\mathrm{norm}}=
\begin{pmatrix}
2t & 2t-3\\
2t-3 & \dfrac{2(t^2-3t+1)}{t}
\end{pmatrix},
\qquad
\det H_W^{\mathrm{norm}}=-5.
\]
Thus the relative Hessian is invertible over $A_{\Gamma}$, and the relative
formal Morse lemma gives a global quadratic formal model along $\Gamma$. Any
Clifford or super-Brauer twist is included in $\MBcat{\Gamma}$.

On the complementary branch $F\neq0$, the equation
$\partial_zW_{MM_{2-13}}=0$ gives
\[
zA=B,
\qquad\text{and hence}\qquad
F=2B.
\]
The $x$-derivative becomes
\[
-\frac{2Fy}{x^2}+B=0.
\]
Since $F\neq0$ implies $B\neq0$, this gives
\[
x^2=4y.
\]
Substituting this relation and $zA=B$ into the remaining equation gives
\[
x^4+4x^3-16x-64=(x+4)(x^3-16)=0.
\]
The factor $x+4$ is a spurious factor introduced by the rational reduction:
if $x=-4$, then $y=4$ and $A=1+y/x=0$, while $B=5/4$, contradicting
$zA=B$. Thus the torus solutions are exactly those with
\[
\xi^3=16.
\]
They are
\[
P_{\xi}=\left(\xi,\frac{\xi^2}{4},1\right),
\qquad
\xi^3=16.
\]
Equivalently, after clearing denominators and saturating by the torus monomial,
the full critical ideal is
\[
\left(x-y^2,\ (y+1)(y^3-4),\ (y+1)(z-1)\right),
\]
which displays the curve $y=-1$, $x=1$, together with the three isolated points
$y^3=4$, $z=1$, $x=y^2$.

The critical values at the isolated points are
\[
W_{MM_{2-13}}(P_{\xi})=\frac{3}{4}\xi^2+3\xi+2,
\]
or, equivalently, the three roots of
\[
w^3-6w^2-96w-332=0.
\]
The Hessian determinant at these points is nonzero. More precisely, modulo
$\xi^3-16$ one obtains
\[
\det H_W(P_{\xi})=
\frac{3(6\xi^2+17\xi+32)}{8\xi},
\]
and the numerator is coprime to $\xi^3-16$. Thus the three isolated points are
Morse singularities of type $A_1$.

By Proposition~\ref{prop:morse_bott_clifford}, the curve contribution is
\[
\overline{MF^{\Gamma}_{-3}((\mathbb{C}^*)^3,W_{MM_{2-13}})}
\cong \MBcat{\Gamma},
\]
and the three isolated Morse points contribute three copies of
$\MorseCat{3}$. Consequently,
\[
\overline{MF((\mathbb{C}^*)^3,W_{MM_{2-13}})}
\cong \MBcat{\Gamma}\oplus \bigoplus_{\xi:\,\xi^3=16}\MorseCat{3}.
\]
\end{example}

\begin{example}
Consider the Mori--Mukai model $MM_{2-14}$ on the algebraic torus
$(\mathbb{C}^*)^3$, with Laurent potential
\begin{align*}
W_{MM_{2-14}}={}&x+\frac{2xz}{y}+\frac{xz^2}{y^2}
+\frac{xz}{y^2}+y+z+\frac{2z}{y}+\frac{3}{y} \\
&+\frac{2y}{xz}+\frac{1}{x}+\frac{3}{xz}
+\frac{y}{x^2z^2}.
\end{align*}
Clearing denominators in the logarithmic critical equations
$\theta_xW=\theta_yW=\theta_zW=0$ and saturating by the torus monomial $xyz$
gives the reduced lexicographic basis
\[
\left(
13x-3z^3-4z^2-4z-22,
26y-7z^3-18z^2+60z-99,
r(z)
\right),
\]
where
\[
r(z)=z^4+2z^3-5z^2+14z-11.
\]
The polynomial $r$ is irreducible over $\mathbb{Q}$ and square-free
($\operatorname{disc}(r)=-1081600$). Hence the critical points are
\[
P_{\eta}=\left(
\frac{3\eta^3+4\eta^2+4\eta+22}{13},
\frac{7\eta^3+18\eta^2-60\eta+99}{26},
\eta
\right),
\qquad
r(\eta)=0.
\]
The torus conditions are automatic on this locus, since $r$ is coprime to
$z$ and to the numerators of the displayed $x$- and $y$-coordinates.

At these points the critical value is
\[
W_{MM_{2-14}}(P_{\eta})
=\frac{17\eta^3+53\eta^2-38\eta+168}{13}.
\]
The four critical values are therefore the roots of
\[
w^4-10w^3-69w^2-30w+209=0.
\]
The discriminant of this polynomial is $-11236000000$, so the four values are
distinct. The logarithmic Hessian determinant reduces modulo $r(\eta)$ to
\[
\det(\theta_i\theta_jW)(P_{\eta})
=22\eta^3+61\eta^2-47\eta+253.
\]
This polynomial is coprime to $r$, so all four critical points are isolated
Morse singularities of type $A_1$.
Consequently,
\[
\overline{MF((\mathbb{C}^*)^3,W_{MM_{2-14}})}
\cong \bigoplus_{\eta:\,r(\eta)=0}\MorseCat{3}.
\]
\end{example}

\begin{example}
Consider the Mori--Mukai model $MM_{2-15}$ on the algebraic torus
$(\mathbb{C}^*)^3$, with Laurent potential
\begin{align*}
W_{MM_{2-15}}={}&x+\frac{x}{z}+\frac{x}{yz}+y+\frac{y}{z}+z
+\frac{2}{z}+\frac{2}{y} \\
&+\frac{y}{xz}+\frac{2}{x}+\frac{z}{xy}.
\end{align*}
This potential is symmetric in $x$ and $y$. It may be grouped as
\[
W_{MM_{2-15}}
=y+\frac{y}{z}+z+\frac{2}{z}+\frac{2}{y}
+x\left(\frac{yz+y+1}{yz}\right)
+\frac{(y+z)^2}{xyz}.
\]

After clearing denominators in the logarithmic critical equations and
saturating by the torus monomial $xyz$, one obtains the ideal
\fitdisplay{\left(
x+y-z^4+2z^3+z^2-2z-1,
(z+1)(y+z^3-z^2-3z-1),
(z+1)(z^4-3z^2-4z-1)
\right)}
This displays two branches. The first branch is obtained from $z=-1$ and
$x+y-1=0$, namely
\[
\Gamma=\{(t,1-t,-1):t\in\mathbb{C}^*,\ t\neq1\}.
\]
It lies in the single fiber
\[
W_{MM_{2-15}}|_{\Gamma}=-1.
\]
The transverse quadratic form is nondegenerate along $\Gamma$. Indeed, use the
global coordinates
\[
u=x+y+z,
\qquad
v=z+1,
\]
and take $x=t$ as a coordinate along the curve, so that
$y=1-t+u-v$ and $z=-1+v$. The normal Hessian of
$W_{MM_{2-15}}+1$ in the variables $(u,v)$ at $u=v=0$ is
\[
\begin{pmatrix}
\dfrac{2}{t(t-1)} & -1\\
-1 & 0
\end{pmatrix},
\qquad
\det=-1.
\]
More precisely, with
\[
A_{\Gamma}=\mathbb C[t^{\pm1},(t-1)^{-1}],
\]
the displayed substitutions identify the completed neighborhood globally as
\[
\widehat X_{\Gamma}\simeq\operatorname{Spf}(A_{\Gamma}[[u,v]]),
\]
and the quadratic expansion is
\[
W_{MM_{2-15}}+1\equiv \frac{u^2}{t(t-1)}-uv
\pmod{(u,v)^3}.
\]
The determinant $-1$ is a unit in $A_{\Gamma}$, so the relative formal Morse
lemma gives a global quadratic formal model along $\Gamma$. Any Clifford or
super-Brauer twist is included in $\MBcat{\Gamma}$.

The second branch is the symmetric branch $x=y$. The point of this branch that
also lies on $\Gamma$ is $(1/2,1/2,-1)$; after removing this curve component,
the isolated branch is
\[
x=y=\alpha,
\qquad
z=\frac{\alpha^2+2\alpha-1}{2},
\qquad
q(\alpha)=0,
\]
where
\[
q(\alpha)=\alpha^4+2\alpha^3-2\alpha^2-6\alpha+1.
\]
Equivalently, before removing the curve component one encounters
\[
2\alpha^5+3\alpha^4-6\alpha^3-10\alpha^2+8\alpha-1
=(2\alpha-1)q(\alpha),
\]
where the root $\alpha=1/2$ gives the point on $\Gamma$. The quartic $q$ is
square-free, with discriminant $-15616$, and its roots give four torus points
\[
P_{\alpha}=\left(\alpha,\alpha,\frac{\alpha^2+2\alpha-1}{2}\right),
\qquad
q(\alpha)=0.
\]
The torus conditions are automatic, since $q$ is coprime to $\alpha$ and to
$\alpha^2+2\alpha-1$.

The critical values at these four isolated points are the roots of
\[
w^4+4w^3-120w^2-816w-688=0.
\]
This polynomial is square-free and does not vanish at $w=-1$, so the isolated
points do not lie in the critical fiber containing $\Gamma$.
At the level of the quartic field one may write
\[
W_{MM_{2-15}}(P_{\alpha})
=\frac{2\alpha(7\alpha^2+18\alpha-3)}{\alpha^2+6\alpha-1}.
\]
The Hessian determinant is nonzero at every isolated point. More precisely,
modulo $q(\alpha)$ one obtains
\[
\det H_W(P_{\alpha})
=2\frac{2391\alpha^3+5550\alpha^2+995\alpha-310}
{791\alpha^3+1242\alpha^2+1949\alpha-346},
\]
and the numerator and denominator are both coprime to $q(\alpha)$. Hence all
four isolated singularities are Morse singularities of type $A_1$.

By Proposition~\ref{prop:morse_bott_clifford}, the curve contribution is
\[
\overline{MF^{\Gamma}_{-1}((\mathbb{C}^*)^3,W_{MM_{2-15}})}
\cong \MBcat{\Gamma},
\]
and the isolated Morse points contribute four copies of $\MorseCat{3}$.
Consequently,
\[
\overline{MF((\mathbb{C}^*)^3,W_{MM_{2-15}})}
\cong \MBcat{\Gamma}\oplus \bigoplus_{\alpha:\,q(\alpha)=0}\MorseCat{3}.
\]
\end{example}

\begin{example}
Consider the Mori--Mukai model $MM_{2-16}$ on the algebraic torus
$(\mathbb{C}^*)^3$, with Laurent potential
\begin{align*}
W_{MM_{2-16}}={}&x+y+z+\frac{1}{z}+\frac{2z}{y}+\frac{2}{y}
+\frac{z}{y^2}+\frac{2y}{x} \\
&+\frac{2y}{xz}+\frac{2}{x}+\frac{y^2}{x^2z}.
\end{align*}
Set
\[
A=1+\frac{1}{y},
\qquad
B=1+\frac{y}{x},
\qquad
F=zA+B.
\]
Then the potential is the exact square completion
\[
W_{MM_{2-16}}
=zA^2+2AB+\frac{B^2}{z}+x+y-2
=\frac{F^2}{z}+x+y-2.
\]

The $z$-derivative factors as
\[
\frac{\partial W_{MM_{2-16}}}{\partial z}
=\frac{F}{z^2}\left(2z\frac{\partial F}{\partial z}-F\right)
=\frac{F}{z^2}(zA-B).
\]
The branch $F=0$ is empty in the torus: on this branch the $x$-derivative
would become $1=0$. Thus every critical point lies on the branch
\[
zA=B,
\qquad\text{so}\qquad
F=2B.
\]
Substituting this into the $x$-derivative gives
\[
y=\frac{x^2}{4}-1.
\]
The remaining equation also gives
\[
y=\frac{4}{x}.
\]
Equating these two expressions yields the cubic
\[
q(x)=x^3-4x-16=0.
\]
Equivalently, clearing denominators and saturating by the torus monomial gives
the reduced critical ideal
\[
\left(
2x-48z^2+125z-76,
2y+16z^2-31z+12,
16z^3-63z^2+80z-32
\right).
\]
On the cubic field generated by a root $\beta$ of $q$, the critical points can
be written as
\[
P_{\beta}=\left(
\beta,
\frac{\beta^2}{4}-1,
\frac{92-4\beta-3\beta^2}{64}
\right),
\qquad
\beta^3-4\beta-16=0.
\]
Using the cubic relation, the second coordinate is also $4/\beta$. The torus
conditions are automatic: $q$ is coprime to $\beta$, to $\beta^2-4$, and to
$92-4\beta-3\beta^2$. The cubic is square-free, with discriminant $-6656$.

The corresponding critical values have the simple form
\[
W_{MM_{2-16}}(P_{\beta})=3\beta+\frac{4}{\beta}+2,
\]
or, equivalently, they are the three roots of
\[
w^3-5w^2-64w-356=0.
\]
The Hessian determinant is nonzero at each critical point. More precisely,
modulo $q(\beta)$ one obtains
\[
\det H_W(P_{\beta})
=\frac{43\beta^2+128\beta+208}{4\beta(\beta+2)},
\]
and the numerator and denominator are both coprime to $q(\beta)$. Therefore
all three critical points are isolated Morse singularities of type $A_1$.
Consequently,
\[
\overline{MF((\mathbb{C}^*)^3,W_{MM_{2-16}})}
\cong \bigoplus_{\beta:\,\beta^3-4\beta-16=0}\MorseCat{3}.
\]
\end{example}

\begin{example}
Consider the Mori--Mukai model $MM_{2-17}$ on the algebraic torus
$(\mathbb{C}^*)^3$, with Laurent potential
\begin{align*}
W_{MM_{2-17}}={}&\frac{x^2}{yz}+\frac{x^2}{yz^2}+x+\frac{2x}{z}
+\frac{x}{yz}+y+z \\
&+\frac{2z}{x}+\frac{1}{x}+\frac{z}{x^2}.
\end{align*}
The logarithmic critical equations, after clearing denominators and saturating
by the torus monomial $xyz$, give the reduced ideal
\[
\left(
7x-88z^3+95z^2-75z+50,
7y+286z^3-350z^2+279z-192,
r(z)
\right),
\]
where
\[
r(z)=11z^4-16z^3+14z^2-10z+2.
\]
The polynomial $r$ is irreducible over $\mathbb{Q}$ and square-free
($\operatorname{disc}(r)=-1651888$). Hence the critical points are
\fitdisplay{P_{\eta}=\left(
\frac{88\eta^3-95\eta^2+75\eta-50}{7},
\frac{-286\eta^3+350\eta^2-279\eta+192}{7},
\eta
\right),
\qquad
r(\eta)=0}
The torus conditions are automatic, since $r$ is coprime to $z$ and to the
numerators of the displayed $x$- and $y$-coordinates.

The corresponding critical values are the four roots of
\[
w^4-90w^2-428w-391=0.
\]
Equivalently, on the quartic field generated by $\eta$,
\[
W_{MM_{2-17}}(P_{\eta})
=\frac{88\eta^3-7\eta^2+46\eta-23}{7}.
\]
The Hessian determinant is nonzero at every critical point. More precisely,
modulo $r(\eta)$ one obtains
\[
\det H_W(P_{\eta})
=-\frac{352\eta^3-1205\eta^2-458\eta+218}{14},
\]
and the numerator is coprime to $r$. Therefore all four singularities are
isolated Morse singularities of type $A_1$. Consequently,
\[
\overline{MF((\mathbb{C}^*)^3,W_{MM_{2-17}})}
\cong \bigoplus_{\eta:\,r(\eta)=0}\MorseCat{3}.
\]
\end{example}

\begin{example}
Consider the Mori--Mukai model $MM_{2-18}$ on the algebraic torus
$(\mathbb{C}^*)^3$, with Laurent potential
\begin{align*}
W_{MM_{2-18}}={}&x+yz^2+2yz+y+2z+\frac{1}{y}
+\frac{2}{x} \\
&+\frac{2}{xz}+\frac{2}{xyz}+\frac{1}{x^2yz^2}.
\end{align*}
Set
\[
A=z+1+\frac{1}{y}+\frac{1}{xyz}.
\]
Then the potential has the exact square completion
\[
W_{MM_{2-18}}=x+yA^2-2.
\]
The $x$-derivative gives
\[
\frac{\partial W_{MM_{2-18}}}{\partial x}
=1-\frac{2A}{x^2z}.
\]
Hence $A\neq0$ at every torus critical point. The remaining two equations
therefore reduce to
\[
xyz^2=1,
\qquad
A-\frac{2}{y}-\frac{2}{xyz}=0.
\]
Substituting the definition of $A$ into the second equation gives
\[
xyz^2+xyz-xz-1=0.
\]
Using $xyz^2=1$, this becomes $xz(y-1)=0$, and hence $y=1$ on the
torus critical locus. Thus $x=1/z^2$, and the $x$-equation becomes
\[
4z^4+4z^3-1=0.
\]

Equivalently, clearing denominators in the logarithmic critical equations and
saturating by the torus monomial $xyz$ gives the reduced ideal
\[
\left(
x-4z^2-4z,
y-1,
s(z)
\right),
\]
where
\[
s(z)=4z^4+4z^3-1.
\]
The polynomial $s$ is irreducible over $\mathbb{Q}$ and square-free
($\operatorname{disc}(s)=-23296$). Therefore the critical points are
\[
P_{\eta}=\left(4\eta^2+4\eta,1,\eta\right)
=\left(\frac{1}{\eta^2},1,\eta\right),
\qquad
s(\eta)=0.
\]
The torus conditions are automatic, since $s$ is coprime to $z$.

The corresponding critical values are the four roots of
\[
w^4-4w^3-32w^2-432w-464=0.
\]
Equivalently, on the quartic field generated by $\eta$,
\[
W_{MM_{2-18}}(P_{\eta})=2(4\eta^2+6\eta+1).
\]
The Hessian determinant is nonzero at every critical point. More precisely,
\[
\det H_W(P_{\eta})=32\eta^4(\eta+1)^2(4\eta+3),
\]
and $s$ is coprime to $z$, $z+1$, and $4z+3$. Hence all four singularities
are isolated Morse singularities of type $A_1$. Consequently,
\[
\overline{MF((\mathbb{C}^*)^3,W_{MM_{2-18}})}
\cong \bigoplus_{\eta:\,s(\eta)=0}\MorseCat{3}.
\]
\end{example}

\begin{example}
Consider the Mori--Mukai model $MM_{2-19}$ on the algebraic torus
$(\mathbb{C}^*)^3$, with Laurent potential
\begin{align*}
W_{MM_{2-19}}={}&\frac{x^2}{yz}+\frac{2x}{yz}+\frac{1}{yz}
+z+\frac{2yz}{x}+\frac{y^2z}{x^2} \\
&+x+y+\frac{y}{x}.
\end{align*}
This potential is the sum of two square blocks and one linear block:
\[
W_{MM_{2-19}}
=\frac{(x+1)^2}{yz}+z\frac{(x+y)^2}{x^2}
+x+y\frac{x+1}{x}.
\]

Clearing denominators in the logarithmic critical equations and saturating by
the torus monomial $xyz$ gives the decomposition
\[
I_{\mathrm{crit}}=I_{\Gamma}\cap I_{\mathrm{iso}},
\]
where
\[
I_{\Gamma}=(x+1,y-1)
\]
and
\[
I_{\mathrm{iso}}=\bigl(
x+yz+y-z,
y+2z^3-z^2-2z-3,
h(z)
\bigr),
\]
with
\[
h(z)=(z^2-z-1)(z^2+z+1)=z^4-z^2-2z-1.
\]
Thus the critical locus has a curve component
\[
\Gamma=\{(-1,1,z):z\in\mathbb{C}^*\},
\qquad
W_{MM_{2-19}}|_{\Gamma}=-1,
\]
and four isolated torus points.

The isolated component can also be obtained from the parametrization
\[
y=t^2,
\qquad
x=t^2+t^3.
\]
Before the torus saturation the remaining equation gives
\[
(t+2)(t^2+t-1)(t^2+t+1)=0.
\]
The root $t=-2$ is extraneous: it makes the formula
$z=xy(x+1)/(x^2-y^2)$ undefined. After saturation, only the quartic
\[
Q(t)=(t^2+t-1)(t^2+t+1)
\]
remains, and on this locus one may take $z=t(t+1)^2$.

Equivalently, using the coordinate $z$ on the isolated component, the four
isolated points are
\[
P_{\rho}=\left(
\rho^3-\rho^2-1,
-2\rho^3+\rho^2+2\rho+3,
\rho
\right),
\qquad
h(\rho)=0.
\]
The torus conditions are automatic, since $h$ is coprime to $z$, to
$z^3-z^2-1$, and to $-2z^3+z^2+2z+3$.

The critical values at the isolated points are the four roots of
\[
w^4+2w^3-73w^2-454w-407=0.
\]
On the two quadratic factors they are given more explicitly by
\[
W_{MM_{2-19}}(P_{\rho})=5\rho+2
\quad\text{if }\rho^2-\rho-1=0,
\]
and
\[
W_{MM_{2-19}}(P_{\rho})=3\rho-4
\quad\text{if }\rho^2+\rho+1=0.
\]
The Hessian determinant is nonzero at the isolated points. Indeed, modulo
$h(\rho)$,
\[
\det H_W(P_{\rho})=2(\rho+1)(19\rho^2+28\rho+10),
\]
and the right-hand side is coprime to $h$.

It remains to identify the transverse type along $\Gamma$. Write
\[
A_{\Gamma}=\mathbb C[t^{\pm1}],
\qquad
u=x+1,
\qquad
v=y-1,
\qquad
t=z.
\]
These coordinates globally split the formal neighborhood of the curve:
\[
\widehat X_{\Gamma}
\simeq
\operatorname{Spf}\!\left(A_{\Gamma}[[u,v]]\right).
\]
Modulo $(u,v)^3$, the shifted potential is
\[
\widehat W_{MM_{2-19}}+1
\equiv
\left(t-1+\frac1t\right)u^2+(2t-1)uv+tv^2.
\]
The corresponding normal Hessian is
\[
\begin{pmatrix}
2t-2+\dfrac{2}{t} & 2t-1\\
2t-1 & 2t
\end{pmatrix},
\qquad
\det=3.
\]
Thus the normal quadratic form is nondegenerate over all of $\Gamma$. The
relative formal Morse lemma therefore gives a global split Morse--Bott model,
and Proposition~\ref{prop:morse_bott_clifford} yields
\[
\overline{MF^{\Gamma}_{-1}((\mathbb C^*)^3,W_{MM_{2-19}})}
\simeq
\MBcat{\Gamma}.
\]
The four isolated points are Morse singularities, so the verified categorical
decomposition is
\[
\overline{MF((\mathbb{C}^*)^3,W_{MM_{2-19}})}
\cong \MBcat{\Gamma} \oplus \bigoplus_{\rho:\,h(\rho)=0}\MorseCat{3}.
\]
\end{example}


\begin{example}
Consider the Mori--Mukai model $MM_{2-20}$ on the algebraic torus
$(\mathbb{C}^*)^3$, with Laurent potential
\begin{align*}
W_{MM_{2-20}}={}&x+\frac{2}{x}+\frac{1}{y}+\frac{1}{z}
+\frac{x}{yz}+y+\frac{y}{x} \\
&+z+\frac{2z}{x}+\frac{z}{x^2}.
\end{align*}
Set
\[
A=1+\frac{1}{x}.
\]
Then the potential can be written in the compact form
\[
W_{MM_{2-20}}
=zA^2+yA+x+\frac{2}{x}+\frac{1}{y}+\frac{1}{z}+\frac{x}{yz}.
\]
The critical equations are zero-dimensional on the torus. More precisely,
clearing denominators in the logarithmic critical equations and saturating by
$xyz$ gives the reduced ideal
\[
\left(
157y-66x^5+34x^4+90x^3+174x^2-137x+94,
\right.
\]
\[
\left.
157z+310x^5-36x^4-594x^3-1274x^2+496x-118,
q(x)
\right),
\]
where
\[
q(x)=2x^6+2x^5-4x^4-12x^3-6x^2+2x-1.
\]
The polynomial $q$ is irreducible over $\mathbb{Q}$ and square-free
($\operatorname{disc}(q)=3112478528$). Thus the torus critical locus consists
of the six points
\[
\begin{aligned}
P_{\alpha}
=\biggl(&\alpha,\,
\frac{66\alpha^5-34\alpha^4-90\alpha^3-174\alpha^2+137\alpha-94}{157},\\
&\frac{-310\alpha^5+36\alpha^4+594\alpha^3+1274\alpha^2-496\alpha+118}{157}
\biggr),
\qquad q(\alpha)=0.
\end{aligned}
\]
The torus conditions are automatic: $q$ is coprime to
$\alpha$ and to the numerators of the displayed $y$- and $z$-coordinates.

The corresponding critical values are the six roots of
\[
w^6+2w^5-61w^4-524w^3-1881w^2-3162w-2059=0.
\]
Equivalently, in the sextic field generated by $\alpha$,
\[
W_{MM_{2-20}}(P_{\alpha})
=-\frac{204\alpha^5+266\alpha^4-478\alpha^3-1594\alpha^2-1118\alpha+523}{157}.
\]
The logarithmic Hessian is nondegenerate at every critical point. Modulo
$q(\alpha)$, its determinant is
\[
\det(\theta_i\theta_jW)(P_{\alpha})
=\frac{2(592\alpha^5+618\alpha^4-950\alpha^3-2988\alpha^2-874\alpha+1293)}{157},
\]
and the numerator is coprime to $q$. Since the logarithmic coordinates are
invertible on the torus, the ordinary Hessian is nondegenerate as well. Hence
all six critical points are isolated Morse singularities of type $A_1$.
Consequently,
\[
\overline{MF((\mathbb{C}^*)^3,W_{MM_{2-20}})}
\cong \bigoplus_{\alpha:\,q(\alpha)=0}\MorseCat{3}.
\]
\end{example}


\begin{example}
Consider the Mori--Mukai model $MM_{2-21}$ on the algebraic torus
$(\mathbb{C}^*)^3$, with Laurent potential
\begin{align*}
W_{MM_{2-21}}={}&x+y+y^2z+2yz+z+\frac{x}{yz}
+\frac{2}{yz}+\frac{1}{xyz}.
\end{align*}
Equivalently,
\[
W_{MM_{2-21}}=x+y+z(y+1)^2+\frac{(x+1)^2}{xyz}.
\]
The $x$-derivative gives
\[
\frac{\partial W_{MM_{2-21}}}{\partial x}
=1+\frac{1}{yz}-\frac{1}{x^2yz},
\]
so that $yz=(1-x^2)/x^2$ on the torus. Combining the $y$- and
$z$-equations gives
\[
z=\frac{y}{1-y^2},
\qquad
x^2+y^2=1,
\qquad
z=\frac{y}{x^2}.
\]
Substituting these relations into the remaining equation gives the finite
intersection
\[
x^3(x+1)^2=y^3(y+1)^2.
\]

Equivalently, clearing denominators in the logarithmic critical equations and
saturating by the torus monomial $xyz$ gives the reduced ideal
\[
\left(
2x^2-3x+3y-1,
-30x^3+82x^2-81x+48z-41,
\right.
\]
\[
\left.
(2x^2-1)(x^2-3x+4)
\right).
\]
Thus the torus critical locus consists of four points. The factor
$2x^2-1$ gives
\[
P_{+}=\left(\frac{\sqrt{2}}{2},\frac{\sqrt{2}}{2},\sqrt{2}\right),
\qquad
P_{-}=\left(-\frac{\sqrt{2}}{2},-\frac{\sqrt{2}}{2},-\sqrt{2}\right).
\]
The factor $x^2-3x+4$ gives the conjugate pair
\[
Q_{+}=\left(
\frac{3+i\sqrt{7}}{2},
\frac{3-i\sqrt{7}}{2},
\frac{-9-5i\sqrt{7}}{32}
\right),
\]
\[
Q_{-}=\left(
\frac{3-i\sqrt{7}}{2},
\frac{3+i\sqrt{7}}{2},
\frac{-9+5i\sqrt{7}}{32}
\right).
\]
These are precisely the two branches of the relation $x^2+y^2=1$: the
symmetric branch $x=y$ and the asymmetric branch $x+y=3$, $xy=4$.

The corresponding critical values are
\[
W(P_{+})=4+4\sqrt{2},
\qquad
W(P_{-})=4-4\sqrt{2},
\qquad
W(Q_{+})=W(Q_{-})=-5.
\]
Equivalently, the critical-value polynomial with multiplicities is
\[
(w+5)^2(w^2-8w-16)=0.
\]
The logarithmic Hessian is nondegenerate at every critical point. If
\[
h(x)=(2x^2-1)(x^2-3x+4),
\]
then, modulo $h(x)$,
\[
\det(\theta_i\theta_jW)
=\frac{2(32x^3-70x^2+50x+83)}{3},
\]
and the numerator is coprime to $h$. Hence all four critical points are
isolated Morse singularities of type $A_1$. Consequently,
\[
\overline{MF((\mathbb{C}^*)^3,W_{MM_{2-21}})}
\cong \bigoplus_{i=1}^{4}\MorseCat{3}.
\]
\end{example}

\begin{example}
Consider the Mori--Mukai model $MM_{2-22}$ on the algebraic torus
$(\mathbb{C}^*)^3$, with Laurent potential
\begin{align*}
W_{MM_{2-22}}={}&x+\frac{x}{y}+y+z+\frac{1}{y}
+\frac{1}{yz}+\frac{z}{x}+\frac{2}{x}+\frac{1}{xz}.
\end{align*}
Equivalently,
\[
W_{MM_{2-22}}
=x\left(1+\frac{1}{y}\right)
+\frac{(z+1)^2}{xz}+y+z+\frac{1}{y}+\frac{1}{yz}.
\]
The critical equations may be written as
\[
x^2z(y+1)=y(z+1)^2,
\qquad
y^2=x+1+\frac{1}{z},
\qquad
yz^2(x+1)=x+y.
\]
After clearing denominators in the logarithmic critical equations and
saturating by the torus monomial $xyz$, the reduced critical ideal is
\[
\left(
2y-x^4-x^3+x^2+3x,
2z-x^2+1,
p(x)
\right),
\]
where
\[
p(x)=x^5+x^4-2x^3-4x^2+x-1.
\]
The polynomial $p$ is irreducible over $\mathbb{Q}$ and square-free
($\operatorname{disc}(p)=220672$). Hence the torus critical locus consists of
the five points
\[
P_{\lambda}=\left(
\lambda,
\frac{\lambda^4+\lambda^3-\lambda^2-3\lambda}{2},
\frac{\lambda^2-1}{2}
\right),
\qquad
p(\lambda)=0.
\]
The torus conditions are automatic, since $p$ is coprime to $x$, to
$x^2-1$, and to $x^4+x^3-x^2-3x$. On this locus one also has
$yz=x$, so this is the saturated version of the coupling variable
$u=yz$.

The corresponding critical values are the five roots of
\[
w^5+2w^4-48w^3-348w^2-640w-400=0.
\]
Equivalently, in the quintic field generated by $\lambda$,
\[
W_{MM_{2-22}}(P_{\lambda})
=\lambda^4+\lambda^3-\lambda-1.
\]
The logarithmic Hessian is nondegenerate at every critical point. More
precisely, modulo $p(\lambda)$ one obtains
\[
\det(\theta_i\theta_jW)(P_{\lambda})
=\lambda^4+3\lambda^3+9\lambda^2+\lambda+2,
\]
and this polynomial is coprime to $p$. Since the logarithmic coordinates are
invertible on the torus, all five critical points are isolated Morse
singularities of type $A_1$. Consequently,
\[
\overline{MF((\mathbb{C}^*)^3,W_{MM_{2-22}})}
\cong \bigoplus_{\lambda:\,p(\lambda)=0}\MorseCat{3}.
\]
\end{example}


\begin{example}
Consider the Mori--Mukai model $MM_{2-23}$ on the algebraic torus
$(\mathbb{C}^*)^3$, with Laurent potential
\[
W_{MM_{2-23}}=x^2y+2xy+x+y+z+\frac{2}{xy}
+\frac{1}{x^2y^2z}.
\]
The $z$-derivative factors as
\[
\frac{\partial W_{MM_{2-23}}}{\partial z}
=1-\frac{1}{x^2y^2z^2},
\]
so any critical point must lie on one of the two branches $xyz=1$ or
$xyz=-1$. The branch $xyz=-1$ is empty in the torus critical locus: after
substituting $z=-1/(xy)$, the $y$-equation gives $(x+1)^2=0$, while the
$x$-equation becomes incompatible. Hence all torus critical points lie on the
branch
\[
z=\frac{1}{xy}.
\]
On this branch the remaining equations reduce to
\[
y^2=\frac{4}{x(x+1)^2},
\qquad
2y(x+1)+1-\frac{4}{x^2y}=0.
\]
Equivalently,
\[
y=\frac{4(1-x)}{x^2(x+1)},
\qquad
x^3-4x^2+8x-4=0.
\]

In reproducible ideal-theoretic form, clearing denominators in the logarithmic
critical equations and saturating by the torus monomial $xyz$ gives the reduced
ideal
\[
\left(
-8x^2+23x+17y-36,
-2x^2+7x+4z-8,
q(x)
\right),
\]
where
\[
q(x)=x^3-4x^2+8x-4.
\]
The polynomial $q$ is irreducible over $\mathbb{Q}$ and square-free
($\operatorname{disc}(q)=-176$). Thus the torus critical locus consists of the
three points
\[
P_{\alpha}=\left(
\alpha,
\frac{4(1-\alpha)}{\alpha^2(\alpha+1)},
\frac{\alpha(\alpha+1)}{4(1-\alpha)}
\right),
\qquad
q(\alpha)=0.
\]
Equivalently, modulo
$q(\alpha)$, the coordinates are
\[
P_{\alpha}=\left(
\alpha,
\frac{8\alpha^2-23\alpha+36}{17},
\frac{2\alpha^2-7\alpha+8}{4}
\right).
\]
The torus conditions are automatic, since $q$ is coprime to $x$, $x+1$, and
$x-1$.

The corresponding critical values are the three roots of
\[
w^3+4w^2-72w-364=0.
\]
Equivalently, in the cubic field generated by $\alpha$,
\[
W_{MM_{2-23}}(P_{\alpha})=4\alpha^2-13\alpha+16.
\]
The Hessian determinant is nonzero at every critical point. More precisely,
modulo $q(\alpha)$,
\[
\det H_W(P_{\alpha})=2(9\alpha^2-24\alpha+35),
\]
and this polynomial is coprime to $q$. Hence all three critical points are
isolated Morse singularities of type $A_1$. Consequently,
\[
\overline{MF((\mathbb{C}^*)^3,W_{MM_{2-23}})}
\cong \bigoplus_{\alpha:\,q(\alpha)=0}\MorseCat{3}.
\]
\end{example}


\begin{example}
Consider the Mori--Mukai model $MM_{2-24}$ on the algebraic torus
$(\mathbb{C}^*)^3$, with Laurent potential
\begin{align*}
W_{MM_{2-24}}={}&y+\frac{2y}{xz}+\frac{y}{x^2z^2}
+z+\frac{z}{y}+x+\frac{1}{x}+\frac{2}{xz}+\frac{1}{y}.
\end{align*}
Equivalently,
\[
W_{MM_{2-24}}
=y\left(1+\frac{1}{xz}\right)^2
+z\left(\frac{y+1}{y}\right)+x+\frac{1}{x}+\frac{2}{xz}+\frac{1}{y}.
\]
Clearing denominators in the logarithmic critical equations and saturating by
$xyz$ gives the reduced ideal
\[
\left(
 x^2-2xz-1,
 y-1,
 z^3-1
\right).
\]
Thus $y=1$, $z$ is a cubic root of unity, and for each such $z$ the remaining
coordinate satisfies
\[
x^2-2zx-1=0.
\]
Since $z^3-1$ is square-free and coprime to $z^2+1$, the quadratic has two
separate roots for every cubic root of unity. Hence the torus critical locus
consists of the six points
\[
P_{\zeta,\varepsilon}=\left(
\zeta+\varepsilon\sqrt{\zeta^2+1},1,\zeta
\right),
\qquad
\zeta^3=1,
\quad
\varepsilon\in\{+1,-1\}.
\]
In particular, for the real branch $\zeta=1$ one obtains
\[
(1+\sqrt{2},1,1),
\qquad
(1-\sqrt{2},1,1).
\]
If $\omega=-\frac12+\frac{\sqrt3}{2}i$, the four remaining points are obtained
from the same formula with $\zeta=\omega$ and $\zeta=\omega^2$.

The critical values are the six roots of
\[
(w^2-6w-23)(w^4+6w^3+35w^2+78w+61)=0.
\]
The two values on the real branch are $3+4\sqrt2$ and $3-4\sqrt2$. The
logarithmic Hessian is nondegenerate at every critical point. More precisely,
modulo the critical ideal, its numerator reduces to
\[
6(33xz^2+18xz+48x+20z^2+16z+5),
\]
and adding this polynomial to the critical ideal gives the unit ideal. Since
$x$, $y$, and $z$ are invertible on the torus, the ordinary Hessian is also
nondegenerate. Therefore all six critical points are isolated Morse
singularities of type $A_1$. Consequently,
\[
\overline{MF((\mathbb{C}^*)^3,W_{MM_{2-24}})}
\cong \bigoplus_{\zeta^3=1}\left(\MorseCat{3}\oplus \MorseCat{3}\right).
\]
\end{example}

\begin{example}
Consider the Mori--Mukai model $MM_{2-25}$ on the algebraic torus
$(\mathbb{C}^*)^3$, with Laurent potential
\[
W_{MM_{2-25}}=x+\frac{x}{y}+y+z+\frac{1}{yz}
+\frac{yz}{x}+\frac{2}{x}+\frac{yz}{x^2}.
\]
The logarithmic critical equations have a useful transverse symmetry. Indeed,
\[
y\frac{\partial W_{MM_{2-25}}}{\partial y}
-z\frac{\partial W_{MM_{2-25}}}{\partial z}
=-\frac{x}{y}+y-z,
\]
so every torus critical point satisfies
\[
x=y(y-z).
\]
Writing $v=y-z$, so that $x=yv$ and $y=v+z$, the remaining equations reduce,
after clearing denominators and removing the nonzero torus factor $v+z$, to
\[
(v-1)(v^3+v^2z+2v^2+vz+2v+2)=0,
\]
and
\[
v^2z^2+vz^2-v+z=0.
\]
The branch $v=1$ gives $(z+1)(2z-1)=0$; the root $z=-1$ is excluded by the
torus condition $y=z+1\neq 0$, hence it gives the rational critical point
\[
\left(\frac{3}{2},\frac{3}{2},\frac{1}{2}\right).
\]
On the complementary branch, saturating by $vz(v+z)$ gives
\[
v^2+1=0,
\qquad
v=2z+1.
\]

Equivalently, clearing denominators in the original logarithmic critical
equations and saturating by $xyz$ gives the reduced ideal
\[
\left(
5x-16z^2-11z+2,
5y+4z^2-11z-3,
q(z)
\right),
\]
where
\[
q(z)=(2z-1)(2z^2+2z+1).
\]
Thus the torus critical locus consists of the three points
\[
P_z=\left(
\frac{16z^2+11z-2}{5},
\frac{-4z^2+11z+3}{5},
z
\right),
\qquad
q(z)=0.
\]
The rational root $z=1/2$ gives the point displayed above, while the two
remaining points correspond to $z=(-1\pm i)/2$.

The critical values are the three roots of
\[
(w-8)(w^2+8w+32)=0.
\]
Equivalently, on the critical locus one has
\[
W_{MM_{2-25}}(P_z)=\frac{8}{5}(2z^2+7z+1),
\]
so the rational point lies on the fiber $W=8$ and the two conjugate values are
$-4\pm 4i$. The logarithmic Hessian is nondegenerate at every critical point:
modulo $q(z)$,
\[
\det(\theta_i\theta_jW)(P_z)=\frac{4}{5}(-4z^2+26z+33),
\]
and this polynomial is coprime to $q$. Since logarithmic coordinates are
invertible on the torus, all three critical points are isolated Morse
singularities of type $A_1$. Consequently,
\[
\overline{MF((\mathbb{C}^*)^3,W_{MM_{2-25}})}
\cong \bigoplus_{z:\,q(z)=0}\MorseCat{3}.
\]
\end{example}

\begin{example}
Consider the Mori--Mukai model $MM_{2-26}$ on the algebraic torus
$(\mathbb{C}^*)^3$, with Laurent potential
\[
W_{MM_{2-26}}=x+\frac{z}{x}+\frac{2}{x}+\frac{1}{xz}
+y+\frac{y}{z}+z+\frac{1}{y}.
\]
Equivalently,
\[
W_{MM_{2-26}}
=x+\frac{(z+1)^2}{xz}+\frac{y(z+1)}{z}+z+\frac{1}{y}.
\]
The $x$- and $y$-derivatives give the two simple relations
\[
x^2=\frac{(z+1)^2}{z},
\qquad
y^2=\frac{z}{z+1}.
\]
Thus, if $u=xy$, then every torus critical point satisfies
\[
z=u^2-1.
\]
The remaining $z$-equation is
\[
z^2(x+1)-xy-1=0.
\]
Since $u=-1$ would force $z=0$, which is forbidden on the torus, this equation
gives
\[
x+1=\frac{u+1}{(u^2-1)^2}
=\frac{1}{(u^2-1)(u-1)},
\]
and hence
\[
x=\frac{-u^3+u^2+u}{(u^2-1)(u-1)}.
\]
Combining this expression with $x^2=u^4/(u^2-1)$ gives the single equation
\[
\frac{(-u^2+u+1)^2}{(u^2-1)^2(u-1)^2}
=\frac{u^2}{u^2-1},
\]
or equivalently
\[
p(u):=u^6-2u^5-u^4+4u^3-2u-1=0.
\]

In ideal-theoretic form, clearing denominators in the logarithmic critical
equations, saturating by $xyz$, and adjoining the relation $u-xy=0$ eliminates
$x,y,z$ to the principal equation $p(u)=0$. The polynomial $p$ is irreducible
over $\mathbb{Q}$ and square-free ($\operatorname{disc}(p)=37568$). It is also
coprime to $u$, $u^2-1$, and $u^2-u-1$, so the torus conditions and all
denominators below are automatic. Therefore the torus critical locus consists
of the six points
\[
P_u=\left(
\frac{-u^3+u^2+u}{(u^2-1)(u-1)},
\frac{u}{x_u},
u^2-1
\right),
\qquad
x_u=\frac{-u^3+u^2+u}{(u^2-1)(u-1)},
\quad
p(u)=0.
\]

The critical values are the six roots of
\[
w^6+4w^5-46w^4-302w^3-431w^2-250w-43=0.
\]
Equivalently, in the quotient $\mathbb{Q}[u]/(p(u))$ one has
\[
W_{MM_{2-26}}(P_u)=(u+1)^2(2u^3-4u^2+2u+1).
\]
The logarithmic Hessian is nondegenerate at every critical point. More
precisely, modulo $p(u)$,
\[
\det(\theta_i\theta_jW)(P_u)
=2(3u^5-2u^4-3u^3+5u^2+3u+1),
\]
and this polynomial is coprime to $p$. Since logarithmic coordinates are
invertible on the torus, all six critical points are isolated Morse
singularities of type $A_1$. Consequently,
\[
\overline{MF((\mathbb{C}^*)^3,W_{MM_{2-26}})}
\cong \bigoplus_{u:\,p(u)=0}\MorseCat{3}.
\]
\end{example}

\begin{example}
Consider the Mori--Mukai model $MM_{2-27}$ on the algebraic torus
$(\mathbb{C}^*)^3$, with Laurent potential
\[
W_{MM_{2-27}}=x+\frac{x}{z}+y+z+\frac{1}{yz}
+\frac{1}{x}+\frac{1}{xy}.
\]
The logarithmic critical equations are
\[
x^2y(z+1)-z(y+1)=0,
\qquad
xy^2z-x-z=0,
\qquad
yz^2-xy-1=0.
\]
Equivalently,
\[
x^2=\frac{z(y+1)}{y(z+1)},
\qquad
y^2=\frac{x+z}{xz},
\qquad
z^2=\frac{xy+1}{y}.
\]
Clearing denominators and saturating by the torus monomial $xyz$ gives the
reduced ideal
\[
\left(
x-z^5+z^2+z,
y-z,
q(z)
\right),
\]
where
\[
q(z)=(z^3-z-1)(z^3+z-1).
\]
Thus every critical point satisfies $y=z$. On the first branch
$z^3-z-1=0$, the relation $x=z^5-z^2-z$ reduces to $x=1$; on the second
branch $z^3+z-1=0$, it reduces to $x=-1$.

Equivalently, introducing the coupling variable $u=xy$, one obtains the
elimination equation
\[
p(u):=u^6-(u+1)^2=(u^3-u-1)(u^3+u+1)=0.
\]
The two cubic factors are irreducible and coprime over $\mathbb{Q}$, and $p$
is square-free ($\operatorname{disc}(p)=45632$). Since $p$ is coprime to $u$
and $u+1$, the torus conditions and the denominators in the following
parametrization are automatic. The six critical points are
\[
P_u=\left(
\frac{u^3}{u+1},
\frac{u+1}{u^2},
\frac{u^4}{u+1}
\right),
\qquad
p(u)=0.
\]
For the branch $u^3-u-1=0$ this becomes
\[
P_u=(1,u,u),
\]
whereas for the branch $u^3+u+1=0$ it becomes
\[
P_u=(-1,-u,-u).
\]

The critical values are the six roots of
\[
(w^3-5w^2-10w-11)(w^3+7w^2+34w+37)=0.
\]
Equivalently, in $\mathbb{Q}[u]/(p(u))$ one has
\[
W_{MM_{2-27}}(P_u)=u^4+2u^3-1.
\]
The logarithmic Hessian is nondegenerate at every critical point. More
precisely, modulo $p(u)$,
\[
\det(\theta_i\theta_jW)(P_u)=2(3u^5-u-1),
\]
and this polynomial is coprime to $p$. Since logarithmic coordinates are
invertible on the torus, all six critical points are isolated Morse
singularities of type $A_1$. Consequently,
\[
\overline{MF((\mathbb{C}^*)^3,W_{MM_{2-27}})}
\cong \bigoplus_{u:\,p(u)=0}\MorseCat{3}.
\]
\end{example}

\begin{example}
Consider the Mori--Mukai model $MM_{2-28}$ on the algebraic torus
$(\mathbb{C}^*)^3$, with Laurent potential
\[
W_{MM_{2-28}}=xyz^2+xyz+x+y+z+\frac{1}{yz}+\frac{1}{xz}.
\]
The logarithmic critical equations are
\[
x^2yz^3+x^2yz^2+x^2z-1=0,
\]
\[
xy^2z^3+xy^2z^2+y^2z-1=0,
\]
and
\[
2x^2y^2z^3+x^2y^2z^2+xyz^2-x-y=0.
\]
After clearing denominators and saturating by the torus monomial $xyz$, the
reduced critical locus decomposes into a smooth curve and four isolated
points. The curve component is
\[
\Gamma=\{x+y-z-1=0,\ xyz+1=0\}\subset(\mathbb{C}^*)^3.
\]
Equivalently, writing $x=z+1-y$, it is cut out in the $(y,z)$-torus by
\[
y^2z-yz^2-yz-1=0.
\]
This affine curve is smooth: the ideal generated by this equation and its two
partial derivatives is the unit ideal after saturation by $yz$. Along
$\Gamma$, the potential is constant,
\[
W_{MM_{2-28}}|_{\Gamma}=-1.
\]
In fact, the relative formal model is globally split. Set
\[
F=x+y-z-1,
\qquad
G=xyz+1.
\]
A direct calculation gives the exact identity
\[
W_{MM_{2-28}}+1
=
\frac{G\bigl(F+(z+1)G\bigr)}{xyz}.
\]
Thus the global functions
\[
u=G,
\qquad
v=\frac{F+(z+1)G}{xyz}
\]
generate the same ideal as $(F,G)$ and satisfy $W_{MM_{2-28}}+1=uv$.
Writing $A_{\Gamma}=\mathcal O(\Gamma)$, the smoothness and affineness of
$\Gamma$ imply that the quotient map from the completed coordinate ring to
$A_{\Gamma}$ admits a section. Since $(u,v)$ is a regular sequence and gives
a basis of the conormal bundle, this yields
\[
\widehat X_{\Gamma}
\simeq
\operatorname{Spf}\!\left(A_{\Gamma}[[u,v]]\right),
\qquad
\widehat W_{MM_{2-28}}+1=uv.
\]
Hence $\Gamma$ is a global split Morse--Bott critical curve.

The isolated component is
\[
\left(
x+2z^3-z^2-4,
y+2z^3-z^2-4,
r(z)
\right),
\qquad
r(z)=z^4-2z-1.
\]
The polynomial $r$ is irreducible over $\mathbb{Q}$ and square-free
($\operatorname{disc}(r)=-688$). It is coprime to $z$ and to
$-2z^3+z^2+4$, so the torus conditions are automatic. Thus the isolated
critical points are
\[
P_{\lambda}=\left(
-2\lambda^3+\lambda^2+4,
-2\lambda^3+\lambda^2+4,
\lambda
\right),
\qquad
r(\lambda)=0.
\]
They do not meet $\Gamma$, since $r$ is coprime to
$3z^3-2z^2+z-8$, the reduction of $xyz+1$ on this branch.

The critical values of the isolated points are the four roots of
\[
w^4-4w^3+6w^2-192w+473=0.
\]
Equivalently, in $\mathbb{Q}[\lambda]/(r(\lambda))$ one has
\[
W_{MM_{2-28}}(P_{\lambda})
=-(\lambda-2)(2\lambda^2+3\lambda+2).
\]
The logarithmic Hessian is nondegenerate at every isolated critical point.
More precisely, modulo $r(\lambda)$,
\[
\det(\theta_i\theta_jW)(P_{\lambda})
=2(38\lambda^3-17\lambda^2+8\lambda-73),
\]
and this polynomial is coprime to $r$. Hence the four isolated critical points
are Morse singularities of type $A_1$.

Proposition~\ref{prop:morse_bott_clifford} identifies the contribution of the
curve with $\MBcat{\Gamma}$. Combining it with the four isolated Morse
contributions gives
\[
\overline{MF((\mathbb{C}^*)^3,W_{MM_{2-28}})}
\cong \MBcat{\Gamma}\oplus\bigoplus_{\lambda:\,r(\lambda)=0}\MorseCat{3}.
\]
\end{example}

\begin{example}
Consider the Mori--Mukai model $MM_{2-29}$ on the algebraic torus
$(\mathbb{C}^*)^3$, with Laurent potential
\[
W_{MM_{2-29}}=x+\frac{x}{y}+y+z+\frac{2}{x}+\frac{1}{x^2z}.
\]
The ordinary critical equations are
\[
1-\frac{1}{x^2z^2}=0,
\qquad
1-\frac{x}{y^2}=0,
\qquad
1+\frac{1}{y}-\frac{2}{x^2}-\frac{2}{x^3z}=0.
\]
Thus every critical point satisfies
\[
x=y^2,
\qquad
xz=\pm1.
\]
On the branch $xz=-1$, the $x$-equation reduces to $1+1/y=0$; hence
\[
(x,y,z)=(1,-1,-1).
\]
On the branch $xz=1$, the same equation becomes
\[
1+\frac{1}{y}-\frac{4}{x^2}=0.
\]
Using $x=y^2$, this is equivalent to
\[
q(y):=y^4+y^3-4=0,
\]
and then $z=1/y^2$.

Equivalently, clearing denominators in the logarithmic critical equations and
saturating by $xyz$ gives the reduced ideal
\[
\begin{aligned}
\bigl(&5x+64z^4+80z^3-32z^2-44z-1,\\
&5y-96z^4-80z^3+48z^2+26z-1,\\
&(z+1)(16z^4-8z^2-z+1)\bigr).
\end{aligned}
\]
This is the disjoint union of the rational point above and the four points
\[
P_{\alpha}=\left(\alpha^2,\alpha,\frac{1}{\alpha^2}\right),
\qquad
q(\alpha)=0.
\]
The polynomial $q$ is irreducible over $\mathbb{Q}$ and square-free
($\operatorname{disc}(q)=-16816$), and it is coprime to $y$. Hence these four
points are all torus points and are pairwise distinct.

The rational point has critical value
\[
W_{MM_{2-29}}(1,-1,-1)=-1.
\]
The remaining critical values are the four roots of
\[
w^4+w^3-32w^2-144w+148=0.
\]
Equivalently, in $\mathbb{Q}[\alpha]/(q(\alpha))$ one has
\[
W_{MM_{2-29}}(P_{\alpha})=\alpha(2\alpha+3).
\]
The logarithmic Hessian is nondegenerate at every critical point. At the
rational point one obtains
\[
\det(\theta_i\theta_jW)(1,-1,-1)=2,
\]
whereas on the quartic branch, modulo $q(\alpha)$,
\[
\det(\theta_i\theta_jW)(P_{\alpha})=2(4\alpha+3),
\]
and $4\alpha+3$ is coprime to $q$. Since logarithmic coordinates are invertible
on the torus, all five critical points are isolated Morse singularities of
type $A_1$. Consequently,
\[
\overline{MF((\mathbb{C}^*)^3,W_{MM_{2-29}})}
\cong \MorseCat{3}\oplus\bigoplus_{\alpha:\,q(\alpha)=0}\MorseCat{3}.
\]
\end{example}

\begin{example}
Consider the Mori--Mukai model $MM_{2-30}$ on the algebraic torus
$(\mathbb{C}^*)^3$, with Laurent potential
\[
W_{MM_{2-30}}=x+y+z+\frac{z}{xy}+\frac{2}{xy}+\frac{1}{xyz}.
\]
The logarithmic critical equations are
\[
x^2yz-(z+1)^2=0,
\qquad
xy^2z-(z+1)^2=0,
\qquad
xyz^2+z^2-1=0.
\]
The first two equations imply $x=y$ on the torus. Substituting $y=x$ gives
\[
x^3z=(z+1)^2,
\qquad
x^2=\frac{1-z^2}{z^2}.
\]
Dividing the first relation by the second gives
\[
x=\frac{z(z+1)}{1-z}.
\]
Substitution back into $x^2=(1-z^2)/z^2$ yields
\[
z^4(z+1)^2=(1-z)^3(1+z).
\]
Since $z=-1$ would force $x=0$, which is forbidden in the torus, this is
equivalent to
\[
p(z):=z^5+z^4+z^3-3z^2+3z-1=0.
\]

In ideal-theoretic form, clearing denominators in the logarithmic critical
equations and saturating by $xyz$ gives the reduced ideal
\[
\left(
x-z^4-2z^3-3z^2+z-1,
y-z^4-2z^3-3z^2+z-1,
p(z)
\right).
\]
The polynomial $p$ is irreducible over $\mathbb{Q}$ and square-free
($\operatorname{disc}(p)=17456$). It is coprime to $z$, $z-1$, and $z+1$, so
the torus conditions and the displayed parametrization have no exceptional
roots. Thus the critical locus consists of the five points
\[
P_{\lambda}=\left(
\frac{\lambda(\lambda+1)}{1-\lambda},
\frac{\lambda(\lambda+1)}{1-\lambda},
\lambda
\right),
\qquad
p(\lambda)=0.
\]

The corresponding critical values are the five roots of
\[
w^5+w^4+4w^3-120w^2-384w-688=0.
\]
Equivalently, in $\mathbb{Q}[\lambda]/(p(\lambda))$ one has
\[
W_{MM_{2-30}}(P_{\lambda})
=3\lambda^4+6\lambda^3+9\lambda^2-2\lambda+3.
\]
The logarithmic Hessian is nondegenerate at every critical point. More
precisely, modulo $p(\lambda)$,
\[
\det(\theta_i\theta_jW)(P_{\lambda})
=2(2\lambda^4+2\lambda^3+2\lambda^2-7\lambda+8),
\]
and this polynomial is coprime to $p$. Since logarithmic coordinates are
invertible on the torus, all five critical points are isolated Morse
singularities of type $A_1$. Consequently,
\[
\overline{MF((\mathbb{C}^*)^3,W_{MM_{2-30}})}
\cong \bigoplus_{\lambda:\,p(\lambda)=0}\MorseCat{3}.
\]
\end{example}

\begin{example}
Consider the Mori--Mukai model $MM_{2-31}$ on the algebraic torus
$(\mathbb{C}^*)^3$, with Laurent potential
\[
W_{MM_{2-31}}=x+\frac{x}{y}+y+z+\frac{1}{yz}+\frac{1}{x}.
\]
The logarithmic critical equations are
\[
x^2y+x^2-y=0,
\qquad
-xz+y^2z-1=0,
\qquad
yz^2-1=0.
\]
The last equation gives
\[
y=\frac{1}{z^2}.
\]
Substituting this into the middle equation gives the cascade relation
\[
x=y^2-\frac{1}{z}=\frac{1-z^3}{z^4}.
\]
The remaining equation $x^2=y/(y+1)$ becomes
\[
\left(\frac{1-z^3}{z^4}\right)^2=\frac{1}{z^2+1},
\]
or equivalently
\[
p(z):=z^6-2z^5-2z^3+z^2+1=0.
\]

In ideal-theoretic form, clearing denominators in the logarithmic critical
equations and saturating by $xyz$ gives the reduced ideal
\[
\left(
x+z^5-3z^4+2z^3-z^2+z-1,
y+z^4-2z^3-2z+1,
p(z)
\right).
\]
The polynomial $p$ is irreducible over $\mathbb{Q}$ and square-free
($\operatorname{disc}(p)=85952$). It is coprime to $z$ and to $1-z^3$, so all
displayed points lie in the torus. Hence the critical locus consists of the
six points
\[
P_{\lambda}=\left(
\frac{1-\lambda^3}{\lambda^4},
\frac{1}{\lambda^2},
\lambda
\right),
\qquad
p(\lambda)=0.
\]

The corresponding critical values are the six roots of
\[
w^6+2w^5-11w^4-142w^3-186w^2-436w-91=0.
\]
Equivalently, in $\mathbb{Q}[\lambda]/(p(\lambda))$ one has
\[
W_{MM_{2-31}}(P_{\lambda})
=-(\lambda+1)(\lambda^2+1)(2\lambda^2-5\lambda+1).
\]
The logarithmic Hessian is nondegenerate at every critical point. More
precisely, modulo $p(\lambda)$,
\[
\det(\theta_i\theta_jW)(P_{\lambda})
=-2\lambda(\lambda^4+\lambda^3-7\lambda^2-2),
\]
and this polynomial is coprime to $p$. Since logarithmic coordinates are
invertible on the torus, all six critical points are isolated Morse
singularities of type $A_1$. Consequently,
\[
\overline{MF((\mathbb{C}^*)^3,W_{MM_{2-31}})}
\cong \bigoplus_{\lambda:\,p(\lambda)=0}\MorseCat{3}.
\]
\end{example}

\begin{example}
Consider the Mori--Mukai model $MM_{2-32}$ on the algebraic torus
$(\mathbb{C}^*)^3$, with Laurent potential
\[
W_{MM_{2-32}}=\frac{xy}{z}+\frac{x}{z}+\frac{y}{z}+\frac{2}{z}
+\frac{1}{xz}+\frac{1}{yz}+\frac{1}{xyz}+z.
\]
The logarithmic derivatives are
\[
x\frac{\partial W}{\partial x}
=\frac{(y+1)(x^2y-1)}{xyz},
\qquad
y\frac{\partial W}{\partial y}
=\frac{(x+1)(xy^2-1)}{xyz},
\]
and
\[
z\frac{\partial W}{\partial z}
=\frac{xy z^2-x^2y^2-x^2y-xy^2-2xy-x-y-1}{xyz}.
\]
After clearing denominators and saturating by the torus monomial $xyz$, the
critical ideal is the radical ideal
\[
\left(
x-y,
3y^2+3y-z^2+2,
(y-1)(z^2+1),
(z^2-8)(z^2+1)
\right).
\]
Equivalently, it decomposes as
\[
(x-1,y-1,z^2-8)
\cap
(x-y,y^2+y+1,z^2+1).
\]
Thus the factor $z^2+1$ is not removed by the torus saturation: it gives
genuine torus critical points. The critical locus consists of the two points
\[
P_{\epsilon}=(1,1,\epsilon),
\qquad \epsilon^2=8,
\]
and the four points
\[
Q_{\omega,\delta}=(\omega,\omega,\delta),
\qquad \omega^2+\omega+1=0,\quad \delta^2+1=0.
\]

The corresponding critical values are
\[
W_{MM_{2-32}}(P_{\epsilon})=\frac{16}{\epsilon}=2\epsilon,
\qquad
W_{MM_{2-32}}(Q_{\omega,\delta})=-\frac{2}{\delta}=2\delta.
\]
Hence the distinct critical values are the roots of
\[
(w^2-32)(w^2+4)=0,
\]
with each value in the branch $w^2+4=0$ occurring at two critical points. The
logarithmic Hessian is nondegenerate at every critical point. More precisely,
\[
\det(\theta_i\theta_jW)(P_{\epsilon})=\frac{24}{\epsilon}=3\epsilon,
\qquad
\det(\theta_i\theta_jW)(Q_{\omega,\delta})=\frac{6}{\delta}=-6\delta.
\]
Since logarithmic coordinates are invertible on the torus, all six critical
points are isolated Morse singularities of type $A_1$. Consequently,
\[
\overline{MF((\mathbb{C}^*)^3,W_{MM_{2-32}})}
\cong \bigoplus_{\epsilon^2=8}\MorseCat{3} \oplus \bigoplus_{\substack{\omega^2+\omega+1=0\\ \delta^2+1=0}}\MorseCat{3}.
\]
\end{example}

\begin{example}
Consider the Mori--Mukai model $MM_{2-33}$ on the algebraic torus
$(\mathbb{C}^*)^3$, with Laurent potential
\[
W_{MM_{2-33}}=x+\frac{x}{z}+y+z+\frac{1}{xy}.
\]
The critical equations are especially sparse:
\[
\frac{\partial W}{\partial z}=1-\frac{x}{z^2}=0,
\qquad
\frac{\partial W}{\partial y}=1-\frac{1}{xy^2}=0,
\qquad
\frac{\partial W}{\partial x}=1+\frac{1}{z}-\frac{1}{x^2y}=0.
\]
Thus
\[
x=z^2,
\qquad
xy^2=1,
\]
and the torus condition splits the solutions into the two branches
\[
y=\frac{1}{z}
\qquad\text{or}\qquad
y=-\frac{1}{z}.
\]
On the positive branch, the remaining equation becomes
\[
z^2(z+1)=1,
\qquad\text{that is}\qquad
p_+(z):=z^3+z^2-1=0,
\]
whereas on the negative branch it becomes
\[
z^2(z+1)=-1,
\qquad\text{that is}\qquad
p_-(z):=z^3+z^2+1=0.
\]

Equivalently, clearing denominators in the logarithmic critical equations and
saturating by $xyz$ gives the reduced ideal
\[
\left(
x-z^2,
y-z^2-z,
q(z)
\right),
\qquad
q(z)=z^6+2z^5+z^4-1=p_+(z)p_-(z).
\]
The cubics $p_+$ and $p_-$ are irreducible and square-free over $\mathbb{Q}$,
with discriminants $-23$ and $-31$, respectively, and they are coprime to each
other and to $z$. Hence all six solutions lie in the torus. The critical
points are
\[
P^+_{\lambda}=\left(\lambda^2,\frac{1}{\lambda},\lambda\right),
\qquad p_+(\lambda)=0,
\]
and
\[
P^-_{\mu}=\left(\mu^2,-\frac{1}{\mu},\mu\right),
\qquad p_-(\mu)=0.
\]

The critical values are given on both branches by the same expression
\[
W_{MM_{2-33}}=3z^2+4z
\]
after reduction modulo the corresponding cubic. More explicitly, the values
on the positive branch are the roots of
\[
w^3+w^2-18w-43=0,
\]
and the values on the negative branch are the roots of
\[
w^3+w^2+18w-11=0.
\]
Thus the six critical values are the roots of
\[
(w^3+w^2-18w-43)(w^3+w^2+18w-11)=0.
\]

The logarithmic Hessian is nondegenerate at every critical point. Indeed, on
the positive branch one has
\[
\det(\theta_i\theta_jW)(P^+_{\lambda})
=\frac{2(3-\lambda^2)}{\lambda},
\]
and on the negative branch one has
\[
\det(\theta_i\theta_jW)(P^-_{\mu})
=\frac{2(\mu^2+3)}{\mu}.
\]
The numerators are coprime to $p_+$ and $p_-$, respectively. Since logarithmic
coordinates are invertible on the torus, all six critical points are isolated
Morse singularities of type $A_1$. Consequently,
\[
\overline{MF((\mathbb{C}^*)^3,W_{MM_{2-33}})}
\cong \bigoplus_{\lambda:\,p_+(\lambda)=0}\MorseCat{3} \oplus \bigoplus_{\mu:\,p_-(\mu)=0}\MorseCat{3}.
\]
\end{example}

\begin{example}
Consider the Mori--Mukai model $MM_{2-34}$ on the algebraic torus
$(\mathbb{C}^*)^3$, with Laurent potential
\[
W_{MM_{2-34}}=x+\frac{1}{x}+y+z+\frac{1}{yz}.
\]
This potential separates as
\[
W_{MM_{2-34}}=\left(x+\frac{1}{x}\right)
+\left(y+z+\frac{1}{yz}\right),
\]
reflecting the product mirror structure of $\mathbb{P}^1\times\mathbb{P}^2$.
The critical equations decouple into
\[
1-\frac{1}{x^2}=0,
\qquad
1-\frac{1}{y^2z}=0,
\qquad
1-\frac{1}{yz^2}=0.
\]
Hence
\[
x^2=1,
\qquad
y^2z=1,
\qquad
yz^2=1.
\]
The last two equations imply $y=z$, and therefore $z^3=1$. Equivalently,
after clearing denominators in the logarithmic critical equations and
saturating by $xyz$, the reduced critical ideal is
\[
\left(
(x-1)(x+1),
y-z,
z^3-1
\right).
\]
Since $z^3-1$ is square-free and coprime to $z$, all critical points lie in
the torus. They are
\[
P_{\epsilon,\rho}=(\epsilon,\rho,\rho),
\qquad
\epsilon^2=1,
\quad
\rho^3=1.
\]
Thus the critical locus consists of six reduced points, namely the product of
the two critical points of the $\mathbb{P}^1$ mirror with the three critical
points of the $\mathbb{P}^2$ mirror.

The critical values are
\[
W_{MM_{2-34}}(P_{\epsilon,\rho})=2\epsilon+3\rho.
\]
For $\epsilon=1$ they are the roots of
\[
(w-2)^3-27=0,
\qquad\text{equivalently}\qquad
(w-5)(w^2-w+7)=0,
\]
whereas for $\epsilon=-1$ they are the roots of
\[
(w+2)^3-27=0,
\qquad\text{equivalently}\qquad
(w-1)(w^2+7w+19)=0.
\]

The logarithmic Hessian is block diagonal with respect to this splitting. Its
determinant is
\[
\det(\theta_i\theta_jW)
=\frac{(x^2+1)(y^2z^2+y+z)}{xyz}.
\]
At the critical points this becomes
\[
\det(\theta_i\theta_jW)(P_{\epsilon,\rho})=\frac{6}{\epsilon\rho},
\]
which is nonzero. Since logarithmic coordinates are invertible on the torus,
all six critical points are isolated Morse singularities of type $A_1$.
Consequently,
\[
\overline{MF((\mathbb{C}^*)^3,W_{MM_{2-34}})}
\cong \bigoplus_{\substack{\epsilon^2=1\\ \rho^3=1}}\MorseCat{3}.
\]
\end{example}

\begin{example}
Consider the Mori--Mukai model $MM_{2-35}$ on the algebraic torus
$(\mathbb{C}^*)^3$, with Laurent potential
\[
W_{MM_{2-35}}=x+\frac{x}{yz}+y+z+\frac{1}{x}.
\]
The critical equations are
\[
1+\frac{1}{yz}-\frac{1}{x^2}=0,
\qquad
1-\frac{x}{y^2z}=0,
\qquad
1-\frac{x}{yz^2}=0.
\]
The last two equations give
\[
x=y^2z=yz^2.
\]
Since we work on the torus, this forces $y=z$, and then $x=y^3$. Substituting
into the remaining equation gives
\[
1+\frac{1}{y^2}-\frac{1}{y^6}=0,
\qquad\text{or equivalently}\qquad
y^6+y^4-1=0.
\]
Thus, after clearing denominators in the logarithmic critical equations and
saturating by $xyz$, the reduced critical ideal is
\[
\left(
x-z^3,
y-z,
z^6+z^4-1
\right).
\]
The polynomial $z^6+z^4-1$ is square-free and coprime to $z$. Equivalently,
with $u=z^2$ it is governed by the irreducible cubic
\[
u^3+u^2-1=0,
\]
whose discriminant is $-23$. Hence the critical locus consists of the six
torus points
\[
P_{\lambda}=(\lambda^3,\lambda,\lambda),
\qquad
\lambda^6+\lambda^4-1=0.
\]

At these points the critical value reduces to
\[
W_{MM_{2-35}}(P_{\lambda})
=2\lambda(\lambda^2+2).
\]
Consequently, the critical values are the roots of
\[
w^6-8w^4-272w^2-1600=0.
\]
The logarithmic Hessian determinant is
\[
\det(\theta_i\theta_jW)
=\frac{x^3y+x^3z+x^2y^2z^2+x^2yz+xy+xz+y^2z^2}{xyz}.
\]
On the critical locus this becomes
\[
\det(\theta_i\theta_jW)(P_{\lambda})
=\frac{2(3-\lambda^4)}{\lambda}.
\]
Its numerator is coprime to $\lambda^6+\lambda^4-1$, so the Hessian is
nondegenerate at all six critical points. Since logarithmic coordinates are
invertible on the torus, all six critical points are isolated Morse
singularities of type $A_1$. Consequently,
\[
\overline{MF((\mathbb{C}^*)^3,W_{MM_{2-35}})}
\cong \bigoplus_{\lambda:\,\lambda^6+\lambda^4-1=0}\MorseCat{3}.
\]
\end{example}

\begin{example}
Consider the Mori--Mukai model $MM_{2-36}$ on the algebraic torus
$(\mathbb{C}^*)^3$, with Laurent potential
\[
W_{MM_{2-36}}=\frac{x^2}{yz}+x+y+z+\frac{1}{x}.
\]
The critical equations are
\[
\frac{2x}{yz}+1-\frac{1}{x^2}=0,
\qquad
1-\frac{x^2}{y^2z}=0,
\qquad
1-\frac{x^2}{yz^2}=0.
\]
The last two equations give
\[
x^2=y^2z=yz^2.
\]
Since $y$ and $z$ are invertible on the torus, this forces $y=z$, and then
\[
x^2=y^3.
\]
Writing
\[
t=\frac{x}{y},
\]
one obtains the torus parametrization
\[
x=t^3,
\qquad
y=z=t^2.
\]
Substitution into the remaining critical equation gives
\[
\frac{2}{t}+1-\frac{1}{t^6}=0,
\qquad\text{or equivalently}\qquad
p(t):=t^6+2t^5-1=0.
\]
The polynomial $p$ is irreducible and square-free over $\mathbb{Q}$
($\operatorname{disc}(p)=246656$), and it is coprime to $t$. Hence the
critical locus consists of the six torus points
\[
P_t=(t^3,t^2,t^2),
\qquad
p(t)=0.
\]
Equivalently, eliminating $t$ from this parametrization gives the reduced
critical ideal
\[
\left(
2x+z^5-4z^4-z^2,
y-z,
z^6-4z^5-2z^3+1
\right),
\]
which is the ideal obtained from the logarithmic critical equations after
clearing denominators and saturating by $xyz$.

At the critical points the critical value reduces to
\[
W_{MM_{2-36}}(P_t)=t^2(2t+5).
\]
Consequently, the six critical values are the roots of
\[
w^6-4w^5-12w^4+50w^3+48w^2-1000w+3061=0.
\]
The logarithmic Hessian determinant is
\[
\det(\theta_i\theta_jW)
=\frac{x^4y+x^4z+4x^3yz+x^2y^2z^2+x^2y+x^2z+y^2z^2}{xyz}.
\]
On the parametrized critical locus this becomes
\[
\det(\theta_i\theta_jW)(P_t)
=2(2t^5+3t-1).
\]
This polynomial is coprime to $p(t)$, so the Hessian is nondegenerate at all
six critical points. Since logarithmic coordinates are invertible on the
torus, all six critical points are isolated Morse singularities of type
$A_1$. Consequently,
\[
\overline{MF((\mathbb{C}^*)^3,W_{MM_{2-36}})}
\cong \bigoplus_{t:\,t^6+2t^5-1=0}\MorseCat{3}.
\]
\end{example}

\begin{example}
Consider the Mori--Mukai model $MM_{3-1}$ on the algebraic torus
$(\mathbb{C}^*)^3$, with Laurent potential
\begin{align*}
W_{MM_{3-1}}={}&\frac{x^2y^3}{z}+\frac{4xy^2}{z}+\frac{6y}{z} \\
&+\frac{4}{xz}+\frac{1}{x^2yz}+xy^2+3xy+x+2y+\frac{2}{y} \\
&+\frac{1}{x}+\frac{3}{xy}+\frac{1}{xy^2}+z+\frac{2z}{y}
+\frac{z}{y^2}.
\end{align*}
Set
\[
u=xy,
\qquad
v=\frac{z}{y},
\qquad
Y=y+2+\frac{1}{y},
\qquad
U=u+2+\frac{1}{u}.
\]
The change of variables $(x,y,z)\mapsto(u,y,v)$ is a torus automorphism, and
the potential becomes
\[
W_{MM_{3-1}}=vY+YU+\frac{U^2}{v}+U-6.
\]
Equivalently, it has the exact factorization
\[
W_{MM_{3-1}}+6
=
(v+U)\left(Y+\frac{U}{v}\right).
\]
The logarithmic critical equations in the coordinates $(u,y,v)$ factor as
\[
v^2Y-U^2=0,
\qquad
(v+U)(y^2-1)=0,
\qquad
\left(Y+\frac{2U}{v}+1\right)(u^2-1)=0.
\]

The branch $y=1$ gives $Y=4$. The first equation and the last equation force
$u=1$ and $v=\pm2$, producing the two isolated critical points
\[
P_+=(1,1,2),
\qquad
P_-=(1,1,-2)
\]
in the original $(x,y,z)$-coordinates. Their critical values are
\[
W_{MM_{3-1}}(P_+)=30,
\qquad
W_{MM_{3-1}}(P_-)=-2.
\]
At these two points the logarithmic Hessian determinants are
\[
\det(\theta_i\theta_jW)(P_+)=3456,
\qquad
\det(\theta_i\theta_jW)(P_-)=-128,
\]
so both isolated points are Morse singularities of type $A_1$.

The branch $y=-1$ gives $Y=0$. Then the first critical equation forces
$U=0$, hence $u=-1$, while $v$ remains free. In the original coordinates this
is the curve
\[
\Gamma_3=\{(1,-1,z):z\in\mathbb{C}^*\},
\qquad
W_{MM_{3-1}}|_{\Gamma_3}=-6.
\]

It remains to consider the branch $v=-U$. Since $v$ is invertible on the
torus, $U\neq0$, and the first critical equation gives $Y=1$. Thus
$y^2+y+1=0$. For each primitive cubic root $\zeta$ of unity, one obtains the
curve
\[
\Gamma_{\zeta}
=\left\{\left(\frac{u}{\zeta},\zeta,
-\zeta\left(u+2+\frac{1}{u}\right)\right):
u\in\mathbb{C}^*,\ u\neq -1\right\},
\qquad
\zeta^2+\zeta+1=0.
\]
Equivalently, in the Laurent coordinate ring of the $(u,y,v)$-torus these two
curves are cut out by
\[
y-\zeta=0,
\qquad
uv+(u+1)^2=0.
\]
Along both curves the potential is again constant:
\[
W_{MM_{3-1}}|_{\Gamma_\zeta}=-6.
\]

The three curve components admit global split relative quadratic models. For
$\Gamma_3$, put $a=u+1$ and $b=y+1$. Its formal neighborhood is
\[
\widehat X_{\Gamma_3}
\simeq
\operatorname{Spf}\!\left(\mathbb C[v^{\pm1}][[a,b]]\right).
\]
Since
\[
U=\frac{a^2}{u},
\qquad
Y=\frac{b^2}{y},
\]
the quadratic part of the shifted potential is
\[
\widehat W_{MM_{3-1}}+6
\equiv
-a^2-vb^2
\pmod{(a,b)^3}.
\]
Its normal Hessian determinant is $4v$, a unit on $\Gamma_3$. The relative
formal Morse lemma therefore gives a global Morse--Bott model on this
component. Its discriminant has the nontrivial square class
$[v]\in\mathbb C[v^{\pm1}]^*/\mathbb C[v^{\pm1}]^{*2}$, so the Clifford twist
cannot be suppressed.

For a primitive cubic root $\zeta$, set
\[
p=v+U,
\qquad
q=Y+\frac{U}{v}.
\]
The factorization above gives $W_{MM_{3-1}}+6=pq$. Moreover,
\[
q=(Y-1)+\frac{p}{v},
\qquad
Y-1=\frac{(y-\zeta)(y-\zeta^2)}{y}.
\]
Since $(y-\zeta^2)/y$ is a unit along $\Gamma_\zeta$, the functions $(p,q)$
generate the same ideal as $(v+U,y-\zeta)$. Writing
\[
A_\zeta=\mathbb C[u^{\pm1},(u+1)^{-1}],
\]
we obtain the global split completion
\[
\widehat X_{\Gamma_\zeta}
\simeq
\operatorname{Spf}\!\left(A_\zeta[[p,q]]\right),
\qquad
\widehat W_{MM_{3-1}}+6=pq.
\]

Consequently, Proposition~\ref{prop:morse_bott_clifford} applies to all three
curves, while the two isolated points contribute Morse categories:
\[
\overline{MF((\mathbb{C}^*)^3,W_{MM_{3-1}})}
\cong \MorseCat{3}\oplus \MorseCat{3}
\oplus \MBtwcat{\Gamma_3}{v}
\oplus\bigoplus_{\zeta^2+\zeta+1=0}\MBcat{\Gamma_\zeta}.
\]
\end{example}

\begin{example}
Consider the Mori--Mukai model $MM_{3-2}$ on the algebraic torus
$(\mathbb{C}^*)^3$, with Laurent potential
\begin{align*}
W_{MM_{3-2}}={}&xy z^2+xy z+3xz+x+\frac{3x}{y}+\frac{x}{y^2z}
+3yz+y+z+\frac{1}{y} \\
&+\frac{3}{yz}+\frac{3y}{x}+\frac{1}{x}+\frac{3}{xz}
+\frac{y}{x^2z}.
\end{align*}
Set
\[
U=1+\frac{x+y}{xyz}.
\]
Then the potential has the compact form
\[
W_{MM_{3-2}}=xy z^2U^3+z(xy+1)U-6.
\]
Equivalently, writing
\[
p=xy,
\qquad
a=x+y,
\]
one has
\[
W_{MM_{3-2}}=F(p,a,z),
\qquad
F=pz^2\left(1+\frac{a}{pz}\right)^3
+z(p+1)\left(1+\frac{a}{pz}\right)-6.
\]
Since
\[
x\frac{\partial W}{\partial x}-y\frac{\partial W}{\partial y}
=(x-y)\frac{\partial F}{\partial a},
\]
the critical equations split into the branch $x=y$ and the symmetric branch
$\partial_aF=0$.

On the symmetric branch, solving
\[
\partial_pF=\partial_aF=\partial_zF=0
\]
on the torus gives
\[
p=-1,
\qquad
a=z.
\]
Thus one obtains the critical curve
\[
\Gamma=\{xy+1=0,\\ z-x-y=0\}
\subset(\mathbb{C}^*)^3,
\]
or, parametrically,
\[
\Gamma=\left\{\left(t,-\frac{1}{t},t-\frac{1}{t}\right):
t\in\mathbb{C}^*,\ t^2\neq1\right\}.
\]
Along this curve
\[
W_{MM_{3-2}}|_{\Gamma}=-6.
\]

On the branch $x=y$, the factor $z^2+4$ corresponds precisely to the two
intersection points of this branch with $\Gamma$. Removing the curve component,
the isolated critical locus is the disjoint union of
\[
P^+_{\lambda}=(1,1,\lambda),
\qquad
p_+(\lambda)=0,
\qquad
p_+(z)=z^3+4z^2-4,
\]
and
\[
P^-_{\mu}=(-1,-1,\mu),
\qquad
p_-(\mu)=0,
\qquad
p_-(z)=z^3-2z^2+4.
\]
Both cubics are square-free and coprime to $z$, with discriminants $592$ and
$-304$, respectively, so all six isolated points lie in the torus.

The critical values on the two isolated branches are given by
\[
W_{MM_{3-2}}(P^+_{\lambda})=3\lambda^2+16\lambda+10
\quad\text{modulo }p_+(\lambda),
\]
and
\[
W_{MM_{3-2}}(P^-_{\mu})=3\mu^2-8\mu+2
\quad\text{modulo }p_-(\mu).
\]
Equivalently, the corresponding critical values are the roots of
\[
w^3-14w^2-308w-1048=0
\]
and
\[
w^3-2w^2-148w-648=0.
\]
The logarithmic Hessian is nondegenerate at all isolated points. More
precisely, modulo $p_+$ one has
\[
\det(\theta_i\theta_jW)(P^+_{\lambda})
=8(22\lambda^2+107\lambda+96),
\]
and modulo $p_-$ one has
\[
\det(\theta_i\theta_jW)(P^-_{\mu})
=-8(19\mu^2-59\mu+66),
\]
with the displayed numerators coprime to the corresponding cubic.

The relative formal model along $\Gamma$ can be identified globally. Set
\[
r=zU,
\qquad
s=xy+1+xyzU^2.
\]
A direct calculation gives the exact factorization
\[
W_{MM_{3-2}}+6=rs.
\]
Indeed, if
\[
F=xy+1,
\qquad
H=xyz+x+y,
\]
then the ideal of $\Gamma$ is $(F,H)$, because
\[
H=zF-(z-x-y).
\]
Moreover,
\[
r=\frac{H}{xy},
\qquad
s-F=\frac{H^2}{xyz}.
\]
Thus $(r,s)$ and $(H,F)$ induce the same basis of the conormal bundle and
generate the same ideal in the formal completion. Writing
\[
A_{\Gamma}=\mathbb C[t^{\pm1},(t^2-1)^{-1}],
\]
we obtain
\[
\widehat X_{\Gamma}
\simeq
\operatorname{Spf}\!\left(A_{\Gamma}[[r,s]]\right),
\qquad
\widehat W_{MM_{3-2}}+6=rs.
\]
Hence $\Gamma$ is a global split Morse--Bott critical curve.

Consequently, Proposition~\ref{prop:morse_bott_clifford} applies to the curve,
while the six isolated points contribute Morse categories:
\[
\overline{MF((\mathbb{C}^*)^3,W_{MM_{3-2}})}
\cong \MBcat{\Gamma} \oplus\bigoplus_{\lambda:\,p_+(\lambda)=0}\MorseCat{3} \oplus\bigoplus_{\mu:\,p_-(\mu)=0}\MorseCat{3}.
\]
\end{example}

\begin{example}
Consider the Mori--Mukai model $MM_{3-3}$ on the algebraic torus
$(\mathbb{C}^*)^3$, with Laurent potential
\begin{align*}
W_{MM_{3-3}}={}&\frac{xy}{z}+x+\frac{x}{z}+\frac{2x}{y}
+\frac{x}{y^2}+y+\frac{y}{z}+z+\frac{3}{y} \\
&+\frac{2y}{x}+\frac{3}{x}+\frac{y}{x^2}.
\end{align*}
Let
\[
F=xy+x+y,
\qquad
A=\frac{F}{xy}=1+\frac{1}{x}+\frac{1}{y}.
\]
Then the potential admits the compact form
\[
W_{MM_{3-3}}
=z+\frac{F}{z}+(x+y)A^2-4
=z+\frac{F}{z}+\frac{(x+y)F^2}{x^2y^2}-4.
\]
The $z$-derivative gives the height relation
\[
z^2=F=xy+x+y.
\]
Clearing denominators in the logarithmic critical equations and saturating by
$xyz$ gives a reduced critical ideal with the decomposition
\[
I_{\mathrm{crit}}=I_{\Delta}\cap I_{\mathrm{off}},
\]
where
\[
I_{\Delta}=\left(
x-y,
(4-x^2)z-x^2(x+1),
q(x)
\right),
\qquad
q(x)=9x^3+16x^2-16x-32,
\]
and
\[
I_{\mathrm{off}}=(z-4,\ x+y+16,\ xy-32).
\]
Thus the diagonal branch consists of the three points
\[
P_{\alpha}=\left(
\alpha,
\alpha,
\frac{\alpha^2(\alpha+1)}{4-\alpha^2}
\right),
\qquad
q(\alpha)=0.
\]
The polynomial $q$ is irreducible and square-free over $\mathbb{Q}$, and it is
coprime to $x(4-x^2)$, so all three points lie in the torus. Notice that the
relation above fixes the sign of $z$; squaring the diagonal equation would
introduce the extraneous opposite sign.

The off-diagonal branch consists of the two points
\[
Q_{\pm}=(-8\pm4\sqrt{2},\ -8\mp4\sqrt{2},\ 4),
\]
equivalently the two solutions of
\[
z=4,
\qquad
x+y=-16,
\qquad
xy=32.
\]
Hence the torus critical locus is finite and consists of five points.

The diagonal critical values are given in $\mathbb{Q}[\alpha]/(q)$ by
\[
W_{MM_{3-3}}(P_{\alpha})
=\frac{9\alpha^2+88\alpha+64}{12},
\]
or, equivalently, they are the roots of
\[
w^3-8w^2-128w-368=0.
\]
On the off-diagonal branch one has
\[
W_{MM_{3-3}}(Q_+)=W_{MM_{3-3}}(Q_-)=0.
\]

The logarithmic Hessian is nondegenerate at every critical point. On the
diagonal branch one has, modulo $q(\alpha)$,
\[
\det(\theta_i\theta_jW)(P_{\alpha})
=\frac{2}{3}(63\alpha^2+238\alpha+232),
\]
and this numerator is coprime to $q$. On the off-diagonal branch,
\[
\det(\theta_i\theta_jW)(Q_\pm)=-64.
\]
Therefore all five critical points are isolated Morse singularities of type
$A_1$. Consequently,
\[
\overline{MF((\mathbb{C}^*)^3,W_{MM_{3-3}})}
\cong \bigoplus_{\alpha:\,q(\alpha)=0}\MorseCat{3} \oplus \bigoplus_{\varepsilon=\pm}\MorseCat{3}.
\]
\end{example}

\begin{example}
Consider the Mori--Mukai model $MM_{3-4}$ on the algebraic torus
$(\mathbb{C}^*)^3$, with Laurent potential
\begin{align*}
W_{MM_{3-4}}={}&x+\frac{x}{y}+2yz+y+z+\frac{2}{y}
+\frac{2}{yz}+\frac{1}{y^2z} \\
&+\frac{y^2z^2}{x}+\frac{4yz}{x}+\frac{6}{x}
+\frac{4}{xyz}+\frac{1}{xy^2z^2}.
\end{align*}
Set
\[
U=yz+2+\frac{1}{yz}.
\]
Then the potential has the compact form
\[
W_{MM_{3-4}}
=x\left(1+\frac{1}{y}\right)+\frac{U^2}{x}
+U\left(2+\frac{1}{y}\right)+y-4.
\]
The $z$-derivative factors as
\[
\frac{\partial W}{\partial z}
=\left(\frac{2U}{x}+2+\frac{1}{y}\right)
\left(y-\frac{1}{yz^2}\right).
\]
The first factor gives no torus critical points. Indeed, if
\[
\frac{2U}{x}+2+\frac{1}{y}=0,
\]
then the case $2y+1=0$ would force $U=0$, contradicting the
$x$-equation. Thus one may solve
\[
x=-\frac{2yU}{2y+1}.
\]
Combining this with
\[
\frac{\partial W}{\partial x}
=1+\frac{1}{y}-\frac{U^2}{x^2}=0
\]
gives
\[
4y(y+1)=(2y+1)^2,
\]
which is the contradiction $0=1$. Thus the first factor is a phantom branch.

The second factor gives
\[
y^2z^2=1,
\qquad\text{hence}\qquad
yz=\pm1.
\]
On the branch $yz=1$, one has $U=4$. Since $\partial U/\partial y=0$ along
this branch, the remaining equations reduce to
\[
1-\frac{x+4}{y^2}=0,
\qquad
1+\frac{1}{y}-\frac{16}{x^2}=0.
\]
Therefore
\[
x=y^2-4,
\qquad
(y^2-4)^2(y+1)=16y,
\]
or equivalently
\[
p(y):=y^5+y^4-8y^3-8y^2+16=0.
\]
The polynomial $p$ is irreducible and square-free over $\mathbb{Q}$, and is
coprime to $y(y^2-4)$. Hence this branch gives the five torus critical points
\[
P_\beta=\left(\beta^2-4,\ \beta,\ \frac{1}{\beta}\right),
\qquad
p(\beta)=0.
\]

There is also one exceptional point on the branch $yz=-1$. Here $U=0$, the
$x$-equation gives $y=-1$, and then $z=1$. The $y$-equation gives $x=1$.
Thus
\[
P_0=(1,-1,1),
\qquad
W_{MM_{3-4}}(P_0)=-5.
\]

The critical values on the five-point branch are given in
$\mathbb{Q}[\beta]/(p)$ by
\[
W_{MM_{3-4}}(P_\beta)
=\frac{(\beta^2+4)(\beta^2+\beta-4)}{4},
\]
or, equivalently, they are the roots of
\[
w^5-11w^4-140w^3+176w^2+4224w+7488=0.
\]

The logarithmic Hessian is nondegenerate at every critical point. On the
five-point branch one has, modulo $p(\beta)$,
\[
\det(\theta_i\theta_jW)(P_\beta)
=2(\beta^4+9\beta^3+3\beta^2-4\beta-16),
\]
and this numerator is coprime to $p$. At the exceptional point,
\[
\det(\theta_i\theta_jW)(P_0)=2.
\]
Therefore all six critical points are isolated Morse singularities of type
$A_1$. Consequently,
\[
\overline{MF((\mathbb{C}^*)^3,W_{MM_{3-4}})}
\cong \MorseCat{3} \oplus \bigoplus_{\beta:\,p(\beta)=0}\MorseCat{3}.
\]
\end{example}

\begin{example}
Consider the Mori--Mukai model $MM_{3-5}$ on the algebraic torus
$(\mathbb{C}^*)^3$, with Laurent potential
\begin{align*}
W_{MM_{3-5}}={}&xyz+xz^2+2xz+x+y+2z+\frac{2}{z}
+\frac{1}{y}+\frac{1}{yz} \\
&+\frac{1}{x}+\frac{2}{xz}+\frac{1}{xz^2}.
\end{align*}
The $y$-derivative is
\[
\frac{\partial W}{\partial y}
=xz+1-\frac{1}{y^2}-\frac{1}{y^2z},
\]
so every torus critical point satisfies
\[
y^2=\frac{z+1}{z(xz+1)}.
\]
Eliminating after clearing denominators and saturating by $xyz$ gives the
reduced critical ideal
\[
\left(
x-\frac{1}{z^2},\ y-1,\ p_+(z)
\right)
\cap
\left(
x-\frac{1}{z^2},\ y+1,\ p_-(z)
\right),
\]
where
\[
p_+(z)=z^4+2z^3-3z-1,
\qquad
p_-(z)=z^4+2z^3-z-1.
\]
Equivalently, the critical locus splits into the two branches
\[
y=1,
\qquad
x=\frac{1}{z^2},
\qquad
p_+(z)=0,
\]
and
\[
y=-1,
\qquad
x=\frac{1}{z^2},
\qquad
p_-(z)=0.
\]
Both quartics are irreducible and square-free over $\mathbb{Q}$, and both are
coprime to $z$, so all eight points lie in the torus. The reduction is also
visible directly: after substituting $y=\pm1$ and $x=z^{-2}$ into the
$x$-equation, one obtains precisely $p_+(z)=0$ on the $y=1$ branch and
$p_-(z)=0$ on the $y=-1$ branch, while the remaining $z$-equation becomes
identically zero.

Thus the critical points are
\[
P^+_\lambda=\left(\frac{1}{\lambda^2},1,\lambda\right),
\qquad
p_+(\lambda)=0,
\]
and
\[
P^-_\mu=\left(\frac{1}{\mu^2},-1,\mu\right),
\qquad
p_-(\mu)=0.
\]

The critical values on the two branches are given by
\[
W_{MM_{3-5}}(P^+_\lambda)=3\lambda^3+8\lambda^2+6\lambda-5
\quad\text{modulo }p_+(\lambda),
\]
and
\[
W_{MM_{3-5}}(P^-_\mu)=\mu^3+4\mu^2+6\mu-1
\quad\text{modulo }p_-(\mu).
\]
Equivalently, these critical values are the roots of
\[
w^4-3w^3-152w^2-768w-1123=0
\]
and
\[
w^4+5w^3-44w^2-360w-691=0,
\]
respectively.

The logarithmic Hessian is nondegenerate at every critical point. More
precisely, modulo $p_+$ one has
\[
\det(\theta_i\theta_jW)(P^+_\lambda)
=2(11\lambda^3+39\lambda^2+52\lambda+23),
\]
and modulo $p_-$ one has
\[
\det(\theta_i\theta_jW)(P^-_\mu)
=-2(17\mu^3+51\mu^2+48\mu+23),
\]
with the displayed numerators coprime to the corresponding quartic. Therefore
all eight critical points are isolated Morse singularities of type $A_1$.
Consequently,
\[
\overline{MF((\mathbb{C}^*)^3,W_{MM_{3-5}})}
\cong \bigoplus_{\lambda:\,p_+(\lambda)=0}\MorseCat{3} \oplus \bigoplus_{\mu:\,p_-(\mu)=0}\MorseCat{3}.
\]
\end{example}

\begin{example}
Consider the Mori--Mukai model $MM_{3-6}$ on the algebraic torus
$(\mathbb{C}^*)^3$, with Laurent potential
\[
W_{MM_{3-6}}=x+\frac{xz}{y}+\frac{2x}{y}+\frac{x}{yz}
+y+z+\frac{1}{y}+\frac{1}{yz}
+\frac{2y}{x}+\frac{2}{x}+\frac{y}{x^2}.
\]
The $z$-derivative gives
\[
\frac{\partial W}{\partial z}
=\frac{x}{y}-\frac{x}{yz^2}-\frac{1}{yz^2}+1,
\]
and hence every torus critical point satisfies
\[
z^2(x+y)=x+1.
\]
The full critical system is slightly more rigid than this equation alone
suggests. Clearing denominators in the logarithmic critical equations and
saturating by $xyz$ gives the reduced critical ideal
\[
I_{\mathrm{crit}}=I_A\cap I_B,
\]
where
\[
I_A=\left(
x-z^2-2z+1,
\ y+\frac{z^2-6}{3},
\ r_A(z)
\right),
\qquad
r_A(z)=z^3+3z^2-3,
\]
and
\fitdisplay{I_B=\left(
x-\frac{z^2-2z+5}{5},
\ y-\frac{z^2+2z+2}{5},
\ r_B(z)
\right),
\qquad
r_B(z)=z^3+z^2+4z+5}
Thus the critical locus consists of the two triples of points
\[
P^A_\lambda=\left(
\lambda^2+2\lambda-1,
-\frac{\lambda^2-6}{3},
\lambda
\right),
\qquad
r_A(\lambda)=0,
\]
and
\[
P^B_\mu=\left(
\frac{\mu^2-2\mu+5}{5},
\frac{\mu^2+2\mu+2}{5},
\mu
\right),
\qquad
r_B(\mu)=0.
\]
Both cubics are irreducible and square-free over $\mathbb{Q}$, and the
displayed coordinate functions are coprime to the corresponding cubic; hence
all six points lie in the torus.

The critical values on the two branches are given by
\[
W_{MM_{3-6}}(P^A_\lambda)=3(\lambda^2+3\lambda+1)
\quad\text{modulo }r_A(\lambda),
\]
and
\[
W_{MM_{3-6}}(P^B_\mu)=\frac{7\mu^2+\mu+5}{5}
\quad\text{modulo }r_B(\mu).
\]
Equivalently, these critical values are the roots of
\[
w^3-9w^2-54w-27=0
\]
and
\[
w^3+7w^2-2w-67=0,
\]
respectively.

The logarithmic Hessian is nondegenerate at every critical point. More
precisely, modulo $r_A$ one has
\[
\det(\theta_i\theta_jW)(P^A_\lambda)
=18(\lambda^2+5\lambda+5),
\]
and modulo $r_B$ one has
\[
\det(\theta_i\theta_jW)(P^B_\mu)
=-\frac{2}{5}(33\mu^2-11\mu+155),
\]
with the displayed numerators coprime to the corresponding cubic. Therefore
all six critical points are isolated Morse singularities of type $A_1$.
Consequently,
\[
\overline{MF((\mathbb{C}^*)^3,W_{MM_{3-6}})}
\cong \bigoplus_{\lambda:\,r_A(\lambda)=0}\MorseCat{3} \oplus \bigoplus_{\mu:\,r_B(\mu)=0}\MorseCat{3}.
\]
\end{example}

\begin{example}
Consider the CCGK Laurent polynomial for $MM_{3-7}$:
\[
W_{MM_{3-7}}
=\frac{xy}{z}+x+\frac{x}{z}+\frac{x}{y}+y+z
+\frac2y+\frac yx+\frac2x+\frac1{xy}.
\]
The saturated logarithmic critical ideal has lexicographic generators
\[
2x+3z-7,
\qquad y+3z-6,
\qquad z^2-6z+7.
\]
Thus the critical points are
\[
P_\sigma=\left(\frac{7-3\sigma}{2},6-3\sigma,\sigma\right),
\qquad \sigma^2-6\sigma+7=0,
\]
with critical values $18-4\sigma$. The logarithmic Hessian determinant reduces
to $-6(9\sigma-35)$ modulo $\sigma^2-6\sigma+7$, hence is nonzero at both
points. Consequently,
\[
\overline{MF((\mathbb{C}^*)^3,W_{MM_{3-7}})}
\cong \bigoplus_{\sigma:\,\sigma^2-6\sigma+7=0}\MorseCat{3}.
\]
\end{example}

\begin{example}
Consider the Mori--Mukai model $MM_{3-8}$ on the algebraic torus
$(\mathbb{C}^*)^3$, with Laurent potential
\begin{align*}
W_{MM_{3-8}}={}&x+\frac{x}{z}+y+z+\frac{1}{z}
+\frac{2}{y}+\frac{1}{yz} \\
&+\frac{y}{x}+\frac{3}{x}+\frac{3}{xy}+\frac{1}{xy^2}.
\end{align*}
The $x$-derivative uses the identity
\[
y+3+\frac{3}{y}+\frac{1}{y^2}=\frac{(y+1)^3}{y^2},
\]
and gives
\[
\frac{z+1}{z}=\frac{(y+1)^3}{x^2y^2}.
\]
The $z$-derivative gives
\[
z^2=x+1+\frac{1}{y},
\qquad\text{hence}\qquad
x=\frac{yz^2-y-1}{y}.
\]
Eliminating $x$ yields the two plane equations
\[
A(y,z):=(z+1)(yz^2-y-1)^2-z(y+1)^3=0
\]
and
\[
B(y,z):=z^3y^3-zy^2-(2z^3+z^2+z-1)y+1=0.
\]
Their resultant with respect to $y$ is
\[
\operatorname{Res}_y(A,B)
=z^{10}(z+1)q(z),
\]
where
\[
q(z)=z^7+z^6-3z^5-3z^4-5z^3+3z^2-z-1.
\]
The factor $z^{10}$ is excluded by the torus condition. The factor $z+1$
gives the single torus point
\[
P_0=(1,-1,-1),
\qquad
W_{MM_{3-8}}(P_0)=-4.
\]
The remaining seven points are
\fitdisplay{P_\lambda=\left(
\frac{\lambda^5+\lambda^4-4\lambda^3-\lambda-1}{4},
\frac{\lambda^6+\lambda^5-2\lambda^4-2\lambda^3-7\lambda^2+\lambda-4}{4},
\lambda
\right),
\qquad
q(\lambda)=0}
The polynomial $q$ is irreducible and square-free over $\mathbb{Q}$, and is
coprime to the displayed coordinate numerators and to $z(z+1)$. Hence the
intersection of the two master curves in the torus consists of exactly eight
points.

On the seven-point branch the critical values are the roots of
\[
w^7-100w^5-584w^4-288w^3+6844w^2+21488w+20572=0.
\]
The logarithmic Hessian is nondegenerate at every critical point. At the
exceptional point one has
\[
\det(\theta_i\theta_jW)(P_0)=2,
\]
whereas on the seven-point branch, modulo $q(\lambda)$,
\[
\det(\theta_i\theta_jW)(P_\lambda)
=\frac{1}{2}(3\lambda^6+6\lambda^5-3\lambda^4+4\lambda^3
-15\lambda^2+6\lambda+3),
\]
and the displayed numerator is coprime to $q$. Therefore all eight critical
points are isolated Morse singularities of type $A_1$. Consequently,
\[
\overline{MF((\mathbb{C}^*)^3,W_{MM_{3-8}})}
\cong \MorseCat{3} \oplus \bigoplus_{\lambda:\,q(\lambda)=0}\MorseCat{3}.
\]
\end{example}

\begin{example}
Consider the Mori--Mukai model $MM_{3-9}$ on the algebraic torus
$(\mathbb{C}^*)^3$, with Laurent potential
\[
W_{MM_{3-9}}=\frac{x^2}{yz}+x+\frac{4x}{yz}+y+z
+\frac{6}{yz}+\frac{1}{x}+\frac{4}{xyz}+\frac{1}{x^2yz}.
\]
Set
\[
U=x+2+\frac{1}{x}=\frac{(x+1)^2}{x}.
\]
Then
\[
W_{MM_{3-9}}=\frac{U^2}{yz}+x+y+z+\frac{1}{x}.
\]
The $y$- and $z$-derivatives give
\[
y^2z=U^2,
\qquad
yz^2=U^2.
\]
Since $y$ and $z$ are invertible on the torus, it follows that
\[
y=z,
\qquad
y^3=U^2.
\]
The $x$-derivative factors as
\[
\frac{\partial W}{\partial x}
=\left(1-\frac{1}{x^2}\right)
\left(1+\frac{2U}{y^2}\right).
\]
Thus the critical equations split into two branches.

On the first branch $x^2=1$. The value $x=-1$ gives $U=0$ and hence
$y=0$, which is excluded from the torus. Therefore $x=1$, $U=4$, and
\[
y=z,
\qquad
y^3=16.
\]
This gives the three critical points
\[
P_\eta=(1,\eta,\eta),
\qquad
\eta^3-16=0.
\]

On the second branch one has $y^2=-2U$. Combining this with $y^3=U^2$ gives
$U=-8$, and hence
\[
x^2+10x+1=0.
\]
The equation $y^2=-2U$ then gives $y^2=16$, while $y^3=64$ forces $y=4$.
Thus the second branch consists of the two points
\[
Q_\alpha=(\alpha,4,4),
\qquad
\alpha^2+10\alpha+1=0.
\]
Consequently, the torus critical locus consists of five points.

The critical values on the first branch are
\[
W_{MM_{3-9}}(P_\eta)=3\eta+2
\quad\text{modulo }\eta^3-16,
\]
or, equivalently, the roots of
\[
w^3-6w^2+12w-440=0.
\]
On the second branch one has the constant critical value
\[
W_{MM_{3-9}}(Q_\alpha)=2.
\]

The logarithmic Hessian is nondegenerate at every critical point. On the first
branch one has, modulo $\eta^3-16$,
\[
\det(\theta_i\theta_jW)(P_\eta)=6(\eta^2+8),
\]
and this numerator is coprime to $\eta^3-16$. On the second branch one has,
modulo $\alpha^2+10\alpha+1$,
\[
\det(\theta_i\theta_jW)(Q_\alpha)=-192.
\]
Therefore all five critical points are isolated Morse singularities of type
$A_1$. Consequently,
\[
\overline{MF((\mathbb{C}^*)^3,W_{MM_{3-9}})}
\cong \bigoplus_{\eta:\,\eta^3=16}\MorseCat{3} \oplus \bigoplus_{\alpha:\,\alpha^2+10\alpha+1=0}\MorseCat{3}.
\]
\end{example}

\begin{example}
Consider the Mori--Mukai model $MM_{3-10}$ on the algebraic torus
$(\mathbb{C}^*)^3$, with Laurent potential
\[
W_{MM_{3-10}}=\frac{xy}{z}+x+\frac{x}{y}+y+z+
\frac{2}{y}+\frac{y}{x}+\frac{2}{x}+\frac{1}{xy}.
\]
The logarithmic critical equations, after clearing denominators, are
\begin{align*}
0={}&x^2y^2+x^2yz+x^2z-y^2z-2yz-z, \\
0={}&x^2y^2-x^2z+xy^2z-2xz+y^2z-z, \\
0={}&z^2-xy.
\end{align*}
Since $x,y,z$ are invertible on the torus, the last equation gives
$x=z^2/y$. Substitution into the first two equations gives the pair
\begin{align*}
A(y,z)={}&y^4+2y^3-y^2z^3+y^2-yz^4-z^4=0, \\
B(y,z)={}&y^4+y^3z^2+y^2z^3-y^2-2yz^2-z^4=0.
\end{align*}
Their sum factors as
\[
A+B=(y-z)(y+z)(2y^2+yz^2+2y+2z^2).
\]
Thus the apparent diagonal collapse $y=z$ gives only one component of the
critical locus. The saturated critical ideal in the Laurent coordinate ring is
\[
I_{\mathrm{crit}}=I_A\cap I_0\cap I_Q,
\]
where
\[
I_A=(x-z,\ y-z,\ 2z^3-2z-1),
\]
\[
I_0=(2x+1,\ 2y+1,\ 2z-1),
\]
and
\[
I_Q=(x+y+9,\ z-4,\ y^2+9y+16).
\]
Equivalently, the critical locus consists of the three diagonal points
\[
P_\lambda=(\lambda,\lambda,\lambda),
\qquad
2\lambda^3-2\lambda-1=0,
\]
the isolated off-diagonal point
\[
Q_0=\left(-\frac12,-\frac12,\frac12\right),
\]
and the two further off-diagonal points
\[
Q_\alpha=(-\alpha-9,\alpha,4),
\qquad
\alpha^2+9\alpha+16=0.
\]
All six points lie in the torus.

The critical values are
\[
W_{MM_{3-10}}(P_\lambda)=2(2\lambda^2+3\lambda-1)
\quad\text{modulo }2\lambda^3-2\lambda-1,
\]
whose eliminant is
\[
w^3-2w^2-76w-212=0,
\]
and
\[
W_{MM_{3-10}}(Q_0)=-2,
\qquad
W_{MM_{3-10}}(Q_\alpha)=1.
\]
The logarithmic Hessian is nondegenerate at every critical point. More
precisely, on the diagonal branch
\[
\det(\theta_i\theta_jW)(P_\lambda)
=2(4\lambda+1)(4\lambda+3),
\]
which is coprime to $2\lambda^3-2\lambda-1$, while
\[
\det(\theta_i\theta_jW)(Q_0)=-2,
\qquad
\det(\theta_i\theta_jW)(Q_\alpha)=-34.
\]
Therefore all six critical points are isolated Morse singularities of type
$A_1$. Consequently,
\[
\overline{MF((\mathbb{C}^*)^3,W_{MM_{3-10}})}
\cong \bigoplus_{\lambda:\,2\lambda^3-2\lambda-1=0}\MorseCat{3} \oplus \MorseCat{3} \oplus \bigoplus_{\alpha:\,\alpha^2+9\alpha+16=0}\MorseCat{3}.
\]
\end{example}

\begin{example}
Consider the Mori--Mukai model $MM_{3-11}$ on the algebraic torus
$(\mathbb{C}^*)^3$, with Laurent potential
\[
W_{MM_{3-11}}=x+\frac{x}{z}+\frac{x}{yz}+y+z+
\frac{1}{y}+\frac{y}{x}+\frac{2}{x}+\frac{1}{xy}.
\]
The logarithmic critical equations, after clearing denominators, are
\begin{align*}
0={}&x^2yz+x^2y+x^2-y^2z-2yz-z, \\
0={}&-x^2+xy^2z-xz+y^2z-z, \\
0={}&yz^2-xy-x.
\end{align*}
The last equation gives
\[
z^2=\frac{x(y+1)}{y}.
\]
The $x$-equation can be written as
\[
x^2(yz+y+1)=z(y+1)^2.
\]
In particular $y+1\neq0$ at every torus critical point. Introduce the master
parameter $v$ by
\[
x=\frac{y+1}{v}.
\]
Then $v$ is invertible, and the $x$-equation gives
\[
z=\frac{y+1}{v^2-y}.
\]
Substituting this expression in the $z$-equation gives the first quadratic
relation
\[
y^2-y(2v^2+v)+v^4=0.
\]
The $y$-equation is equivalently
\[
(y^2-1)(x+1)=\frac{x^2}{z}.
\]
After the same substitution, and using $y+1\neq0$, this becomes the second
quadratic relation
\[
vy^2+y(v^2+1)-2v^2-v=0.
\]
Eliminating $y$ from the two quadratics gives
\[
\operatorname{Res}_y=v^4\Phi(v),
\]
where the factor $v^4$ is excluded by the torus condition and
\[
\Phi(v)=v^6+2v^5+2v^4-2v^3-11v^2-10v-1.
\]
Equivalently, subtracting the second quadratic from $v$ times the first one
gives
\[
y=\frac{v(v^4+2v+1)}{2v^3+2v^2+1}.
\]
Thus the critical points are
\[
P_\nu=(X(\nu),Y(\nu),Z(\nu)),
\qquad
\Phi(\nu)=0,
\]
where
\begin{align*}
Y(v)&=\frac{v(v^4+2v+1)}{2v^3+2v^2+1}, \\
X(v)&=\frac{v^5+2v^3+4v^2+v+1}{v(2v^3+2v^2+1)}, \\
Z(v)&=\frac{v^5+2v^3+4v^2+v+1}{v(v^4+2v^3-v-1)}.
\end{align*}
The polynomial $\Phi$ is square-free over $\mathbb{Q}$, and the displayed
coordinate numerators and denominators are coprime to $\Phi$. Hence all six
critical points lie in the torus and no additional branch survives the
saturation.

The critical values are the roots of
\[
w^6+2w^5-43w^4-426w^3-1662w^2-2972w-167=0.
\]
Equivalently, modulo $\Phi(\nu)$ one has
\[
W_{MM_{3-11}}(P_\nu)=
\frac{60\nu^5+37\nu^4+80\nu^3-129\nu^2-204\nu-25}{61}.
\]
Finally, the logarithmic Hessian determinant reduces modulo $\Phi$ to
\[
\det(\theta_i\theta_jW)(P_\nu)
=\frac{2}{61}\left(10\nu^5+179\nu^4+115\nu^3
+314\nu^2-217\nu-360\right),
\]
and the displayed numerator is coprime to $\Phi$. Therefore all six critical
points are isolated Morse singularities of type $A_1$. Consequently,
\[
\overline{MF((\mathbb{C}^*)^3,W_{MM_{3-11}})}
\cong \bigoplus_{\nu:\,\Phi(\nu)=0}\MorseCat{3}.
\]
\end{example}

\begin{example}
Consider the Mori--Mukai model $MM_{3-12}$ on the algebraic torus
$(\mathbb{C}^*)^3$, with Laurent potential
\[
W_{MM_{3-12}}=xy+x+y+z+\frac{1}{y}+\frac{y}{xz}
+\frac{3}{x}+\frac{2}{xz}+\frac{3}{x^2z}+\frac{1}{x^3z^2}.
\]
The coefficients $1,3,3,1$ give the compact form
\[
z+\frac{3}{x}+\frac{3}{x^2z}+\frac{1}{x^3z^2}
=z\left(1+\frac{1}{xz}\right)^3.
\]
Set
\[
U=1+\frac{1}{xz}.
\]
Then $U-1$ is invertible on the torus, $z=1/(x(U-1))$, and
\[
W_{MM_{3-12}}=zU^3+U(y+2)+x(y+1)+\frac{1}{y}-2.
\]
In the coordinates $(x,y,U)$ the critical equations are equivalent to
\begin{align*}
0={}&y^2(x+U)-1, \\
0={}&x^2(y+1)(U-1)-U^3, \\
0={}&x(y+2)(U-1)^2-U^2(3-2U).
\end{align*}
Eliminating $x$ from the last two equations gives the first master curve
\[
A(U,y):=(U-1)^3(y+2)^2-U(3-2U)^2(y+1)=0,
\]
while eliminating $x$ using the first and third equations gives the second
master curve
\[
B(U,y):=U(U-1)^2y^3+U(2-U)y^2-(U-1)^2y-2(U-1)^2=0.
\]
Their resultant with respect to $y$ is
\[
\operatorname{Res}_y(A,B)
=U(U-1)^2(2U-3)^2\Phi_A(U)\Phi_B(U),
\]
where the extraneous factors are removed by the full critical system and
\[
\Phi_A(U)=U^4-7U^3+16U^2-12U+1,
\]
\[
\Phi_B(U)=U^4-U^3+2U^2-4U+1.
\]
The two quartics are irreducible and square-free over $\mathbb{Q}$. The
saturated torus critical locus is the disjoint union of the four points
\fitdisplay{P^A_\lambda=\left(
\lambda^3-7\lambda^2+11\lambda-1,
-\lambda^2+4\lambda-4,
-6\lambda^3+42\lambda^2-95\lambda+67
\right),
\qquad
\Phi_A(\lambda)=0}
and the four points
\fitdisplay{P^B_\mu=\left(
\frac{\mu^3+\mu^2+\mu+1}{3},
\frac{2\mu^3-\mu^2+2\mu-4}{3},
\frac{4\mu^3-2\mu^2+7\mu-11}{3}
\right),
\qquad
\Phi_B(\mu)=0}
The displayed coordinate functions are coprime to the corresponding quartic,
so all eight points lie in the torus.

The critical values on the first branch are
\[
W_{MM_{3-12}}(P^A_\lambda)=3\lambda^3-15\lambda^2+19\lambda-4
\quad\text{modulo }\Phi_A(\lambda),
\]
or, equivalently, the roots of
\[
w^4+9w^3+22w^2-24w-91=0.
\]
On the second branch one has
\[
W_{MM_{3-12}}(P^B_\mu)=\frac{7\mu^3+\mu^2+19\mu-14}{3}
\quad\text{modulo }\Phi_B(\mu),
\]
or, equivalently, the roots of
\[
w^4-3w^3-50w^2-156w-171=0.
\]
The logarithmic Hessian is nondegenerate at every critical point. More
precisely,
\[
\det(\theta_i\theta_jW)(P^A_\lambda)
=-2(3\lambda^3-22\lambda^2+39\lambda-3),
\]
which is coprime to $\Phi_A$, and
\[
\det(\theta_i\theta_jW)(P^B_\mu)
=\frac{2}{3}(19\mu^3+4\mu^2+43\mu-11),
\]
which is coprime to $\Phi_B$. Therefore all eight critical points are isolated
Morse singularities of type $A_1$. Consequently,
\[
\overline{MF((\mathbb{C}^*)^3,W_{MM_{3-12}})}
\cong \bigoplus_{\lambda:\,\Phi_A(\lambda)=0}\MorseCat{3} \oplus \bigoplus_{\mu:\,\Phi_B(\mu)=0}\MorseCat{3}.
\]
\end{example}

\begin{example}
Consider the Mori--Mukai model $MM_{3-13}$ on the algebraic torus
$(\mathbb{C}^*)^3$, with Laurent potential
\[
W_{MM_{3-13}}=x+y+z+\frac{2}{yz}+\frac{yz}{x}
+\frac{3}{x}+\frac{3}{xyz}+\frac{1}{xy^2z^2}.
\]
The last four terms contain the binomial cube
\[
\frac{yz}{x}+\frac{3}{x}+\frac{3}{xyz}+\frac{1}{xy^2z^2}
=\frac{yz}{x}\left(1+\frac{1}{yz}\right)^3.
\]
Thus, with
\[
U=1+\frac{1}{yz},
\]
one may write
\[
W_{MM_{3-13}}=x+y+z+\frac{2}{yz}+\frac{yz}{x}U^3.
\]
The logarithmic $y$- and $z$-derivatives satisfy
\[
y\frac{\partial W}{\partial y}-z\frac{\partial W}{\partial z}=y-z.
\]
Hence every torus critical point lies on the diagonal $y=z$. After imposing
$z=y$, the remaining logarithmic critical equations are
\begin{align*}
0={}&x^2y^4-(y^2+1)^3, \\
0={}&xy^2(y^3-2)+(y^2-2)(y^2+1)^2.
\end{align*}
The first equation gives
\[
x^2=\frac{(y^2+1)^3}{y^4},
\]
whereas the second gives
\[
x=-\frac{(y^2-2)(y^2+1)^2}{y^2(y^3-2)}.
\]
Combining these two expressions and clearing the torus denominators gives
\[
y^3(3y-4)(y^2+1)^3=0.
\]
The factors $y$ and $y^2+1$ do not give torus critical points: $y=0$ is
excluded, and $y^2+1=0$ would force $x=0$ in the first equation. Therefore
\[
y=z=\frac{4}{3},
\qquad
x=\frac{125}{48}.
\]
Equivalently, clearing denominators in the logarithmic critical equations and
saturating by $xyz$ gives the reduced ideal
\[
I_{\mathrm{crit}}=(48x-125,\ 3y-4,\ 3z-4).
\]
Thus the torus critical locus consists of the single rational point
\[
P=\left(\frac{125}{48},\frac{4}{3},\frac{4}{3}\right).
\]
Its critical value is
\[
W_{MM_{3-13}}(P)=9.
\]
The logarithmic Hessian is nondegenerate at this point; explicitly,
\[
\det(\theta_i\theta_jW)(P)=50.
\]
Therefore $P$ is an isolated Morse singularity of type $A_1$. Consequently,
\[
\overline{MF((\mathbb{C}^*)^3,W_{MM_{3-13}})}
\cong \MorseCat{3}.
\]
\end{example}

\begin{example}
Consider the Mori--Mukai model $MM_{3-14}$ on the algebraic torus
$(\mathbb{C}^*)^3$, with Laurent potential
\[
W_{MM_{3-14}}=\frac{x^2}{yz}+x+\frac{3x}{yz}+y+z
+\frac{3}{yz}+\frac{1}{x}+\frac{1}{xyz}.
\]
The terms with denominator $yz$ form the binomial cube
\[
\frac{x^2}{yz}+\frac{3x}{yz}+\frac{3}{yz}+\frac{1}{xyz}
=\frac{(x+1)^3}{xyz}.
\]
Thus
\[
W_{MM_{3-14}}=x+y+z+\frac{1}{x}+\frac{(x+1)^3}{xyz}.
\]
The logarithmic $y$- and $z$-derivatives satisfy
\[
y\frac{\partial W}{\partial y}-z\frac{\partial W}{\partial z}=y-z.
\]
Hence every torus critical point lies on the diagonal $y=z$. On this diagonal
the $y$-equation gives
\[
y^3=\frac{(x+1)^3}{x},
\]
whereas the $x$-equation factors as
\[
(x+1)\bigl(y^2(x-1)+(x+1)(2x-1)\bigr)=0.
\]
The factor $x+1$ would force $y=0$ in the preceding equation, and is excluded
from the torus. Therefore
\[
y^2=\frac{(x+1)(1-2x)}{x-1}.
\]
Combining the expressions for $y^3$ and $y^2$ gives the rational coordinate
\[
y=z=Y(x):=\frac{(x+1)^2(x-1)}{x(1-2x)}.
\]
Substitution gives the characteristic polynomial
\[
\Phi(x)=x^6+8x^5-15x^4+6x^3+2x^2-1.
\]
The polynomial $\Phi$ is irreducible and square-free over $\mathbb{Q}$, and
is coprime to $x(x+1)(x-1)(1-2x)$. Hence the torus critical locus consists of
the six points
\[
P_\lambda=(\lambda,Y(\lambda),Y(\lambda)),
\qquad
\Phi(\lambda)=0.
\]
Equivalently, clearing denominators in the logarithmic critical equations and
saturating by $xyz$ gives a lexicographic basis of the form
\[
\begin{aligned}
&\Phi(x),\\
&3y-(5x^5+44x^4-41x^3-13x^2+11x+7),\\
&3z-(5x^5+44x^4-41x^3-13x^2+11x+7),
\end{aligned}
\]
in the quotient by $\Phi$.

The critical values are the roots of
\[
w^6-4w^5-2w^4-182w^3-443w^2+374w+4549=0.
\]
Equivalently, modulo $\Phi(\lambda)$ one has
\[
W_{MM_{3-14}}(P_\lambda)=
6\lambda^5+52\lambda^4-56\lambda^3-7\lambda^2+14\lambda+7.
\]
The logarithmic Hessian is nondegenerate at every critical point. More
precisely,
\[
\det(\theta_i\theta_jW)(P_\lambda)
=2(5\lambda^5+46\lambda^4-21\lambda^3-7\lambda^2+10\lambda+9),
\]
and the displayed numerator is coprime to $\Phi$. Therefore all six critical
points are isolated Morse singularities of type $A_1$. Consequently,
\[
\overline{MF((\mathbb{C}^*)^3,W_{MM_{3-14}})}
\cong \bigoplus_{\lambda:\,\Phi(\lambda)=0}\MorseCat{3}.
\]
\end{example}

\begin{example}
Consider the Mori--Mukai model $MM_{3-15}$ on the algebraic torus
$(\mathbb{C}^*)^3$, with Laurent potential
\[
W_{MM_{3-15}}=x+\frac{x}{yz}+y+z+\frac{1}{y}
+\frac{y}{x}+\frac{2}{x}+\frac{1}{xy}.
\]
The logarithmic critical equations, after clearing denominators, are
\begin{align*}
0={}&x^2yz+x^2-y^2z-2yz-z, \\
0={}&-x^2+xy^2z-xz+y^2z-z, \\
0={}&-x+yz^2.
\end{align*}
The last equation gives $x=yz^2$. Substitution into the first two equations
gives the two equations
\begin{align*}
0={}&y^3z^4+y^2z^3-y^2-2y-1, \\
0={}&y^3z^2-y^2z^3+y^2-yz^2-1.
\end{align*}
The apparent diagonal specialization is not a component of the full saturated
system. Clearing denominators in the logarithmic critical equations and
saturating by $xyz$ gives a zero-dimensional ideal with lexicographic basis
\[
\Phi(y),
\qquad
82x-(21y^6-14y^5+9y^4+3y^3+106y^2+71y-284),
\]
\[
164z+55y^6-64y^5+6y^4+2y^3+221y^2+102y-818,
\]
where
\[
\Phi(y)=y^7-y^6+5y^3+3y^2-14y-2.
\]
The polynomial $\Phi$ is irreducible and square-free over $\mathbb{Q}$. Hence
the torus critical locus consists of the seven points
\[
P_\lambda=(X(\lambda),\lambda,Z(\lambda)),
\qquad
\Phi(\lambda)=0,
\]
where
\begin{align*}
X(y)&=\frac{21y^6-14y^5+9y^4+3y^3+106y^2+71y-284}{82}, \\
Z(y)&=\frac{-55y^6+64y^5-6y^4-2y^3-221y^2-102y+818}{164}.
\end{align*}
The displayed coordinate functions and $y$ are coprime to $\Phi$, so all seven
points lie in the torus.

The critical values are the roots of
\[
w^7+4w^6-43w^5-362w^4-1132w^3-1644w^2+196w+1268=0.
\]
Equivalently, modulo $\Phi(\lambda)$ one has
\[
W_{MM_{3-15}}(P_\lambda)=
\frac{\lambda^6+54\lambda^5+18\lambda^4+6\lambda^3
+171\lambda^2+388\lambda+170}{164}.
\]
The logarithmic Hessian is nondegenerate at every critical point. More
precisely,
\[
\det(\theta_i\theta_jW)(P_\lambda)
=\frac{84\lambda^6-15\lambda^5-5\lambda^4-29\lambda^3
+465\lambda^2+612\lambda-562}{41},
\]
and the displayed numerator is coprime to $\Phi$. Therefore all seven critical
points are isolated Morse singularities of type $A_1$. Consequently,
\[
\overline{MF((\mathbb{C}^*)^3,W_{MM_{3-15}})}
\cong \bigoplus_{\lambda:\,\Phi(\lambda)=0}\MorseCat{3}.
\]
\end{example}

\begin{example}
Consider the Mori--Mukai model $MM_{3-16}$ on the algebraic torus
$(\mathbb{C}^*)^3$, with Laurent potential
\[
W_{MM_{3-16}}=x+\frac{x}{y}+y+\frac{y}{z}+z+\frac{1}{y}
+\frac{y}{xz}+\frac{1}{x}.
\]
The logarithmic critical equations, after clearing denominators, are
\begin{align*}
0={}&x^2yz+x^2z-y^2-yz, \\
0={}&-x^2z+xy^2z+xy^2-xz+y^2, \\
0={}&-xy+xz^2-y.
\end{align*}
The last equation gives
\[
z^2=\frac{y(x+1)}{x}.
\]
Using this relation in the $y$-equation gives
\[
y+z=\frac{x+1}{y}.
\]
Substituting this identity in the $x$-equation then gives
\[
z=xy(y+1).
\]
The preceding relation and $y+z=(x+1)/y$ determine
\[
x=\frac{1-y^2}{y^3+y^2-1}.
\]
Substituting this expression in $z^2=y(x+1)/x$ gives the characteristic
equation
\[
\Phi(y):=2y^8+4y^7-y^6-8y^5-2y^4+6y^3+3y^2-2y-1=0.
\]
Equivalently, the identity before expansion is
\[
(1-y^2)^3(y+1)^2=y^2(y^3+y^2-1)^2.
\]
The polynomial $\Phi$ is irreducible and square-free over $\mathbb{Q}$, and
is coprime to $y(y-1)(y+1)(y^3+y^2-1)$. Hence all eight roots give torus
points. The critical locus consists of the eight points
\[
P_\lambda=(X(\lambda),\lambda,Z(\lambda)),
\qquad
\Phi(\lambda)=0,
\]
where
\[
X(y)=\frac{1-y^2}{y^3+y^2-1},
\qquad
Z(y)=\frac{y(1-y^2)(y+1)}{y^3+y^2-1}.
\]
Equivalently, modulo $\Phi$, these coordinates are represented by
\[
X(y)=2y^6+4y^5-y^4-6y^3+3y,
\]
and
\[
Z(y)=1+2y^7+4y^6+y^5-4y^4-3y^3+2y.
\]

The critical values are the roots of
\[
w^8+6w^7-9w^6-330w^5-2081w^4-7470w^3
-16278w^2-20580w-12391=0.
\]
Equivalently, modulo $\Phi(\lambda)$ one has
\[
W_{MM_{3-16}}(P_\lambda)=
8\lambda^7+18\lambda^6+8\lambda^5-17\lambda^4
-16\lambda^3+2\lambda^2+12\lambda+2.
\]
The logarithmic Hessian is nondegenerate at every critical point. More
precisely,
\[
\det(\theta_i\theta_jW)(P_\lambda)
=2(14\lambda^7+30\lambda^6+5\lambda^5-37\lambda^4
-12\lambda^3+18\lambda^2+12\lambda+1),
\]
and the displayed numerator is coprime to $\Phi$. Therefore all eight critical
points are isolated Morse singularities of type $A_1$. Consequently,
\[
\overline{MF((\mathbb{C}^*)^3,W_{MM_{3-16}})}
\cong \bigoplus_{\lambda:\,\Phi(\lambda)=0}\MorseCat{3}.
\]
\end{example}

\begin{example}
Consider the Mori--Mukai model $MM_{3-17}$ on the algebraic torus
$(\mathbb{C}^*)^3$, with Laurent potential
\[
W_{MM_{3-17}}=x+y+z+\frac{1}{y}+\frac{2}{yz}
+\frac{1}{x}+\frac{2}{xyz}+\frac{1}{xy^2z^2}.
\]
The last three terms form the square
\[
\frac{1}{x}+\frac{2}{xyz}+\frac{1}{xy^2z^2}
=\frac{1}{x}\left(1+\frac{1}{yz}\right)^2.
\]
Set
\[
U=1+\frac{1}{yz}.
\]
Then
\[
W_{MM_{3-17}}=x+y+z+\frac{1}{y}+\frac{2}{yz}+\frac{U^2}{x}.
\]
The $x$-derivative gives
\[
0=1-\frac{U^2}{x^2},
\qquad\text{hence}\qquad
x=\pm U.
\]
The negative branch is empty. Indeed, using
$\partial U/\partial z=-1/(yz^2)$, the $z$-derivative is
\[
\frac{\partial W}{\partial z}
=1-\frac{2}{yz^2}\left(1+\frac{U}{x}\right),
\]
which becomes $1=0$ after substituting $x=-U$. Therefore every torus critical
point lies on the positive branch $x=U$.

On this branch the $z$-equation becomes
\[
y=\frac{4}{z^2},
\]
and the $y$-equation becomes
\[
y^2=1+\frac{4}{z}.
\]
Combining these two equations gives
\[
\Phi(z):=z^4+4z^3-16=0.
\]
The polynomial $\Phi$ is irreducible and square-free over $\mathbb{Q}$, and is
coprime to $z(z+4)$. Thus the torus critical locus consists of the four
points
\[
P_\lambda=\left(1+\frac{\lambda}{4},\frac{4}{\lambda^2},\lambda\right),
\qquad
\Phi(\lambda)=0.
\]
Equivalently, clearing denominators in the logarithmic critical equations and
saturating by $xyz$ gives the reduced lexicographic basis
\[
4x-z-4,
\qquad
4y-z^2-4z,
\qquad
z^4+4z^3-16.
\]

The critical values are the roots of
\[
w^4-4w^3-8w^2-96w-176=0.
\]
Equivalently, modulo $\Phi(\lambda)$ one has
\[
W_{MM_{3-17}}(P_\lambda)=\frac{\lambda^2+6\lambda+4}{2}.
\]
The logarithmic Hessian is nondegenerate at every critical point. More
precisely,
\[
\det(\theta_i\theta_jW)(P_\lambda)
=\frac{3\lambda^3+12\lambda^2+16}{2},
\]
and the displayed numerator is coprime to $\Phi$. Therefore all four critical
points are isolated Morse singularities of type $A_1$. Consequently,
\[
\overline{MF((\mathbb{C}^*)^3,W_{MM_{3-17}})}
\cong \bigoplus_{\lambda:\,\Phi(\lambda)=0}\MorseCat{3}.
\]
\end{example}

\begin{example}
Consider the Mori--Mukai model $MM_{3-18}$ on the algebraic torus
$(\mathbb{C}^*)^3$, with Laurent potential
\[
W_{MM_{3-18}}=x+\frac{x}{y}+y+z+\frac{z}{x}+\frac{2}{x}
+\frac{1}{xz}.
\]
The logarithmic critical equations, after clearing denominators, are
\begin{align*}
0={}&x^2yz+x^2z-yz^2-2yz-y, \\
0={}&-x+y^2, \\
0={}&xz^2+z^2-1.
\end{align*}
The second equation gives
\[
x=y^2.
\]
The third equation then gives
\[
z^2=\frac{1}{y^2+1}.
\]
Finally, the $x$-derivative gives
\[
x^2+\frac{x^2}{y}=z+2+\frac{1}{z}.
\]
Using $x=y^2$ and $1/z=z(y^2+1)$, this becomes
\[
y^4+y^3-2=z(y^2+2),
\]
and therefore
\[
z=\frac{y^4+y^3-2}{y^2+2}.
\]
Substitution into $z^2=1/(y^2+1)$ gives
\[
(y^4+y^3-2)^2(y^2+1)=(y^2+2)^2.
\]
After removing the forbidden torus factor $y^3$, one obtains the characteristic
polynomial
\[
\Phi(y):=y^7+2y^6+2y^5+2y^4-3y^3-4y^2-5y-4=0.
\]
The polynomial $\Phi$ is irreducible and square-free over $\mathbb{Q}$, and is
coprime to $y(y^2+1)(y^2+2)(y^4+y^3-2)$. Hence the torus critical locus consists
of the seven points
\[
P_\lambda=\left(\lambda^2,\lambda,
\frac{\lambda^4+\lambda^3-2}{\lambda^2+2}\right),
\qquad
\Phi(\lambda)=0.
\]
Equivalently, clearing denominators in the logarithmic critical equations and
saturating by $xyz$ gives the lexicographic basis
\[
x-y^2,
\qquad
3z+y^6-4y^4-2y^3-2y^2+3y+7,
\qquad
\Phi(y).
\]

The critical values are the roots of
\[
w^7+5w^6-22w^5-336w^4-1225w^3-1141w^2+2811w+4673=0.
\]
Equivalently, modulo $\Phi(\lambda)$ one has
\[
W_{MM_{3-18}}(P_\lambda)=
\frac{-\lambda^6+4\lambda^4+2\lambda^3+8\lambda^2+6\lambda-7}{3}.
\]
The logarithmic Hessian is nondegenerate at every critical point. More
precisely,
\[
\det(\theta_i\theta_jW)(P_\lambda)
=-\frac{2}{3}(5\lambda^6+15\lambda^5+13\lambda^4-4\lambda^3
-22\lambda^2-45\lambda-31),
\]
and the displayed numerator is coprime to $\Phi$. Therefore all seven critical
points are isolated Morse singularities of type $A_1$. Consequently,
\[
\overline{MF((\mathbb{C}^*)^3,W_{MM_{3-18}})}
\cong \bigoplus_{\lambda:\,\Phi(\lambda)=0}\MorseCat{3}.
\]
\end{example}

\begin{example}
Consider the Mori--Mukai model $MM_{3-19}$ on the algebraic torus
$(\mathbb{C}^*)^3$, with Laurent potential
\[
W_{MM_{3-19}}=x+\frac{2x}{yz}+\frac{x}{y^2z^2}+y+z
+\frac{2}{yz}+\frac{1}{x}.
\]
The first three terms form the square
\[
x+\frac{2x}{yz}+\frac{x}{y^2z^2}
=x\left(1+\frac{1}{yz}\right)^2.
\]
Set
\[
U=1+\frac{1}{yz}.
\]
Then
\[
W_{MM_{3-19}}=xU^2+y+z+\frac{2}{yz}+\frac{1}{x}.
\]
The $y$- and $z$-derivatives are
\[
\frac{\partial W}{\partial y}
=-\frac{2xU}{y^2z}+1-\frac{2}{y^2z},
\qquad
\frac{\partial W}{\partial z}
=-\frac{2xU}{yz^2}+1-\frac{2}{yz^2}.
\]
After multiplying by $y^2z$ and $yz^2$, respectively, they give
\[
y^2z=2xU+2,
\qquad
yz^2=2xU+2.
\]
Hence every torus critical point satisfies
\[
y=z.
\]
On this diagonal one has
\[
U=1+\frac{1}{y^2}=\frac{y^2+1}{y^2},
\]
and the transverse equation becomes
\[
y^3-2=2xU.
\]
The $x$-derivative gives
\[
U^2-\frac{1}{x^2}=0,
\qquad\text{hence}\qquad
x=\pm\frac{1}{U}=\pm\frac{y^2}{y^2+1}.
\]
The negative branch is empty: if $xU=-1$, then the equation above gives
$y^3=0$, impossible in the torus. On the positive branch, $xU=1$, and hence
\[
\Phi(y):=y^3-4=0.
\]
The polynomial $\Phi$ is irreducible and square-free over $\mathbb{Q}$, and is
coprime to $y(y^2+1)$. Thus the torus critical locus consists of the three
points
\[
P_\lambda=\left(\frac{\lambda^2}{\lambda^2+1},\lambda,\lambda\right),
\qquad
\Phi(\lambda)=0.
\]
Equivalently, clearing denominators in the logarithmic critical equations and
saturating by $xyz$ gives the lexicographic basis
\[
17x-y^2+4y-16,
\qquad
z-y,
\qquad
y^3-4.
\]

The critical values are the roots of
\[
w^3-6w^2+12w-116=0.
\]
Equivalently, modulo $\Phi(\lambda)$ one has
\[
W_{MM_{3-19}}(P_\lambda)=3\lambda+2.
\]
The logarithmic Hessian is nondegenerate at every critical point. More
precisely,
\[
\det(\theta_i\theta_jW)(P_\lambda)=6(\lambda^2+1),
\]
and $\lambda^2+1$ is coprime to $\Phi$. Therefore all three critical points are
isolated Morse singularities of type $A_1$. Consequently,
\[
\overline{MF((\mathbb{C}^*)^3,W_{MM_{3-19}})}
\cong \bigoplus_{\lambda:\,\Phi(\lambda)=0}\MorseCat{3}.
\]
\end{example}

\begin{example}
Consider the Mori--Mukai model $MM_{3-20}$ on the algebraic torus
$(\mathbb{C}^*)^3$, with Laurent potential
\[
W_{MM_{3-20}}=x+y+z+\frac{1}{yz}+\frac{y}{x}
+\frac{2}{x}+\frac{1}{xy}.
\]
The last three terms have the common factor $1/x$ and satisfy
\[
\frac{y}{x}+\frac{2}{x}+\frac{1}{xy}
=\frac{1}{x}\left(y+2+\frac{1}{y}\right)
=\frac{1}{x}\frac{(y+1)^2}{y}.
\]
Set
\[
U=\frac{(y+1)^2}{y}.
\]
Then
\[
W_{MM_{3-20}}=x+y+z+\frac{1}{yz}+\frac{U}{x}.
\]
The logarithmic critical equations, after clearing denominators, are
\begin{align*}
0={}&x^2y-y^2-2y-1, \\
0={}&xy^2z-x+y^2z-z, \\
0={}&yz^2-1.
\end{align*}
The last equation gives
\[
y=\frac{1}{z^2}.
\]
The $x$-equation gives
\[
x^2=U.
\]
Substituting $y=1/z^2$ in the $y$-equation gives
\[
1-z^3+\frac{1}{x}(1-z^4)=0,
\]
or equivalently
\[
x(1-z^3)=z^4-1.
\]
This relation has two branches. If $z=1$, then $y=1$, $U=4$, and hence
\[
(x,y,z)=(2,1,1)\quad\text{or}\quad(-2,1,1).
\]
If $z\ne1$, cancellation of the common factor gives
\[
x=-\frac{(z+1)(z^2+1)}{z^2+z+1}.
\]
Together with $x^2=U=(z^2+1)^2/z^2$, this implies
\[
z^2(z+1)^2=(z^2+z+1)^2.
\]
Since the denominators and $z^2+1$ are nonzero on this branch, this reduces to
\[
\Psi(z):=2z^2+2z+1=0.
\]
The polynomial $\Psi$ is irreducible and square-free over $\mathbb{Q}$, and it
is coprime to $z(z-1)(z+1)(z^2+1)(z^2+z+1)$. Thus the remaining critical
points are
\[
Q_\mu=(-\mu-2,4\mu+2,\mu),
\qquad
\Psi(\mu)=0,
\]
where $4\mu+2=1/\mu^2$ modulo $\Psi(\mu)$. Equivalently, clearing denominators
in the logarithmic critical equations and saturating by $xyz$ gives the
lexicographic basis
\[
x^2+z^2-2z-3,
\qquad
(z-1)(x+z+2),
\]
\[
y+2z^2-2z-1,
\qquad
(z-1)(2z^2+2z+1).
\]

The critical values at the two rational points are
\[
W_{MM_{3-20}}(2,1,1)=7,
\qquad
W_{MM_{3-20}}(-2,1,1)=-1.
\]
On the quadratic branch, modulo $\Psi(\mu)$ one has
\[
W_{MM_{3-20}}(Q_\mu)=4\mu-2,
\]
and these two values are the roots of
\[
w^2+8w+20=0.
\]
Equivalently, all critical values are the roots of
\[
(w-7)(w+1)(w^2+8w+20)=0.
\]
The logarithmic Hessian is nondegenerate at every critical point. More
precisely,
\[
\det(\theta_i\theta_jW)(2,1,1)=20,
\qquad
\det(\theta_i\theta_jW)(-2,1,1)=-4,
\]
and on the quadratic branch
\[
\det(\theta_i\theta_jW)(Q_\mu)=10(2\mu+1),
\]
which is coprime to $\Psi$. Therefore all four critical points are isolated
Morse singularities of type $A_1$. Consequently,
\[
\overline{MF((\mathbb{C}^*)^3,W_{MM_{3-20}})}
\cong \left(\bigoplus_{\epsilon=\pm1}\MorseCat{3}\right) \oplus \bigoplus_{\mu:\,\Psi(\mu)=0}\MorseCat{3}.
\]
\end{example}

\begin{example}
Consider the Mori--Mukai model $MM_{3-21}$ on the algebraic torus
$(\mathbb{C}^*)^3$, with Laurent potential
\[
W_{MM_{3-21}}=x+y+z+\frac{1}{y}+\frac{y}{x}+\frac{2}{x}
+\frac{1}{xyz}+\frac{1}{x^2z}.
\]
It is convenient to pass to inverse torus coordinates
\[
u=\frac{1}{x},
\qquad
v=\frac{1}{y},
\qquad
\omega=\frac{1}{z}.
\]
In these coordinates the potential becomes
\[
W_{MM_{3-21}}=\frac{1}{u}+\frac{1}{v}+\frac{1}{\omega}
+v+\frac{u}{v}+2u+uv\omega+u^2\omega.
\]
The critical equations are
\begin{align*}
0={}&-\frac{1}{u^2}+\frac{1}{v}+2+v\omega+2u\omega, \\
0={}&-\frac{1}{v^2}+1-\frac{u}{v^2}+u\omega, \\
0={}&-\frac{1}{\omega^2}+uv+u^2.
\end{align*}
The last two equations give
\[
\omega^2=\frac{1}{u(u+v)},
\qquad
v^2(1+u\omega)=1+u.
\]
Set
\[
T=u\omega.
\]
The equation in $\omega$ gives
\[
T^2+Tv\omega=1,
\qquad\text{hence}\qquad
v=\frac{1-T^2}{T\omega}.
\]
Substituting this expression in the remaining two equations gives the rational
parametrization
\[
\omega=\frac{(1-T)(1+T)^3(1-T-T^2)}{T^3},
\]
and the characteristic equation
\[
\Phi(T):=T^8+4T^7+4T^6-6T^5-8T^4+2T^3+5T^2-1=0.
\]
The polynomial $\Phi$ is irreducible and square-free over $\mathbb{Q}$, and is
coprime to
\[
T(1-T)(1+T)(1-T-T^2).
\]
Thus all eight roots give torus points. Returning to the original coordinates,
the critical locus consists of the eight points
\fitdisplay{P_\tau=\left(
\frac{(1-\tau)(1+\tau)^3(1-\tau-\tau^2)}{\tau^4},
\frac{(1+\tau)^2(1-\tau-\tau^2)}{\tau^2},
\frac{\tau^3}{(1-\tau)(1+\tau)^3(1-\tau-\tau^2)}
\right)}
where $\Phi(\tau)=0$. Equivalently, adjoining $T=u\omega$ to the logarithmic
critical equations in the inverse coordinates and saturating by $uv\omega T$
gives the lexicographic basis
\[
u+T^7+3T^6-10T^4-2T^3+10T^2+3T-4,
\]
\[
v+2T^7+5T^6-3T^5-21T^4+2T^3+18T^2+2T-6,
\]
\[
\omega-2T^7-9T^6-13T^5+4T^4+17T^3+7T^2-5T-3,
\qquad
\Phi(T).
\]

The critical values are the roots of
\[
\xi^8+6\xi^7-41\xi^6-324\xi^5-511\xi^4+20\xi^3
+2655\xi^2+5478\xi+3509=0.
\]
Equivalently, modulo $\Phi(\tau)$ one has
\[
W_{MM_{3-21}}(P_\tau)=
(\tau+1)(8\tau^6+31\tau^5+31\tau^4-40\tau^3
-46\tau^2+7\tau+17).
\]
The logarithmic Hessian is nondegenerate at every critical point. More
precisely,
\[
\det(\theta_i\theta_jW)(P_\tau)
=2(19\tau^7+81\tau^6+106\tau^5-52\tau^4-129\tau^3
-32\tau^2+42\tau+15),
\]
and the displayed numerator is coprime to $\Phi$. Therefore all eight critical
points are isolated Morse singularities of type $A_1$. Consequently,
\[
\overline{MF((\mathbb{C}^*)^3,W_{MM_{3-21}})}
\cong \bigoplus_{\tau:\,\Phi(\tau)=0}\MorseCat{3}.
\]
\end{example}

\begin{example}
Consider the Mori--Mukai model $MM_{3-22}$ on the algebraic torus
$(\mathbb{C}^*)^3$, with Laurent potential
\[
W_{MM_{3-22}}=\frac{x^2}{yz}+x+\frac{2x}{yz}+y+z
+\frac{1}{yz}+\frac{1}{x}.
\]
The terms with denominator $yz$ form the square
\[
\frac{x^2}{yz}+\frac{2x}{yz}+\frac{1}{yz}
=\frac{(x+1)^2}{yz},
\]
so that
\[
W_{MM_{3-22}}=x+y+z+\frac{1}{x}+\frac{(x+1)^2}{yz}.
\]
The $y$- and $z$-derivatives give
\[
y^2z=(x+1)^2,
\qquad
yz^2=(x+1)^2.
\]
Hence every torus critical point satisfies
\[
y=z,
\]
and therefore
\[
y^3=(x+1)^2.
\]
On the diagonal, the $x$-derivative becomes
\[
1-\frac{1}{x^2}+\frac{2(x+1)}{y^2}=0.
\]
The branch $x=-1$ would force $y^3=0$, which is impossible in the torus. Thus
one may divide by $x+1$, obtaining
\[
\frac{x-1}{x^2}+\frac{2}{y^2}=0,
\qquad\text{or equivalently}\qquad
y^2=\frac{2x^2}{1-x}.
\]
Combining $y^3=(x+1)^2$ with $y^2=2x^2/(1-x)$ gives
\[
8x^6=(x+1)^4(1-x)^3.
\]
Equivalently, the characteristic equation is
\[
\Phi(x):=x^7+9x^6-3x^5-3x^4+3x^3+3x^2-x-1=0.
\]
The polynomial $\Phi$ is irreducible and square-free over $\mathbb{Q}$, and is
coprime to $x(x-1)(x+1)$. Moreover the transverse coordinate is recovered
rationally as
\[
y=z=\frac{(x+1)^2(1-x)}{2x^2}.
\]
Thus the torus critical locus consists of the seven points
\[
P_\lambda=\left(\lambda,
\frac{(\lambda+1)^2(1-\lambda)}{2\lambda^2},
\frac{(\lambda+1)^2(1-\lambda)}{2\lambda^2}\right),
\qquad
\Phi(\lambda)=0.
\]
Equivalently, clearing denominators in the logarithmic critical equations and
saturating by $xyz$ gives the lexicographic basis
\[
x^5+9x^4-3x^3-3x^2+2x-2y+2,
\]
\[
x^5+9x^4-3x^3-3x^2+2x-2z+2,
\qquad
\Phi(x).
\]

The critical values are the roots of
\[
w^7-2w^6-16w^5-24w^4-424w^3-180w^2+4624w+2764=0.
\]
Equivalently, modulo $\Phi(\lambda)$ one has
\[
W_{MM_{3-22}}(P_\lambda)=
\frac{2\lambda^6+21\lambda^5+21\lambda^4-15\lambda^3
-3\lambda^2+14\lambda+4}{2}.
\]
The logarithmic Hessian is nondegenerate at every critical point. More
precisely,
\[
\det(\theta_i\theta_jW)(P_\lambda)
=\frac{3\lambda^6+30\lambda^5+21\lambda^4+12\lambda^3
+17\lambda^2+30\lambda+3}{2},
\]
and the displayed numerator is coprime to $\Phi$. Therefore all seven critical
points are isolated Morse singularities of type $A_1$. Consequently,
\[
\overline{MF((\mathbb{C}^*)^3,W_{MM_{3-22}})}
\cong \bigoplus_{\lambda:\,\Phi(\lambda)=0}\MorseCat{3}.
\]
\end{example}

\begin{example}
Consider the Mori--Mukai model $MM_{3-23}$ on the algebraic torus
$(\mathbb{C}^*)^3$, with Laurent potential
\[
W_{MM_{3-23}}=xz+x+\frac{x}{y}+y+z+\frac{1}{yz}+\frac{1}{x}.
\]
Introduce the master variables
\[
U=xz,
\qquad
V=yz.
\]
The $y$-derivative gives
\[
y^2z=xz+1,
\qquad\text{hence}\qquad
z=\frac{V^2}{U+1}.
\]
The $z$-derivative gives
\[
yz^2(x+1)=1,
\qquad\text{hence}\qquad
z=\frac{1-UV}{V}.
\]
Equating these two expressions for $z$ gives the first master equation
\[
V^3+V(U^2+U)-(U+1)=0.
\]
The $x$-derivative is
\[
x^2(yz+y+1)=y.
\]
Using $x=U/z$, $y=V/z$, and $z=(1-UV)/V$, this gives the second master equation
\[
U^3V^2+U^3V-U^2V-U^2-2UV+1=0.
\]
Eliminating $U$ and saturating by the torus factor $V$ gives
\[
\Psi(V):=V^8+2V^7+V^6+2V^5-4V^3-2V+1=0.
\]
Moreover, modulo $\Psi$ one has the rational expression
\[
U=\frac{N(V)}{11},
\]
where
\[
N(V)=V^7+4V^6+9V^5+9V^4+7V^3-V^2-13V-6.
\]
Set
\[
D(V)=VN(V)-11.
\]
Then the original coordinates are recovered rationally as
\[
z=-\frac{D(V)}{11V},
\qquad
x=-\frac{VN(V)}{D(V)},
\qquad
y=-\frac{11V^2}{D(V)}.
\]
The polynomial $\Psi$ is irreducible and square-free over $\mathbb{Q}$, and is
coprime to $V N(V)D(V)$. Thus the torus critical locus consists of the eight
points
\[
P_\lambda=\left(-\frac{\lambda N(\lambda)}{D(\lambda)},
-\frac{11\lambda^2}{D(\lambda)},
-\frac{D(\lambda)}{11\lambda}\right),
\qquad
\Psi(\lambda)=0.
\]
Equivalently, after adjoining $U=xz$ and $V=yz$, clearing denominators, and
saturating by the torus variables, the reduced elimination basis in the master
variables is
\[
11U-N(V),
\qquad
\Psi(V).
\]

The critical values are the roots of
\[
w^8+6w^7+9w^6-160w^5-1174w^4-4666w^3
-11104w^2-17660w-11075=0.
\]
Equivalently, modulo $\Psi(\lambda)$ one has
\[
W_{MM_{3-23}}(P_\lambda)=
(\lambda+1)(\lambda^2+
\lambda-1)(\lambda^3+2\lambda+1).
\]
The logarithmic Hessian is nondegenerate at every critical point. More
precisely,
\[
\det(\theta_i\theta_jW)(P_\lambda)
=\frac{2}{11}(15\lambda^7+27\lambda^6+14\lambda^5+58\lambda^4
+17\lambda^3-48\lambda^2+58\lambda-35),
\]
and the displayed numerator is coprime to $\Psi$. Therefore all eight critical
points are isolated Morse singularities of type $A_1$. Consequently,
\[
\overline{MF((\mathbb{C}^*)^3,W_{MM_{3-23}})}
\cong \bigoplus_{\lambda:\,\Psi(\lambda)=0}\MorseCat{3}.
\]
\end{example}

\begin{example}
Consider the Mori--Mukai model $MM_{3-24}$ on the algebraic torus
$(\mathbb{C}^*)^3$, with Laurent potential
\[
W_{MM_{3-24}}=x+y+z+\frac{1}{y}+\frac{y}{x}+\frac{1}{x}
+\frac{1}{xyz}.
\]
The last three terms have the common factor $1/x$:
\[
\frac{y}{x}+\frac{1}{x}+\frac{1}{xyz}
=\frac{1}{x}\left(y+1+\frac{1}{yz}\right).
\]
Set
\[
U=y+1+\frac{1}{yz}.
\]
Then
\[
W_{MM_{3-24}}=x+y+z+\frac{1}{y}+\frac{U}{x}.
\]
The $z$-derivative gives
\[
x=\frac{1}{yz^2},
\qquad\text{hence}\qquad
y=\frac{1}{xz^2}.
\]
The $x$-derivative gives
\[
x^2=U.
\]
Introduce the master parameter
\[
T=xz.
\]
Substituting $y=1/(xz^2)$ into $x^2=U$ gives
\[
T^3-z=z^2T(T+1).
\]
The $y$-derivative is
\[
\frac{1}{x}\left(1-\frac{1}{y^2z}\right)+1-\frac{1}{y^2}=0.
\]
Using $y=1/(xz^2)$ and $T=xz$, this becomes
\[
T+z=z^2T^2(T+1).
\]
Multiplying the first of these two equations by $T$ and comparing with the
second gives
\[
z(T+1)=T^4-T,
\]
and therefore
\[
z=\frac{T(T^3-1)}{T+1}.
\]
Substituting this expression back into $T+z=z^2T^2(T+1)$ gives
\[
\Phi(T):=T^8-2T^5-1=0.
\]
The polynomial $\Phi$ is irreducible and square-free over $\mathbb{Q}$, and is
coprime to $T(T+1)(T^3-1)$. Hence all eight roots give torus points. The
critical locus consists of the eight points
\[
P_\tau=\left(
\frac{\tau+1}{(\tau-1)(\tau^2+\tau+1)},
\frac{\tau+1}{\tau^2(\tau-1)(\tau^2+\tau+1)},
\frac{\tau(\tau-1)(\tau^2+\tau+1)}{\tau+1}
\right),
\]
where $\Phi(\tau)=0$. Equivalently, adjoining $T=xz$ to the logarithmic
critical equations and saturating by $xyzT$ gives the lexicographic basis
\[
x-T^6+T^3+1,
\qquad
y+T^6-T^4-2T^3+T,
\]
\[
z+T^7-T^6+T^5-3T^4+2T^3-2T^2+2T-1,
\qquad
\Phi(T).
\]

The critical values are the roots of
\[
w^8+4w^7-24w^6-160w^5-494w^4-1304w^3
-720w^2+128w+1973=0.
\]
Equivalently, modulo $\Phi(\tau)$ one has
\[
W_{MM_{3-24}}(P_\tau)=\tau^6+2\tau^4-2\tau-2.
\]
The logarithmic Hessian is nondegenerate at every critical point. More
precisely,
\[
\det(\theta_i\theta_jW)(P_\tau)=2\tau^2(4\tau^3-5),
\]
and $\tau^2(4\tau^3-5)$ is coprime to $\Phi$. Therefore all eight critical
points are isolated Morse singularities of type $A_1$. Consequently,
\[
\overline{MF((\mathbb{C}^*)^3,W_{MM_{3-24}})}
\cong \bigoplus_{\tau:\,\Phi(\tau)=0}\MorseCat{3}.
\]
\end{example}

\begin{example}
Consider the Mori--Mukai model $MM_{3-25}$ on the algebraic torus
$(\mathbb{C}^*)^3$, with Laurent potential
\[
W_{MM_{3-25}}=x+\frac{x}{z}+y+z+\frac{1}{x}+\frac{1}{xy}.
\]
The $z$- and $y$-derivatives immediately give
\[
x=z^2,
\qquad
xy^2=1.
\]
Thus
\[
(yz)^2=1,
\]
and the critical locus splits into two branches.

On the positive branch $yz=1$, one has $y=1/z$ and $x=z^2$. The remaining
logarithmic $x$-equation becomes
\[
z^2+z-\frac{1}{z^2}-\frac{1}{z}=0,
\]
or equivalently
\[
q_+(z):=z^4+z^3-z-1=(z^3-1)(z+1)=0.
\]
On the negative branch $yz=-1$, one has $y=-1/z$ and $x=z^2$. The same equation
becomes
\[
z^2+z-\frac{1}{z^2}+\frac{1}{z}=0,
\]
or equivalently
\[
q_-(z):=z^4+z^3+z-1=(z^2+1)(z^2+z-1)=0.
\]
The polynomials $q_+$ and $q_-$ are square-free and coprime. Hence the torus
critical locus consists of the eight points
\[
P^+_\alpha=\left(\alpha^2,\frac{1}{\alpha},\alpha\right),
\qquad q_+(\alpha)=0,
\]
and
\[
P^-_\beta=\left(\beta^2,-\frac{1}{\beta},\beta\right),
\qquad q_-(\beta)=0.
\]
Equivalently, clearing denominators in the logarithmic critical equations and
saturating by $xyz$ gives the lexicographic basis
\[
x-z^2,
\qquad
y-z^6-2z^5-z^4+z^2+z+1,
\]
\[
(z-1)(z+1)(z^2+1)(z^2+z-1)(z^2+z+1).
\]

On the positive branch the critical values are the roots of
\[
(w-6)(w+2)(w+3)^2=0,
\]
and, modulo $q_+(\alpha)$,
\[
W_{MM_{3-25}}(P^+_\alpha)=\alpha^3+3\alpha^2+3\alpha-1.
\]
On the negative branch the critical values are the roots of
\[
(w-1)^2(w^2+4w+20)=0,
\]
and, modulo $q_-(\beta)$,
\[
W_{MM_{3-25}}(P^-_\beta)=-\beta^3+\beta^2+3\beta-1.
\]
The logarithmic Hessian is nondegenerate at every critical point. More
precisely,
\[
\det(\theta_i\theta_jW)(P^+_\alpha)
=-2(\alpha^3-3\alpha^2-3\alpha-1),
\]
and
\[
\det(\theta_i\theta_jW)(P^-_\beta)
=-2(\beta^3+5\beta^2+3\beta+1),
\]
with the displayed factors coprime to $q_+$ and $q_-$, respectively. Therefore
all eight critical points are isolated Morse singularities of type $A_1$.
Consequently,
\[
\overline{MF((\mathbb{C}^*)^3,W_{MM_{3-25}})}
\cong \left(\bigoplus_{\alpha:\,q_+(\alpha)=0}\MorseCat{3}\right) \oplus \left(\bigoplus_{\beta:\,q_-(\beta)=0}\MorseCat{3}\right).
\]
\end{example}

\begin{example}
Consider the Mori--Mukai model $MM_{3-26}$ on the algebraic torus
$(\mathbb{C}^*)^3$, with Laurent potential
\[
W_{MM_{3-26}}=xy+x+y+z+\frac{1}{yz}+\frac{1}{x}.
\]
The critical equations form a rational cascade. The $z$-derivative gives
\[
yz^2=1,
\]
and hence
\[
y=\frac{1}{z^2}.
\]
The $y$-derivative then gives
\[
x+1-\frac{1}{y^2z}=0,
\qquad\text{so}\qquad
x=z^3-1.
\]
Substituting these expressions into the $x$-derivative
\[
y+1-\frac{1}{x^2}=0
\]
gives
\[
(z^3-1)^2=\frac{z^2}{z^2+1}.
\]
Equivalently,
\[
\Psi(z):=z^8+z^6-2z^5-2z^3+1=0.
\]
The polynomial $\Psi$ is irreducible and square-free over $\mathbb{Q}$, and is
coprime to $z(z^3-1)(z^2+1)$. Hence all eight roots give torus points. The
critical locus consists of the eight points
\[
P_\lambda=\left(\lambda^3-1,\frac{1}{\lambda^2},\lambda\right),
\qquad
\Psi(\lambda)=0.
\]
Equivalently, clearing denominators in the logarithmic critical equations and
saturating by $xyz$ gives the lexicographic basis
\[
x-z^3+1,
\qquad
y+z^6+z^4-2z^3-2z,
\]
\[
\Psi(z).
\]

The critical values are the roots of
\[
w^8+4w^7-4w^6-94w^5-554w^4-2176w^3
-4911w^2-7526w-407=0.
\]
Equivalently, modulo $\Psi(\lambda)$ one has
\[
W_{MM_{3-26}}(P_\lambda)=\lambda^6+\lambda^4+2\lambda-2.
\]
The logarithmic Hessian is nondegenerate at every critical point. More
precisely,
\[
\det(\theta_i\theta_jW)(P_\lambda)
=2(4\lambda^5+3\lambda^3-5\lambda^2-3),
\]
and the displayed factor is coprime to $\Psi$. Therefore all eight critical
points are isolated Morse singularities of type $A_1$. Consequently,
\[
\overline{MF((\mathbb{C}^*)^3,W_{MM_{3-26}})}
\cong \bigoplus_{\lambda:\,\Psi(\lambda)=0}\MorseCat{3}.
\]
\end{example}

\begin{example}
Consider the Mori--Mukai model $MM_{3-27}$ on the algebraic torus
$(\mathbb{C}^*)^3$, with Laurent potential
\[
W_{MM_{3-27}}=x+y+z+\frac{1}{x}+\frac{1}{y}+\frac{1}{z}.
\]
This potential is the direct sum of three one-dimensional Laurent blocks. Its
critical equations are completely decoupled:
\[
1-\frac{1}{x^2}=0,
\qquad
1-\frac{1}{y^2}=0,
\qquad
1-\frac{1}{z^2}=0.
\]
Equivalently,
\[
x^2=1,
\qquad
y^2=1,
\qquad
z^2=1.
\]
Thus the torus critical locus consists of the eight vertices
\[
P_\epsilon=(\epsilon_1,\epsilon_2,\epsilon_3),
\qquad
\epsilon_i\in\{1,-1\}.
\]
Equivalently, clearing denominators in the logarithmic critical equations and
saturating by $xyz$ gives the lexicographic basis
\[
x^2-1,
\qquad
y^2-1,
\qquad
z^2-1.
\]

The critical values are the roots of
\[
(w-6)(w-2)^3(w+2)^3(w+6)=0,
\]
or, equivalently,
\[
w^8-48w^6+480w^4-1792w^2+2304=0.
\]
They are distributed as follows:
\[
W(1,1,1)=6,
\qquad
W(-1,-1,-1)=-6,
\]
and the three vertices with exactly one negative coordinate have value $2$,
whereas the three vertices with exactly two negative coordinates have value
$-2$.

The logarithmic Hessian is diagonal:
\[
(\theta_i\theta_jW)=
\begin{pmatrix}
x+x^{-1} & 0 & 0 \\
0 & y+y^{-1} & 0 \\
0 & 0 & z+z^{-1}
\end{pmatrix}.
\]
Therefore
\[
\det(\theta_i\theta_jW)(P_\epsilon)
=8\epsilon_1\epsilon_2\epsilon_3,
\]
which is nonzero at every critical point. Hence all eight critical points are
isolated Morse singularities of type $A_1$. Consequently,
\[
\overline{MF((\mathbb{C}^*)^3,W_{MM_{3-27}})}
\cong \bigoplus_{\epsilon\in\{1,-1\}^3}\MorseCat{3}.
\]
\end{example}

\begin{example}
Consider the Mori--Mukai model $MM_{3-28}$ on the algebraic torus
$(\mathbb{C}^*)^3$, with Laurent potential
\[
W_{MM_{3-28}}=x+\frac{x}{z}+y+z+\frac{1}{y}+\frac{1}{x}.
\]
The variable $y$ is completely decoupled from the $(x,z)$-subsystem. Its
critical equation is
\[
1-\frac{1}{y^2}=0,
\qquad\text{hence}\qquad
y^2=1.
\]
The remaining two equations are
\[
1-\frac{x}{z^2}=0,
\qquad
1+\frac{1}{z}-\frac{1}{x^2}=0.
\]
The first gives
\[
x=z^2.
\]
Substitution into the second gives
\[
1+\frac{1}{z}-\frac{1}{z^4}=0,
\]
or equivalently, since $z\neq 0$ on the torus,
\[
\Phi(z):=z^4+z^3-1=0.
\]
The polynomial $\Phi$ is irreducible and square-free over $\mathbb{Q}$, and is
coprime to $z(z+1)$. Hence all four roots give torus points, and the independent
choice of $y=\pm1$ gives eight critical points:
\[
P_{\lambda,\epsilon}=\left(\lambda^2,\epsilon,\lambda\right),
\qquad
\Phi(\lambda)=0,
\qquad
\epsilon\in\{1,-1\}.
\]
Equivalently, clearing denominators in the logarithmic critical equations and
saturating by $xyz$ gives the lexicographic basis
\[
x-z^2,
\qquad
y^2-1,
\qquad
\Phi(z).
\]

On the branch $y=1$, the critical values are the roots of
\[
w^4-7w^3+10w^2-24w+37=0,
\]
and, modulo $\Phi(\lambda)$,
\[
W_{MM_{3-28}}(P_{\lambda,1})=2\lambda^2+3\lambda+2.
\]
On the branch $y=-1$, the critical values are the roots of
\[
w^4+9w^3+22w^2-24w-91=0,
\]
and, modulo $\Phi(\lambda)$,
\[
W_{MM_{3-28}}(P_{\lambda,-1})=2\lambda^2+3\lambda-2.
\]
Equivalently, the full critical-value polynomial is
\[
\bigl(w^4-7w^3+10w^2-24w+37\bigr)
\bigl(w^4+9w^3+22w^2-24w-91\bigr)=0.
\]

The logarithmic Hessian is nondegenerate at every critical point. More
precisely, modulo $\Phi(\lambda)$ one has
\[
\det(\theta_i\theta_jW)(P_{\lambda,\epsilon})
=2\epsilon\lambda^2(4\lambda+3),
\]
and the displayed factor is coprime to $\Phi$. Therefore all eight critical
points are isolated Morse singularities of type $A_1$. Consequently,
\[
\overline{MF((\mathbb{C}^*)^3,W_{MM_{3-28}})}
\cong \bigoplus_{\epsilon\in\{1,-1\}} \left(\bigoplus_{\lambda:\,\Phi(\lambda)=0}\MorseCat{3}\right).
\]
\end{example}

\begin{example}
Consider the Mori--Mukai model $MM_{3-29}$ on the algebraic torus
$(\mathbb{C}^*)^3$, with Laurent potential
\[
W_{MM_{3-29}}=xy+x+\frac{x}{yz}+y+z+\frac{1}{x}.
\]
The critical equations again collapse by a rational cascade. The $z$-derivative
gives
\[
x=yz^2.
\]
The $y$-derivative gives
\[
x+1=\frac{x}{y^2z}.
\]
Substituting $x=yz^2$ into this equation yields
\[
z=y(x+1).
\]
Returning to $x=yz^2$, one obtains
\[
x=y^3(x+1)^2,
\qquad\text{hence}\qquad
y^3=\frac{x}{(x+1)^2}.
\]
The remaining $x$-derivative is
\[
y+1+\frac{1}{yz}-\frac{1}{x^2}=0.
\]
Using $z=y(x+1)$ and multiplying by $y^2$, this becomes
\[
y^3+\frac{1}{x+1}+y^2\left(1-\frac{1}{x^2}\right)=0.
\]
After substituting $y^3=x/(x+1)^2$, one gets the second rational relation
\[
y^2=\frac{x^2(2x+1)}{(1-x)(x+1)^3}.
\]
Equating $(y^3)^2$ and $(y^2)^3$ gives
\[
(1-x)^3(x+1)^5=x^4(2x+1)^3.
\]
Equivalently,
\[
\Phi(x):=x^8+10x^7+10x^6+x^4+6x^3+2x^2-2x-1=0.
\]
The polynomial $\Phi$ is irreducible and square-free over $\mathbb{Q}$, and is
coprime to $x(x-1)(x+1)(2x+1)$. Hence all eight roots give torus points. Since
$y=y^3/y^2$, the remaining coordinates are rational functions of the root:
\[
y=\frac{1-x^2}{x(2x+1)},
\qquad
z=y(x+1)=-\frac{(x-1)(x+1)^2}{x(2x+1)}.
\]
Thus the critical locus consists of the eight points
\[
P_\lambda=\left(
\lambda,
\frac{1-\lambda^2}{\lambda(2\lambda+1)},
-\frac{(\lambda-1)(\lambda+1)^2}{\lambda(2\lambda+1)}
\right),
\qquad
\Phi(\lambda)=0.
\]
Equivalently, clearing denominators in the logarithmic critical equations and
saturating by $xyz$ gives the lexicographic basis
\[
9y+55x^7+518x^6+246x^5-168x^4+139x^3+256x^2-45x-92,
\]
\[
9z+23x^7+214x^6+78x^5-84x^4+65x^3+101x^2-27x-37,
\]
\[
\Phi(x).
\]

The critical values are the roots of
\[
w^8-22w^6-129w^4-812w^3+1790w^2+8076w+15749=0.
\]
Equivalently, modulo $\Phi(\lambda)$ one has
\fitdisplay{W_{MM_{3-29}}(P_\lambda)=
-\frac{1}{3}\left(20\lambda^7+184\lambda^6+48\lambda^5-84\lambda^4
+62\lambda^3+83\lambda^2-36\lambda-31\right)}
The logarithmic Hessian is nondegenerate at every critical point. More
precisely,
\[
\det(\theta_i\theta_jW)(P_\lambda)
=\frac{2}{3}\left(8\lambda^7+73\lambda^6+18\lambda^5+3\lambda^4
+26\lambda^3+41\lambda^2-7\right),
\]
and the displayed factor is coprime to $\Phi$. Therefore all eight critical
points are isolated Morse singularities of type $A_1$. Consequently,
\[
\overline{MF((\mathbb{C}^*)^3,W_{MM_{3-29}})}
\cong \bigoplus_{\lambda:\,\Phi(\lambda)=0}\MorseCat{3}.
\]
\end{example}

\begin{example}
Consider the Mori--Mukai model $MM_{3-30}$ on the algebraic torus
$(\mathbb{C}^*)^3$, with Laurent potential
\[
W_{MM_{3-30}}=x+\frac{x}{y}+y+\frac{y}{z}+z+\frac{1}{x}.
\]
The critical equations resolve by a direct cascade. The $z$-derivative gives
\[
y=z^2.
\]
The $y$-derivative gives
\[
-\frac{x}{y^2}+1+\frac{1}{z}=0,
\qquad\text{hence}\qquad
x=y^2\frac{z+1}{z}.
\]
Substituting $y=z^2$ gives
\[
x=z^3(z+1).
\]
The remaining $x$-derivative is
\[
1+\frac{1}{y}-\frac{1}{x^2}=0,
\qquad\text{or}\qquad
x^2=\frac{y}{y+1}.
\]
Substitution of $x=z^3(z+1)$ and $y=z^2$ yields
\[
z^6(z+1)^2=\frac{z^2}{z^2+1}.
\]
Since $z\neq0$ on the torus, this is equivalent to
\[
z^4(z+1)^2(z^2+1)=1.
\]
Thus the characteristic equation is
\[
\Phi(z):=z^8+2z^7+2z^6+2z^5+z^4-1=0.
\]
The polynomial $\Phi$ is irreducible and square-free over $\mathbb{Q}$, and is
coprime to $z(z+1)(z^2+1)$. Hence all eight roots give torus points. The
critical locus consists of the eight points
\[
P_\lambda=\left(\lambda^3(\lambda+1),\lambda^2,\lambda\right),
\qquad
\Phi(\lambda)=0.
\]
Equivalently, clearing denominators in the logarithmic critical equations and
saturating by $xyz$ gives the lexicographic basis
\[
x-z^4-z^3,
\qquad
y-z^2,
\qquad
\Phi(z).
\]

The critical values are the roots of
\[
w^8+4w^7-6w^6-100w^5-545w^4-1080w^3
-482w^2+672w+9389=0.
\]
Equivalently, modulo $\Phi(\lambda)$ one has
\[
W_{MM_{3-30}}(P_\lambda)=2\lambda^4+2\lambda^3+3\lambda^2+4\lambda.
\]
The logarithmic Hessian is nondegenerate at every critical point. More
precisely, modulo $\Phi(\lambda)$ one has
\[
\det(\theta_i\theta_jW)(P_\lambda)
=2\lambda^3(\lambda+1)(4\lambda^3+3\lambda^2+3\lambda+2),
\]
and the displayed factor is coprime to $\Phi$. Therefore all eight critical
points are isolated Morse singularities of type $A_1$. Consequently,
\[
\overline{MF((\mathbb{C}^*)^3,W_{MM_{3-30}})}
\cong \bigoplus_{\lambda:\,\Phi(\lambda)=0}\MorseCat{3}.
\]
\end{example}

\begin{example}
Consider the Mori--Mukai model $MM_{3-31}$ on the algebraic torus
$(\mathbb{C}^*)^3$, with Laurent potential
\[
W_{MM_{3-31}}=\frac{x^2}{y^2z}+x+\frac{2x}{y}+y+z+\frac{1}{x}.
\]
The $z$-derivative gives
\[
x^2=y^2z^2,
\]
so a priori there are two branches, $x=yz$ and $x=-yz$. The $y$-derivative is
\[
1-\frac{2x^2}{y^3z}-\frac{2x}{y^2}=0.
\]
After multiplying by $y^3z$ and using $x^2=y^2z^2$, this becomes
\[
y^3z-2y^2z^2-2xyz=0,
\]
or equivalently
\[
y-2z-2\frac{x}{y}=0.
\]
On the branch $x=-yz$ this equation gives $y=0$, which is impossible on the
torus. Hence the negative branch is empty. On the branch $x=yz$ one obtains
\[
y=4z,
\qquad
x=4z^2.
\]
The remaining $x$-derivative is
\[
\frac{2x}{y^2z}+1+\frac{2}{y}-\frac{1}{x^2}=0.
\]
Using $x=yz$, this reduces to
\[
1+\frac{4}{y}-\frac{1}{x^2}=0.
\]
Substituting $y=4z$ and $x=4z^2$ gives
\[
1+\frac{1}{z}-\frac{1}{16z^4}=0.
\]
Thus the characteristic equation is
\[
\Phi(z):=16z^4+16z^3-1=0.
\]
The polynomial $\Phi$ is irreducible and square-free over $\mathbb{Q}$, and is
coprime to $z$. Hence all four roots give torus points. The critical locus
consists of the four points
\[
P_\lambda=\left(4\lambda^2,4\lambda,\lambda\right),
\qquad
\Phi(\lambda)=0.
\]
Equivalently, clearing denominators in the logarithmic critical equations and
saturating by $xyz$ gives the lexicographic basis
\[
x-4z^2,
\qquad
y-4z,
\qquad
\Phi(z).
\]

The critical values are the roots of
\[
w^4+4w^3-8w^2-144w-416=0.
\]
Equivalently, modulo $\Phi(\lambda)$ one has
\[
W_{MM_{3-31}}(P_\lambda)=8\lambda^2+12\lambda.
\]
The logarithmic Hessian is nondegenerate at every critical point. More
precisely, modulo $\Phi(\lambda)$ one has
\[
\det(\theta_i\theta_jW)(P_\lambda)
=-8(4\lambda^3-1),
\]
and the displayed factor is coprime to $\Phi$. Therefore all four critical
points are isolated Morse singularities of type $A_1$. Consequently,
\[
\overline{MF((\mathbb{C}^*)^3,W_{MM_{3-31}})}
\cong \bigoplus_{\lambda:\,\Phi(\lambda)=0}\MorseCat{3}.
\]
\end{example}

\begin{example}
Consider the Mori--Mukai model $MM_{4-1}$ on the algebraic torus
$(\mathbb{C}^*)^3$, with Laurent potential
\[
W_{MM_{4-1}}=x+yz^2+2yz+y+2z+\frac{1}{y}+\frac{2}{yz}
+\frac{1}{x}+\frac{2}{xyz}+\frac{1}{xy^2z^2}.
\]
The last three terms form a square after extracting the factor $1/x$:
\[
\frac{1}{x}+\frac{2}{xyz}+\frac{1}{xy^2z^2}
=\frac{1}{x}\left(1+\frac{1}{yz}\right)^2.
\]
Set
\[
U=1+\frac{1}{yz}.
\]
Since also
\[
yz^2+2yz+y=y(z+1)^2,
\]
the potential becomes
\[
W_{MM_{4-1}}=x+y(z+1)^2+2z+\frac{1}{y}+\frac{2}{yz}+\frac{U^2}{x}.
\]
The $x$-derivative gives
\[
1-\frac{U^2}{x^2}=0,
\qquad\text{hence}\qquad
x=\pm U.
\]

On the negative branch $x=-U$, the factor $1+U/x$ vanishes. The remaining
transverse equations reduce to
\[
(z+1)^2-\frac{1}{y^2}=0,
\qquad
2y(z+1)+2=0.
\]
Thus
\[
y=-\frac{1}{z+1},
\qquad
x=-U=\frac{1}{z}.
\]
This gives the critical curve
\[
\Gamma=\left\{xz=1,\ y(z+1)+1=0\right\}
\subset (\mathbb{C}^*)^3,
\]
parametrized by
\[
\gamma(t)=\left(\frac{1}{t},-\frac{1}{t+1},t\right),
\qquad
t\in\mathbb{C}^*,\quad t\neq -1.
\]
Along this curve the potential is constant:
\[
W_{MM_{4-1}}|_{\Gamma}=-4.
\]
The transverse formal model is global. Set
\[
A=x+U,
\qquad
B=y(z+1)+1.
\]
Then one has the exact identity
\[
W_{MM_{4-1}}+4=\frac{A^2}{x}+\frac{B^2}{y}.
\]
Moreover, if $F=xz-1$, then
\[
F=zA-\frac{B}{y}.
\]
Since $y$ and $z$ are units on the torus, $(A,B)=(F,B)$ is the ideal of
$\Gamma$. With
\[
A_{\Gamma}=\mathbb C[t^{\pm1},(t+1)^{-1}],
\]
the completed neighborhood is therefore
\[
\widehat X_{\Gamma}
\simeq
\operatorname{Spf}\!\left(A_{\Gamma}[[A,B]]\right).
\]
The quadratic part of the shifted potential is
\[
tA^2-(t+1)B^2,
\]
whose relative Hessian determinant is $-4t(t+1)$, a unit in
$A_{\Gamma}$. Its square class is the nontrivial class $[t(t+1)]$. Hence the
relative formal Morse lemma gives a global nondegenerate quadratic bundle with
that Clifford twist.

On the positive branch $x=U$, the factor $1+U/x$ equals $2$. The two transverse
equations become
\[
y^2z(z+1)^2=z+4,
\qquad
y^2z^2(z+1)+yz^2=2.
\]
After saturating by the torus variables, their lexicographic basis is
\[
y-2,
\qquad
2z-1.
\]
Hence the positive branch consists of the single point
\[
P_0=\left(2,2,\frac{1}{2}\right),
\]
and
\[
W_{MM_{4-1}}(P_0)=12.
\]

Equivalently, clearing denominators in the logarithmic critical equations and
saturating by $xyz$ gives the reduced critical locus determined by
\[
xz-1,
\qquad
xy+x-2yz-y-2,
\]
\[
(y-2)(yz+y+1),
\qquad
(2z-1)(yz+y+1),
\]
which decomposes as the disjoint union of the curve $\Gamma$ and the point
$P_0$.

The logarithmic Hessian has determinant zero along $\Gamma$, while the minor
with rows $(x,y)$ and columns $(x,z)$ is the unit $-4$, consistently with the
global Morse--Bott model above. At the isolated point,
\[
(\theta_i\theta_jW)(P_0)=
\begin{pmatrix}
4 & 2 & 2 \\
2 & 10 & 8 \\
2 & 8 & 10
\end{pmatrix},
\qquad
\det(\theta_i\theta_jW)(P_0)=128.
\]
Therefore $P_0$ is an isolated Morse singularity of type $A_1$. By
Proposition~\ref{prop:morse_bott_clifford}, the curve contributes
$\MBtwcat{\Gamma}{t(t+1)}$, while the isolated point contributes
$\MorseCat{3}$.
Consequently,
\[
\overline{MF((\mathbb{C}^*)^3,W_{MM_{4-1}})}
\cong \MBtwcat{\Gamma}{t(t+1)}\oplus \MorseCat{3}.
\]
\end{example}

\begin{example}
Consider the Mori--Mukai model $MM_{4-2}$ on the algebraic torus
$(\mathbb{C}^*)^3$, with Laurent potential
\[
W_{MM_{4-2}}=x+y+\frac{y}{z}+z+\frac{1}{z}+\frac{2}{y}
+\frac{y}{x}+\frac{3}{x}+\frac{3}{xy}+\frac{1}{xy^2}.
\]
The last four terms with denominator $x$ form a cubic block:
\[
\frac{y}{x}+\frac{3}{x}+\frac{3}{xy}+\frac{1}{xy^2}
=\frac{1}{x}\left(y+3+\frac{3}{y}+\frac{1}{y^2}\right)
=\frac{(y+1)^3}{xy^2}.
\]
Thus
\[
W_{MM_{4-2}}=x+y+\frac{y}{z}+z+\frac{1}{z}+\frac{2}{y}
+\frac{(y+1)^3}{xy^2}.
\]
The $z$-derivative gives
\[
1-\frac{y+1}{z^2}=0,
\qquad\text{hence}\qquad
y=z^2-1.
\]
The $x$-derivative gives
\[
1-\frac{(y+1)^3}{x^2y^2}=0.
\]
Using $y+1=z^2$ and $y=z^2-1$, this becomes
\[
x^2=\frac{z^6}{(z^2-1)^2},
\]
so the critical locus splits into two branches
\[
x=\pm\frac{z^3}{z^2-1}.
\]
The remaining $y$-derivative is
\[
1+\frac{1}{z}-\frac{2}{y^2}
+\frac{(y+1)^2(y-2)}{xy^3}=0.
\]

On the positive branch $x=z^3/(z^2-1)$, substituting $y=z^2-1$ gives
\[
\frac{z^3-3z+1}{z(z-1)^2}=0.
\]
Equivalently, after multiplying by $z(z^2-1)^2$, the forbidden factor
$(z+1)^2$ appears and is removed by the torus condition $y\neq0$. Hence the
positive branch is governed by
\[
p_+(z):=z^3-3z+1=0.
\]
On the negative branch $x=-z^3/(z^2-1)$, one obtains
\[
\frac{z^3+2z^2+z+1}{z(z+1)^2}=0,
\]
and the forbidden factor $(z-1)^2$ is removed in the same way. Hence the
negative branch is governed by
\[
p_-(z):=z^3+2z^2+z+1=0.
\]
The polynomials $p_+$ and $p_-$ are irreducible and square-free over
$\mathbb{Q}$, are coprime to each other, and are coprime to $z(z^2-1)$. Thus
all six roots give torus points. The critical points are
\[
P^+_\alpha=\left(
\frac{\alpha^3}{\alpha^2-1},\alpha^2-1,\alpha
\right),
\qquad
p_+(\alpha)=0,
\]
and
\[
P^-_\beta=\left(
-\frac{\beta^3}{\beta^2-1},\beta^2-1,\beta
\right),
\qquad
p_-(\beta)=0.
\]
Equivalently, after clearing denominators and saturating by the torus variables,
the two branch lexicographic bases are
\[
3x-z^2-5z+2,
\qquad
y-z^2+1,
\qquad
p_+(z),
\]
and
\[
5x-3z^2+z-2,
\qquad
y-z^2+1,
\qquad
p_-(z).
\]

On the positive branch, the critical values are the roots of
\[
w^3-3w^2-78w-163=0,
\]
and, modulo $p_+(\alpha)$,
\[
W_{MM_{4-2}}(P^+_\alpha)=3\alpha^2+6\alpha-5.
\]
On the negative branch, the critical values are the roots of
\[
w^3+w^2-18w-43=0,
\]
and, modulo $p_-(\beta)$,
\[
W_{MM_{4-2}}(P^-_\beta)=3\beta^2+2\beta-1.
\]

The logarithmic Hessian is nondegenerate at every critical point. More
precisely,
\[
\det(\theta_i\theta_jW)(P^+_\alpha)
=6(4\alpha^2+5\alpha-2),
\]
and
\[
\det(\theta_i\theta_jW)(P^-_\beta)
=-2(12\beta^2+5\beta+8),
\]
with the displayed factors coprime to $p_+$ and $p_-$, respectively. Therefore
all six critical points are isolated Morse singularities of type $A_1$.
Consequently,
\[
\overline{MF((\mathbb{C}^*)^3,W_{MM_{4-2}})}
\cong \left(\bigoplus_{\alpha:\,p_+(\alpha)=0}\MorseCat{3}\right) \oplus \left(\bigoplus_{\beta:\,p_-(\beta)=0}\MorseCat{3}\right).
\]
\end{example}

\begin{example}
Consider the Mori--Mukai model $MM_{4-3}$ on the algebraic torus
$(\mathbb{C}^*)^3$, with Laurent potential
\[
W_{MM_{4-3}}=x+\frac{2x}{yz}+\frac{x}{y^2z^2}
+yz^2+2yz+y+\frac{2}{yz}+\frac{1}{x}.
\]
Set
\[
U=1+\frac{1}{yz}.
\]
Then the two square blocks in the potential give
\[
x+\frac{2x}{yz}+\frac{x}{y^2z^2}=xU^2,
\qquad
yz^2+2yz+y=y(z+1)^2,
\]
and therefore
\[
W_{MM_{4-3}}=xU^2+y(z+1)^2+2U-2+\frac{1}{x}.
\]
The $x$-derivative gives
\[
U^2-\frac{1}{x^2}=0,
\qquad\text{hence}\qquad
xU=\pm1.
\]
The remaining derivatives can be written as
\[
y^2z(z+1)^2=2(xU)+2,
\qquad
2y^2z^2(z+1)=2(xU)+2.
\]

On the negative branch $xU=-1$, both right-hand sides vanish. Since $y$ and
$z$ are invertible on the torus, this forces $z=-1$. The equation $xU=-1$
then gives
\[
x=\frac{y}{1-y}.
\]
Thus this branch is the critical curve
\[
\Gamma=\left\{z+1=0,\ x(1-y)-y=0\right\}
\subset (\mathbb{C}^*)^3,
\]
parametrized by
\[
\gamma(t)=\left(\frac{t}{1-t},t,-1\right),
\qquad
t\in\mathbb{C}^*,\quad t\neq1.
\]
Along this curve the potential is constant:
\[
W_{MM_{4-3}}|_{\Gamma}=-2.
\]
The transverse formal model is global. Set
\[
A=xU+1,
\qquad
B=z+1.
\]
Completing the square gives the exact identity
\[
W_{MM_{4-3}}+2=\frac{A^2}{x}+yB^2.
\]
If $F=x(1-y)-y$, then
\[
F=yzA-y(x+1)B.
\]
Thus $(A,B)=(F,B)$ is the ideal of $\Gamma$. With
\[
A_{\Gamma}=\mathbb C[t^{\pm1},(t-1)^{-1}],
\]
one obtains
\[
\widehat X_{\Gamma}
\simeq
\operatorname{Spf}\!\left(A_{\Gamma}[[A,B]]\right).
\]
The quadratic part of the shifted potential is
\[
\frac{1-t}{t}A^2+tB^2,
\]
whose relative Hessian determinant is $4(1-t)$, a unit in $A_{\Gamma}$.
Its square class $[1-t]$ is nontrivial. Therefore the relative formal Morse
lemma gives a global nondegenerate quadratic bundle with the corresponding
Clifford twist.

On the positive branch $xU=1$, both right-hand sides are equal to $4$. Since
the left-hand sides are then nonzero, one may divide by $y^2z(z+1)$ and obtain
\[
z+1=2z,
\qquad\text{hence}\qquad
z=1.
\]
Substitution in $y^2z(z+1)^2=4$ gives $y^2=1$. The solution $y=-1$ gives
$U=0$, contradicting $xU=1$. Hence $y=1$, $U=2$, and
\[
P_0=\left(\frac{1}{2},1,1\right).
\]
At this point
\[
W_{MM_{4-3}}(P_0)=10.
\]

Equivalently, clearing denominators in the logarithmic critical equations and
saturating by $xyz$ gives the reduced critical locus determined by
\[
2xy-2x+2y-z-1,
\qquad
(2x-1)(z+1),
\]
\[
(y-1)(z+1),
\qquad
(z-1)(z+1),
\]
which decomposes as the disjoint union of the curve $\Gamma$ and the point
$P_0$.

The logarithmic Hessian has determinant zero along $\Gamma$, while the minor
with rows $(x,z)$ and columns $(x,z)$ is the unit $-4(y-1)$ on $\Gamma$,
consistently with the global Morse--Bott model above. At the isolated point,
\[
(\theta_i\theta_jW)(P_0)=
\begin{pmatrix}
4 & -2 & -2 \\
-2 & 9 & 9 \\
-2 & 9 & 11
\end{pmatrix},
\qquad
\det(\theta_i\theta_jW)(P_0)=64.
\]
Therefore $P_0$ is an isolated Morse singularity of type $A_1$. By
Proposition~\ref{prop:morse_bott_clifford}, the curve contributes
$\MBtwcat{\Gamma}{1-t}$ and the isolated point contributes $\MorseCat{3}$.
Consequently,
\[
\overline{MF((\mathbb{C}^*)^3,W_{MM_{4-3}})}
\cong \MBtwcat{\Gamma}{1-t}\oplus \MorseCat{3}.
\]
\end{example}

\begin{example}
Consider the Mori--Mukai model $MM_{4-4}$ on the algebraic torus
$(\mathbb{C}^*)^3$, with Laurent potential
\[
W_{MM_{4-4}}=x+\frac{x}{y}+y+z+\frac{1}{z}+\frac{1}{y}
+\frac{2}{x}+\frac{2}{xz}+\frac{1}{x^2z}.
\]
The terms with denominator $z$ contain the square
\[
\frac{1}{z}+\frac{2}{xz}+\frac{1}{x^2z}
=\frac{1}{z}\left(1+\frac{1}{x}\right)^2
=\frac{1}{z}\left(\frac{x+1}{x}\right)^2.
\]
Also
\[
\frac{x}{y}+\frac{1}{y}=\frac{x+1}{y}.
\]
Thus the potential can be written as
\[
W_{MM_{4-4}}=x+z+\frac{2}{x}+y+\frac{x+1}{y}
+\frac{1}{z}\left(\frac{x+1}{x}\right)^2.
\]
The $y$-derivative gives
\[
1-\frac{x+1}{y^2}=0,
\qquad\text{hence}\qquad
y^2=x+1.
\]
The $z$-derivative gives
\[
1-\frac{1}{z^2}\left(\frac{x+1}{x}\right)^2=0,
\]
so the critical locus splits into two possible branches
\[
z=\pm\frac{x+1}{x}.
\]
The remaining $x$-derivative is
\[
1+\frac{1}{y}-\frac{2}{x^2}
-\frac{2(x+1)}{x^3z}=0.
\]

On the negative branch $z=-(x+1)/x$, this equation reduces to
\[
1+\frac{1}{y}=0,
\]
and hence $y=-1$. But the equation $y^2=x+1$ then gives $x=0$, which is
forbidden on the torus. Therefore the negative branch is empty.

On the positive branch $z=(x+1)/x$, the $x$-equation becomes
\[
1+\frac{1}{y}-\frac{4}{x^2}=0,
\]
so
\[
y=\frac{x^2}{4-x^2}.
\]
Combining this with $y^2=x+1$ gives
\[
\left(\frac{x^2}{4-x^2}\right)^2=x+1.
\]
Equivalently,
\[
\Phi(x):=x^5-8x^3-8x^2+16x+16=0.
\]
The polynomial $\Phi$ is irreducible and square-free over $\mathbb{Q}$, and is
coprime to $x(x+1)(4-x^2)$. Hence all five roots give torus points. The
critical locus consists of the five points
\[
P_\lambda=\left(
\lambda,
\frac{\lambda^2}{4-\lambda^2},
\frac{\lambda+1}{\lambda}
\right),
\qquad
\Phi(\lambda)=0.
\]
Equivalently, clearing denominators in the logarithmic critical equations and
saturating by $xyz$ gives the lexicographic basis
\[
x^3-4x+4y-4,
\qquad
x^4-8x^2-8x+16z,
\]
\[
\Phi(x).
\]

The critical values are the roots of
\[
w^5-4w^4-44w^3-80w^2+192w+592=0.
\]
Equivalently, modulo $\Phi(\lambda)$ one has
\[
W_{MM_{4-4}}(P_\lambda)=
-\frac{1}{4}\lambda(\lambda^3+2\lambda^2-8\lambda-20).
\]
The logarithmic Hessian is nondegenerate at every critical point. More
precisely,
\[
\det(\theta_i\theta_jW)(P_\lambda)
=-2(\lambda^4+\lambda^3-4\lambda^2-11\lambda-8),
\]
and the displayed factor is coprime to $\Phi$. Therefore all five critical
points are isolated Morse singularities of type $A_1$. Consequently,
\[
\overline{MF((\mathbb{C}^*)^3,W_{MM_{4-4}})}
\cong \bigoplus_{\lambda:\,\Phi(\lambda)=0}\MorseCat{3}.
\]
\end{example}

\begin{example}
Consider the Mori--Mukai model $MM_{4-5}$ on the algebraic torus
$(\mathbb{C}^*)^3$, with Laurent potential
\[
W_{MM_{4-5}}=x+\frac{x}{z}+y+z+\frac{1}{y}
+\frac{2y}{x}+\frac{2}{x}+\frac{y}{x^2}.
\]
The three terms with numerator $y$ form the square
\[
y+\frac{2y}{x}+\frac{y}{x^2}
=y\left(1+\frac{2}{x}+\frac{1}{x^2}\right)
=y\left(\frac{x+1}{x}\right)^2.
\]
Thus
\[
W_{MM_{4-5}}=x+\frac{x}{z}+z+\frac{1}{y}+\frac{2}{x}
+y\left(\frac{x+1}{x}\right)^2.
\]
The $z$-derivative gives
\[
1-\frac{x}{z^2}=0,
\qquad\text{hence}\qquad
x=z^2.
\]
The $y$-derivative gives
\[
-\frac{1}{y^2}+\left(\frac{x+1}{x}\right)^2=0,
\qquad\text{hence}\qquad
y=\pm\frac{x}{x+1}.
\]
The remaining $x$-derivative is
\[
1+\frac{1}{z}-\frac{2}{x^2}
-\frac{2y(x+1)}{x^3}=0.
\]

On the negative branch $y=-x/(x+1)$, this equation reduces to
\[
1+\frac{1}{z}=0,
\]
so $z=-1$. Since $x=z^2$, this gives $x=1$, and hence
$y=-1/2$. Thus the negative branch contributes the rational critical point
\[
P_0=\left(1,-\frac{1}{2},-1\right).
\]

On the positive branch $y=x/(x+1)$, the $x$-equation becomes
\[
1+\frac{1}{z}-\frac{4}{x^2}=0.
\]
Using $x=z^2$, this is equivalent to
\[
q(z):=z^4+z^3-4=0.
\]
The polynomial $q$ is irreducible and square-free over $\mathbb{Q}$, and is
coprime to $z(z^2+1)$. Hence its four roots give torus points
\[
P_\alpha=\left(
\alpha^2,
\frac{\alpha^2}{\alpha^2+1},
\alpha
\right),
\qquad
q(\alpha)=0.
\]
Equivalently, clearing denominators in the logarithmic critical equations and
saturating by $xyz$ gives the reduced critical locus determined by the
lexicographic basis
\[
x-z^2,
\qquad
20y-9z^4-11z^3+4z^2+8z+12,
\]
\[
(z+1)(z^4+z^3-4).
\]

The rational point has critical value
\[
W_{MM_{4-5}}(P_0)=-3.
\]
On the quartic branch, the critical values are the roots of
\[
w^4-7w^3-14w^2-36w+316=0,
\]
and, modulo $q(\alpha)$,
\[
W_{MM_{4-5}}(P_\alpha)=2\alpha^2+3\alpha+2.
\]
The logarithmic Hessian is nondegenerate at every critical point. At the
rational point,
\[
(\theta_i\theta_jW)(P_0)=
\begin{pmatrix}
-1 & 2 & 1 \\
2 & -4 & 0 \\
1 & 0 & -2
\end{pmatrix},
\qquad
\det(\theta_i\theta_jW)(P_0)=4.
\]
On the quartic branch,
\[
\det(\theta_i\theta_jW)(P_\alpha)
=2(4\alpha+3)(\alpha^2+1),
\]
and the displayed factor is coprime to $q$. Therefore all five critical points
are isolated Morse singularities of type $A_1$. Consequently,
\[
\overline{MF((\mathbb{C}^*)^3,W_{MM_{4-5}})}
\cong \MorseCat{3}\oplus \bigoplus_{\alpha:\,q(\alpha)=0}\MorseCat{3}.
\]
\end{example}

\begin{example}
Consider the Mori--Mukai model $MM_{4-6}$ on the algebraic torus
$(\mathbb{C}^*)^3$, with Laurent potential
\[
W_{MM_{4-6}}=x+\frac{x}{y}+y+z+\frac{1}{y}
+\frac{y}{x}+\frac{2}{x}+\frac{y}{x^2z}+\frac{1}{x^2z}.
\]
The $z$-derivative gives
\[
1-\frac{y+1}{x^2z^2}=0,
\qquad\text{hence}\qquad
x^2z^2=y+1.
\]
Introduce the torus coordinate
\[
T=xz.
\]
Then, on the $z$-critical locus,
\[
y=T^2-1,
\qquad
z=\frac{T}{x}.
\]
Substituting these relations into the potential gives the reduced two-variable
potential
\[
W_{MM_{4-6}}(x,T)=
\frac{xT^2}{T^2-1}
+\frac{T^4-2T^2+2}{T^2-1}
+\frac{(T+1)^2}{x}.
\]
Indeed, the terms with denominator $x$ combine as
\[
\frac{y}{x}+\frac{2}{x}+\frac{y+1}{x^2z}
=\frac{T^2-1}{x}+\frac{2}{x}+\frac{T^2}{xT}
=\frac{(T+1)^2}{x}.
\]

The $x$-derivative of the reduced potential gives
\[
\frac{T^2}{T^2-1}-\frac{(T+1)^2}{x^2}=0,
\qquad\text{hence}\qquad
x^2=\frac{(T-1)(T+1)^3}{T^2}.
\]
The $T$-derivative gives
\[
-\frac{2Tx}{(T^2-1)^2}
+\frac{2T^3(T^2-2)}{(T^2-1)^2}
+\frac{2(T+1)}{x}=0.
\]
After multiplying by $x(T^2-1)^2/(2T)$, this becomes
\[
x^2-T^2(T^2-2)x-\frac{(T-1)^2(T+1)^3}{T}=0.
\]
Using the first equation to eliminate $x^2$ gives
\[
x=-\frac{(T-1)(T+1)^3(T^2-T-1)}{T^4(T^2-2)}.
\]
Squaring this expression and comparing again with
$x^2=(T-1)(T+1)^3/T^2$ yields the characteristic polynomial
\[
P(T):=T^{10}-5T^8+9T^6+2T^5-8T^4
-6T^3+3T^2+4T+1=0.
\]
The polynomial $P$ is irreducible and square-free over $\mathbb{Q}$, and is
coprime to
\[
T(T^2-1)(T^2-2)(T^2-T-1).
\]
Hence all ten roots give torus points. Explicitly, the critical points are
\fitdisplay{P_\tau=\left(
-\frac{(\tau-1)(\tau+1)^3(\tau^2-\tau-1)}{\tau^4(\tau^2-2)},
\tau^2-1,
-\frac{\tau^5(\tau^2-2)}{(\tau-1)(\tau+1)^3(\tau^2-\tau-1)}
\right)}
where $P(\tau)=0$.
Equivalently, after adjoining $T=xz$ and saturating by the torus variables, a
lexicographic basis for the reduced critical scheme is
\[
x-2T^9+T^8+9T^7-5T^6-14T^5+4T^4+12T^3+5T^2-7T-3,
\]
\[
y-T^2+1,
\qquad
z-T^8+T^7+3T^6-3T^5-3T^4+T^3+4T^2-1,
\]
\[
P(T).
\]

The critical values are the roots of
\[
\begin{aligned}
0={}&w^{10}+10w^9-37w^8-808w^7-3137w^6-48w^5
+26995w^4 \\
&+65954w^3+50346w^2-18344w-34171.
\end{aligned}
\]
Equivalently, modulo $P(\tau)$ one has
\[
\begin{aligned}
W_{MM_{4-6}}(P_\tau)={}&
\tau^8+2\tau^7-4\tau^6-6\tau^5+7\tau^4+8\tau^3 \\
&-2\tau^2-8\tau-5.
\end{aligned}
\]
The logarithmic Hessian is nondegenerate at every critical point. More
precisely, modulo $P(\tau)$,
\[
\det(\theta_i\theta_jW)(P_\tau)
=2(\tau^9-\tau^8-7\tau^7+7\tau^6+16\tau^5-11\tau^4
-13\tau^3+13\tau+5),
\]
and the displayed factor is coprime to $P$. Therefore all ten critical points
are isolated Morse singularities of type $A_1$. Consequently,
\[
\overline{MF((\mathbb{C}^*)^3,W_{MM_{4-6}})}
\cong \bigoplus_{\tau:\,P(\tau)=0}\MorseCat{3}.
\]
\end{example}

\begin{example}
Consider the Mori--Mukai model $MM_{4-7}$ on the algebraic torus
$(\mathbb{C}^*)^3$, with Laurent potential
\[
W_{MM_{4-7}}=x+\frac{x}{y}+y+z+\frac{1}{y}
+\frac{z}{x}+\frac{2}{x}+\frac{1}{xz}.
\]
The three terms with denominator $x$ form the square
\[
\frac{z}{x}+\frac{2}{x}+\frac{1}{xz}
=\frac{1}{x}\left(z+2+\frac{1}{z}\right)
=\frac{(z+1)^2}{xz}.
\]
Also
\[
y+\frac{x}{y}+\frac{1}{y}=y+\frac{x+1}{y}.
\]
Thus
\[
W_{MM_{4-7}}=x+y+z+\frac{x+1}{y}+\frac{(z+1)^2}{xz}.
\]
The $y$-derivative gives
\[
1-\frac{x+1}{y^2}=0,
\qquad\text{hence}\qquad
y^2=x+1.
\]
The $z$-derivative gives
\[
1+\frac{1}{x}-\frac{1}{xz^2}=0,
\qquad\text{hence}\qquad
z^2(x+1)=1.
\]
Combining the two equations gives $(yz)^2=1$, so the critical locus splits
into the two branches
\[
yz=1,
\qquad
yz=-1.
\]
The remaining $x$-derivative is
\[
1+\frac{1}{y}-\frac{(z+1)^2}{x^2z}=0.
\]

On the positive branch $yz=1$, so $z=1/y$. The $x$-equation becomes
\[
x^2(y+1)=(y+1)^2.
\]
If $y=-1$, then $y^2=x+1$ gives $x=0$, which is forbidden on the torus.
Thus $y\neq -1$, and the equation reduces to $x^2=y+1$. Together with
$y^2=x+1$ this gives
\[
(x-y)(x+y+1)=0.
\]
The subbranch $x+y+1=0$ gives $x(x+1)=0$, hence either $x=0$ or $y=0$;
both are forbidden on the torus. Therefore $x=y$, and the positive branch is
governed by
\[
p_+(y):=y^2-y-1=0.
\]
It contributes the two critical points
\[
P^+_\alpha=\left(\alpha,\alpha,\frac{1}{\alpha}\right),
\qquad
p_+(\alpha)=0.
\]

On the negative branch $yz=-1$, so $z=-1/y$. The $x$-equation becomes
\[
x^2(y+1)=-(y-1)^2.
\]
Using $x=y^2-1$, this is
\[
(y-1)^2(y+1)^3=-(y-1)^2.
\]
If $y=1$, then $x=0$, again forbidden on the torus. Hence
\[
(y+1)^3+1=0,
\]
or equivalently
\[
p_-(y):=(y+2)(y^2+y+1)=0.
\]
The negative branch contributes the three critical points
\[
P^-_\beta=\left(\beta^2-1,\beta,-\frac{1}{\beta}\right),
\qquad
p_-(\beta)=0.
\]
The polynomials $p_+$ and $p_-$ are square-free, and their roots avoid the
torus boundary. Thus the critical locus consists of five torus points.
Equivalently, clearing denominators in the logarithmic critical equations and
saturating by $xyz$ gives the lexicographic basis
\[
x+8z^4-2z^3-9z^2+18z-10,
\]
\[
y-18z^4-z^3+17z^2-36z+17,
\]
\[
(2z-1)(z^2-z+1)(z^2+z-1).
\]

On the positive branch, the critical values are the roots of
\[
w^2-5w-25=0,
\]
and, modulo $p_+(\alpha)$,
\[
W_{MM_{4-7}}(P^+_\alpha)=5\alpha.
\]
On the negative branch, the critical values are the roots of
\[
(w-1)(w^2+7w+19)=0,
\]
and, modulo $p_-(\beta)$,
\[
W_{MM_{4-7}}(P^-_\beta)=3\beta^2+6\beta+1.
\]

The logarithmic Hessian is nondegenerate at every critical point. More
precisely,
\[
\det(\theta_i\theta_jW)(P^+_\alpha)=2(7\alpha+4),
\]
and
\[
\det(\theta_i\theta_jW)(P^-_\beta)=-6(\beta-1)(\beta+1),
\]
with the displayed factors coprime to $p_+$ and $p_-$, respectively. Therefore
all five critical points are isolated Morse singularities of type $A_1$.
Consequently,
\[
\overline{MF((\mathbb{C}^*)^3,W_{MM_{4-7}})}
\cong \left(\bigoplus_{\alpha:\,p_+(\alpha)=0}\MorseCat{3}\right) \oplus \left(\bigoplus_{\beta:\,p_-(\beta)=0}\MorseCat{3}\right).
\]
\end{example}

\begin{example}
Consider the Mori--Mukai model $MM_{4-8}$ on the algebraic torus
$(\mathbb{C}^*)^3$, with Laurent potential
\[
W_{MM_{4-8}}=x+y+z+\frac{1}{z}+\frac{z}{y}+\frac{1}{y}
+\frac{1}{x}+\frac{1}{xz}.
\]
The terms with denominator $y$ and $x$ combine as
\[
\frac{z}{y}+\frac{1}{y}=\frac{z+1}{y},
\qquad
\frac{1}{x}+\frac{1}{xz}=\frac{z+1}{xz}.
\]
Thus
\[
W_{MM_{4-8}}=x+y+z+\frac{1}{z}+\frac{z+1}{y}
+\frac{z+1}{xz}.
\]
The $y$-derivative gives
\[
1-\frac{z+1}{y^2}=0,
\qquad\text{hence}\qquad
y^2=z+1.
\]
The $x$-derivative gives
\[
1-\frac{z+1}{x^2z}=0,
\qquad\text{hence}\qquad
x^2=\frac{z+1}{z}.
\]
Dividing these two equations gives $y^2/x^2=z$. Introduce the torus coordinate
\[
T=\frac{y}{x}.
\]
Then
\[
z=T^2,
\qquad
y=xT,
\qquad
y^2=T^2+1,
\qquad
x^2=\frac{T^2+1}{T^2}.
\]
The remaining $z$-derivative is
\[
1-\frac{1}{z^2}+\frac{1}{y}-\frac{1}{xz^2}=0.
\]
After multiplying by $z^2$ and substituting $z=T^2$ and $y=xT$, this becomes
\[
T^4-1+\frac{T^3-1}{x}=0.
\]
Factoring the two differences of powers gives
\[
(T-1)\left((T+1)(T^2+1)+\frac{T^2+T+1}{x}\right)=0.
\]

The first branch is $T=1$. Then $z=1$ and $x=y$, while $y^2=2$. Hence it
contributes the two quadratic critical points
\[
P^0_\sigma=(\sigma,\sigma,1),
\qquad
\sigma^2-2=0.
\]

On the second branch,
\[
\frac{1}{x}=-\frac{(T+1)(T^2+1)}{T^2+T+1}.
\]
Squaring and comparing with $1/x^2=T^2/(T^2+1)$ yields
\[
(T+1)^2(T^2+1)^3=T^2(T^2+T+1)^2.
\]
Equivalently,
\[
\begin{aligned}
R(T):={}&T^8+2T^7+3T^6+4T^5+3T^4+4T^3 \\
&+3T^2+2T+1=0.
\end{aligned}
\]
This reciprocal polynomial satisfies, after setting $u=T+T^{-1}$,
\[
T^{-4}R(T)=u^4+2u^3-u^2-2u-1,
\]
or equivalently
\[
(u^2+u-1)^2=2.
\]
The polynomial $R$ is irreducible and square-free over $\mathbb{Q}$, and is
coprime to
\[
T(T+1)(T^2+1)(T^2+T+1).
\]
Hence all eight roots give torus points. Explicitly, they are
\[
P_\tau=\left(
-\frac{\tau^2+\tau+1}{(\tau+1)(\tau^2+1)},
-\frac{\tau(\tau^2+\tau+1)}{(\tau+1)(\tau^2+1)},
\tau^2
\right),
\]
where $R(\tau)=0$.
Equivalently, after adjoining $T=y/x$ and saturating by the torus variables,
a lexicographic basis for the main branch is
\[
x+T^7+2T^6+3T^5+3T^4+2T^3+3T^2+2T+2,
\]
\[
y-T^5-T^4-T^3-T^2-1,
\qquad
z-T^2,
\qquad
R(T).
\]

On the identity branch, the critical values are the roots of
\[
w^2-4w-28=0,
\]
and, modulo $\sigma^2-2$, one has
\[
W_{MM_{4-8}}(P^0_\sigma)=4\sigma+2.
\]
On the main branch, the critical-value resultant is
\[
\left(w^4+6w^3+7w^2-6w-31\right)^2,
\]
and, modulo $R(\tau)$,
\[
W_{MM_{4-8}}(P_\tau)
=-\tau^6+\tau^4+2\tau^2-2\tau-1.
\]

The logarithmic Hessian is nondegenerate at every critical point. On the
identity branch,
\[
\det(\theta_i\theta_jW)(P^0_\sigma)=2(3\sigma+8),
\]
which is coprime to $\sigma^2-2$. On the main branch,
\[
\det(\theta_i\theta_jW)(P_\tau)
=-2\tau(\tau^6-3\tau^5-3\tau^4-3\tau^3-9\tau^2+2\tau-8),
\]
and the displayed factor is coprime to $R$. Therefore all ten critical points
are isolated Morse singularities of type $A_1$. Consequently,
\[
\overline{MF((\mathbb{C}^*)^3,W_{MM_{4-8}})}
\cong \left(\bigoplus_{\sigma^2=2}\MorseCat{3}\right) \oplus \left(\bigoplus_{\tau:\,R(\tau)=0}\MorseCat{3}\right).
\]
\end{example}

\begin{example}
Consider the Mori--Mukai model $MM_{4-9}$ on the algebraic torus
$(\mathbb{C}^*)^3$, with Laurent potential
\[
W_{MM_{4-9}}=xy+x+y+z+\frac{1}{y}+\frac{2}{x}+\frac{1}{x^2z}.
\]
The $z$-derivative gives
\[
1-\frac{1}{x^2z^2}=0,
\qquad\text{hence}\qquad
x^2z^2=1.
\]
Thus the critical locus splits into two possible branches
\[
z=\pm\frac{1}{x}.
\]
The $y$-derivative gives
\[
x+1-\frac{1}{y^2}=0,
\qquad\text{hence}\qquad
y^2=\frac{1}{x+1}.
\]
The remaining $x$-derivative is
\[
y+1-\frac{2}{x^2}-\frac{2}{x^3z}=0.
\]

On the negative branch $z=-1/x$, this equation reduces to
\[
y+1=0.
\]
Hence $y=-1$. Substituting into $y^2=1/(x+1)$ gives $x=0$, which is
forbidden on the torus. Therefore the negative branch is empty.

On the positive branch $z=1/x$, the $x$-equation becomes
\[
y+1-\frac{4}{x^2}=0,
\]
so
\[
y=\frac{4-x^2}{x^2}.
\]
Combining this with $y^2=1/(x+1)$ gives
\[
\left(\frac{4-x^2}{x^2}\right)^2=\frac{1}{x+1}.
\]
Equivalently,
\[
\Phi(x):=x^5-8x^3-8x^2+16x+16=0.
\]
This is the same quintic that governs the $x$-coordinates in the model
$MM_{4-4}$. The polynomial $\Phi$ is irreducible and square-free over
$\mathbb{Q}$, and is coprime to $x(x+1)(4-x^2)$. Hence all five roots give
torus points
\[
P_\lambda=\left(
\lambda,
\frac{4-\lambda^2}{\lambda^2},
\frac{1}{\lambda}
\right),
\qquad
\Phi(\lambda)=0.
\]
Equivalently, clearing denominators in the logarithmic critical equations and
saturating by $xyz$ gives the basis
\[
-x^4+x^3+8x^2+4y-20,
\]
\[
x^4-8x^2-8x+16z+16,
\qquad
\Phi(x).
\]

The critical values are the roots of
\[
w^5+6w^4-36w^3-360w^2-704w+272=0.
\]
Equivalently, modulo $\Phi(\lambda)$ one has
\[
W_{MM_{4-9}}(P_\lambda)=
-\frac{1}{4}(\lambda^4+2\lambda^3-8\lambda^2-20\lambda+8).
\]
The logarithmic Hessian is nondegenerate at every critical point. More
precisely,
\[
\det(\theta_i\theta_jW)(P_\lambda)
=2(\lambda^4-2\lambda^3-7\lambda^2+3\lambda+24),
\]
and the displayed factor is coprime to $\Phi$. Therefore all five critical
points are isolated Morse singularities of type $A_1$. Consequently,
\[
\overline{MF((\mathbb{C}^*)^3,W_{MM_{4-9}})}
\cong \bigoplus_{\lambda:\,\Phi(\lambda)=0}\MorseCat{3}.
\]
\end{example}

\begin{example}
Consider the Mori--Mukai model $MM_{4-10}$ on the algebraic torus
$(\mathbb{C}^*)^3$, with Laurent potential
\[
W_{MM_{4-10}}=xy+x+y+z+\frac{1}{y}+\frac{1}{yz}+\frac{1}{x}.
\]
The $z$-derivative gives
\[
1-\frac{1}{yz^2}=0,
\qquad\text{hence}\qquad
y=\frac{1}{z^2}.
\]
The $x$-derivative gives
\[
y+1-\frac{1}{x^2}=0,
\qquad\text{hence}\qquad
x^2=\frac{1}{y+1}.
\]
The $y$-derivative gives
\[
x+1-\frac{1}{y^2}-\frac{1}{y^2z}=0,
\qquad\text{hence}\qquad
x+1=\frac{1}{y^2}\left(1+\frac{1}{z}\right).
\]
Substituting $y=z^{-2}$ into the last equation yields
\[
x+1=z^4\left(\frac{z+1}{z}\right)=z^4+z^3,
\]
so
\[
x=z^4+z^3-1.
\]
Thus the critical equations reduce to a single equation in $z$. Imposing
$x^2=1/(y+1)$ gives
\[
(z^4+z^3-1)^2=\frac{z^2}{z^2+1},
\]
or equivalently
\[
\Psi(z):=z^{10}+2z^9+2z^8+2z^7-z^6-2z^5
-2z^4-2z^3+1=0.
\]
The polynomial $\Psi$ is irreducible and square-free over $\mathbb{Q}$, and is
coprime to $z(z^2+1)(z^4+z^3-1)$. Hence all ten roots give torus points
\[
P_\gamma=\left(
\gamma^4+\gamma^3-1,
\frac{1}{\gamma^2},
\gamma
\right),
\qquad
\Psi(\gamma)=0.
\]
Equivalently, clearing denominators in the logarithmic critical equations and
saturating by $xyz$ gives the lexicographic basis
\[
x-z^4-z^3+1,
\]
\[
y+z^8+2z^7+2z^6+2z^5-z^4-2z^3-2z^2-2z,
\]
\[
\Psi(z).
\]

The critical values are the roots of
\[
\begin{aligned}
0={}&w^{10}+8w^9-4w^8-314w^7-1768w^6-4310w^5 \\
&-1815w^4+17052w^3+46298w^2+34800w-23887.
\end{aligned}
\]
Equivalently, modulo $\Psi(\gamma)$ one has
\[
\begin{aligned}
W_{MM_{4-10}}(P_\gamma)={}&
\gamma^8+2\gamma^7+2\gamma^6+2\gamma^5+\gamma^4 \\
&+\gamma^2+2\gamma-2.
\end{aligned}
\]
The logarithmic Hessian is nondegenerate at every critical point. More
precisely,
\[
\det(\theta_i\theta_jW)(P_\gamma)
=2(5\gamma^7+9\gamma^6+8\gamma^5+7\gamma^4
-3\gamma^3-5\gamma^2-4\gamma-3),
\]
and the displayed factor is coprime to $\Psi$. Therefore all ten critical
points are isolated Morse singularities of type $A_1$. Consequently,
\[
\overline{MF((\mathbb{C}^*)^3,W_{MM_{4-10}})}
\cong \bigoplus_{\gamma:\,\Psi(\gamma)=0}\MorseCat{3}.
\]
\end{example}

\begin{example}
Consider the Mori--Mukai model $MM_{4-11}$ on the algebraic torus
$(\mathbb{C}^*)^3$, with Laurent potential
\[
W_{MM_{4-11}}=x+y+z+\frac{1}{y}+\frac{z}{x}+\frac{2}{x}+\frac{1}{xz}.
\]
The three terms with denominator $x$ form the square
\[
\frac{z}{x}+\frac{2}{x}+\frac{1}{xz}
=\frac{1}{x}\left(z+2+\frac{1}{z}\right)
=\frac{(z+1)^2}{xz}.
\]
Thus
\[
W_{MM_{4-11}}=x+z+\frac{(z+1)^2}{xz}+y+\frac{1}{y},
\]
so the $y$-coordinate is decoupled from the $(x,z)$-subsystem. The
$y$-derivative gives
\[
1-\frac{1}{y^2}=0,
\qquad\text{hence}\qquad
y=\pm1.
\]
The $x$-derivative gives
\[
1-\frac{(z+1)^2}{x^2z}=0,
\qquad\text{hence}\qquad
x^2=\frac{(z+1)^2}{z}.
\]
The $z$-derivative gives
\[
1+\frac{1}{x}-\frac{1}{xz^2}=0,
\]
or, equivalently,
\[
z^2(x+1)=1,
\qquad\text{so}\qquad
x=\frac{1-z^2}{z^2}.
\]
Substituting this expression for $x$ into $x^2=(z+1)^2/z$ gives
\[
\left(\frac{1-z^2}{z^2}\right)^2=\frac{(z+1)^2}{z}.
\]
The factor $(z+1)^2$ corresponds to $z=-1$, which would force $x=0$ in
$x=(1-z^2)/z^2$ and is therefore removed by the torus condition. After this
saturation one obtains
\[
(1-z)^2=z^3,
\]
or
\[
p(z):=z^3-z^2+2z-1=0.
\]
The polynomial $p$ is irreducible and square-free over $\mathbb{Q}$, and is
coprime to $z(z^2-1)$. Hence its three roots give torus points on each of the
two sheets $y=\pm1$. The critical points are
\[
P^{\varepsilon}_\alpha=\left(
\frac{1-\alpha^2}{\alpha^2},
\varepsilon,
\alpha
\right),
\qquad
\varepsilon\in\{\pm1\},
\qquad
p(\alpha)=0.
\]
Equivalently, clearing denominators in the logarithmic critical equations and
saturating by $xyz$ gives the lexicographic basis
\[
x-2z^2+z-2,
\qquad
(y-1)(y+1),
\qquad
p(z).
\]

On the sheet $y=1$, the critical values are the roots of
\[
w^3-5w^2-10w-11=0,
\]
and, modulo $p(\alpha)$,
\[
W_{MM_{4-11}}(P^+_\alpha)=4\alpha^2-\alpha+6.
\]
On the sheet $y=-1$, the critical values are the roots of
\[
w^3+7w^2-2w-67=0,
\]
and
\[
W_{MM_{4-11}}(P^-_\alpha)=4\alpha^2-\alpha+2.
\]
The logarithmic Hessian is nondegenerate at every critical point. More
precisely,
\[
\det(\theta_i\theta_jW)(P^{\varepsilon}_\alpha)
=2\varepsilon(3\alpha^2-4\alpha+8),
\]
and the displayed factor is coprime to $p$. Therefore all six critical points
are isolated Morse singularities of type $A_1$. Consequently,
\[
\overline{MF((\mathbb{C}^*)^3,W_{MM_{4-11}})}
\cong \bigoplus_{\varepsilon\in\{\pm1\}} \bigoplus_{\alpha:\,p(\alpha)=0}\MorseCat{3}.
\]
\end{example}

\begin{example}
Consider the Mori--Mukai model $MM_{4-12}$ on the algebraic torus
$(\mathbb{C}^*)^3$, with Laurent potential
\[
W_{MM_{4-12}}=xy+x+\frac{x}{z}+y+z+\frac{1}{y}+\frac{1}{x}.
\]
The $z$-derivative gives
\[
-\frac{x}{z^2}+1=0,
\qquad\text{hence}\qquad
x=z^2.
\]
The $y$-derivative gives
\[
x+1-\frac{1}{y^2}=0,
\qquad\text{hence}\qquad
y^2=\frac{1}{x+1}.
\]
The $x$-derivative gives
\[
y+1+\frac{1}{z}-\frac{1}{x^2}=0,
\qquad\text{hence}\qquad
y=\frac{1}{x^2}-1-\frac{1}{z}.
\]
Substituting $x=z^2$ gives
\[
y^2=\frac{1}{z^2+1},
\qquad
y=\frac{1-z^3-z^4}{z^4}.
\]
Thus the critical equations reduce to a single equation in $z$. Equating these
two expressions for $y^2$ gives
\[
\left(\frac{1-z^3-z^4}{z^4}\right)^2=\frac{1}{z^2+1},
\]
or equivalently
\[
\Xi(z):=z^{10}+2z^9+z^8+2z^7-z^6-2z^5
-2z^4-2z^3+z^2+1=0.
\]
The polynomial $\Xi$ is irreducible and square-free over $\mathbb{Q}$, and is
coprime to $z(z^2+1)(1-z^3-z^4)$. Hence all ten roots give torus points
\[
P_\delta=\left(
\delta^2,
\frac{1-\delta^3-\delta^4}{\delta^4},
\delta
\right),
\qquad
\Xi(\delta)=0.
\]
Equivalently, clearing denominators in the logarithmic critical equations and
saturating by $xyz$ gives the lexicographic basis
\[
x-z^2,
\]
\[
y+z^9+z^8-z^7+2z^6-z^5+2z^3-z^2+z-2,
\]
\[
\Xi(z).
\]

The critical values are the roots of
\[
\begin{aligned}
0={}&w^{10}+8w^9-10w^8-344w^7-1339w^6-386w^5 \\
&+9156w^4+23694w^3+16220w^2-21120w-44407.
\end{aligned}
\]
Equivalently, modulo $\Xi(\delta)$ one has
\[
\begin{aligned}
W_{MM_{4-12}}(P_\delta)={}&
-2\delta^9-5\delta^8-4\delta^7-7\delta^6-4\delta^5 \\
&+3\delta^4+2\delta^3+7\delta^2+4\delta+1.
\end{aligned}
\]
The logarithmic Hessian is nondegenerate at every critical point. More
precisely,
\[
\det(\theta_i\theta_jW)(P_\delta)
=-2\delta(\delta+2)(5\delta^7+3\delta^6+4\delta^5+3\delta^4
-\delta^3-4\delta^2-3\delta-1),
\]
and the displayed factor is coprime to $\Xi$. Therefore all ten critical
points are isolated Morse singularities of type $A_1$. Consequently,
\[
\overline{MF((\mathbb{C}^*)^3,W_{MM_{4-12}})}
\cong \bigoplus_{\delta:\,\Xi(\delta)=0}\MorseCat{3}.
\]
\end{example}

\begin{example}
Consider the Mori--Mukai model $MM_{4-13}$ on the algebraic torus
$(\mathbb C^*)^3$:
\[
W_{MM_{4-13}}
=xy+\frac{xy}{z}+x+y+z+\frac1y+\frac1x.
\]
This is the Laurent polynomial listed for $MM_{4-13}$ in Appendix~A of
\cite{CoatesCortiGalkinKasprzyk2016}.  After clearing Laurent denominators,
the logarithmic critical equations are
\[
x^2yz+x^2y+x^2z-z=0,\qquad
xy^2z+xy^2+y^2z-z=0,\qquad
-xy+z^2=0.
\]
Their saturated lexicographic Gr\"obner basis may be taken to contain
\[
\begin{gathered}
x-yz^2-z^4-z^3,\qquad y^2+yz^2+yz-1,\\
(z+1)(z^2+1)(y-z^3+z-1),\\
(z+1)(z^2+1)(z^2-z+1)(z^2+z-1).
\end{gathered}
\]
Thus the critical locus consists of ten points. Six are indexed by
\[
z=-1\ \text{or}\ z^2+1=0,\qquad
y^2+yz^2+yz-1=0,\qquad x=\frac{z^2}{y},
\]
and the remaining four are indexed by
\[
(z^2-z+1)(z^2+z-1)=0,\qquad
y=z^3-z+1,\qquad x=\frac{z^2}{y}.
\]
Their critical values are the roots, with multiplicity one, of
\[
(w-3)(w+5)(w^2-w-31)(w^2+2w+5)(w^2+3w+9).
\]
To check nondegeneracy, reduce the determinant of the logarithmic Hessian on
the three branches above. The respective reductions are
\[
4(y+2),\qquad -4(3yz-y+z-2),\qquad 2z(4z^2-z+1).
\]
Each is coprime to the corresponding branch equations. Hence all ten critical
points are Morse and
\[
\overline{MF((\mathbb C^*)^3,W_{MM_{4-13}})}
\cong
\bigoplus_{p\in\CritLoc(W_{MM_{4-13}})}\MorseCat{3}.
\]
\end{example}

\begin{example}
Consider the Mori--Mukai model $MM_{5-1}$ on the algebraic torus
$(\mathbb{C}^*)^3$, with Laurent potential
\[
W_{MM_{5-1}}=x+y+\frac{y}{z}+z+\frac{2}{y}
+\frac{y}{x}+\frac{3}{x}+\frac{3}{xy}+\frac{1}{xy^2}.
\]
The $z$-derivative gives
\[
-\frac{y}{z^2}+1=0,
\qquad\text{hence}\qquad
y=z^2.
\]
The $x$-derivative gives
\[
1-\frac{y}{x^2}-\frac{3}{x^2}-\frac{3}{x^2y}
-\frac{1}{x^2y^2}=0,
\]
or equivalently
\[
x^2y^2=(y+1)^3.
\]
The $y$-derivative gives
\[
1+\frac{1}{z}-\frac{2}{y^2}+\frac{1}{x}
-\frac{3}{xy^2}-\frac{2}{xy^3}=0.
\]
After multiplying by $xy^3$ and using $y=z^2$, this becomes
\[
xy(y^2+yz-2)+(y^3-3y-2)=0.
\]
Since
\[
y^3-3y-2=(y+1)^2(y-2),
\]
we obtain
\[
xy(y^2+yz-2)=-(y+1)^2(y-2).
\]
Squaring this equation and using $x^2y^2=(y+1)^3$ gives
\[
(y^2+yz-2)^2=(y+1)(y-2)^2.
\]
Here $y+1$ is nonzero, since $x^2y^2=(y+1)^3$ and $x,y\neq0$ on the torus.
Substituting $y=z^2$ yields
\[
(z^4+z^3-2)^2=(z^2+1)(z^2-2)^2,
\]
and hence, after removing the torus factor $z^3$,
\[
p(z):=z^5+2z^4-z-4=0.
\]
The polynomial $p$ is irreducible and square-free over $\mathbb{Q}$, and is
coprime to
\[
z(z^2+1)(z^2-2)(z^4+z^3-2).
\]
Thus all five roots give torus points. The critical points are
\[
P_\alpha=\left(
-\frac{(\alpha^2+1)^2(\alpha^2-2)}{\alpha^2(\alpha^4+\alpha^3-2)},
\alpha^2,
\alpha
\right),
\qquad
p(\alpha)=0.
\]
Equivalently, clearing denominators in the logarithmic critical equations and
saturating by $xyz$ gives the lexicographic basis
\[
16x+23z^4+2z^3-24z^2+16z-71,
\qquad
y-z^2,
\qquad
p(z).
\]

The critical values are the roots of
\[
w^5-5w^4-74w^3+14w^2+1333w+3499=0.
\]
Equivalently, modulo $p(\alpha)$ one has
\[
W_{MM_{5-1}}(P_\alpha)
=2x(P_\alpha)+\alpha^2+2\alpha+\frac{2}{\alpha^2}
=-3\alpha^4+5\alpha^2+9.
\]
The logarithmic Hessian is nondegenerate at every critical point. More
precisely,
\[
\det(\theta_i\theta_jW)(P_\alpha)
=2(16\alpha^4+12\alpha^3-3\alpha^2+22\alpha-31),
\]
and the displayed factor is coprime to $p$. Therefore all five critical points
are isolated Morse singularities of type $A_1$. Consequently,
\[
\overline{MF((\mathbb{C}^*)^3,W_{MM_{5-1}})}
\cong \bigoplus_{\alpha:\,p(\alpha)=0}\MorseCat{3}.
\]
\end{example}

\begin{example}
Consider the Mori--Mukai model $MM_{5-2}$ on the algebraic torus
$(\mathbb{C}^*)^3$, with Laurent potential
\[
W_{MM_{5-2}}=x+y+\frac{y}{z}+z+\frac{1}{y}
+\frac{y}{x}+\frac{2}{x}+\frac{1}{xy}.
\]
The $z$-derivative gives
\[
-\frac{y}{z^2}+1=0,
\qquad\text{hence}\qquad
y=z^2.
\]
The $x$-derivative gives
\[
1-\frac{y}{x^2}-\frac{2}{x^2}-\frac{1}{x^2y}=0,
\]
or equivalently
\[
x^2y=(y+1)^2.
\]
The $y$-derivative gives
\[
1+\frac{1}{z}-\frac{1}{y^2}+\frac{1}{x}-\frac{1}{xy^2}=0.
\]
Using $y=z^2$, this equation can be written as
\[
(z^4-1)(x+1)+xz^3=0,
\]
so that
\[
x=-\frac{z^4-1}{z^4+z^3-1}.
\]
On the other hand, the $x$-equation gives
\[
x^2=\frac{(z^2+1)^2}{z^2}.
\]
Equating these two expressions for $x^2$ and cancelling the nonzero factor
$(z^2+1)^2$ gives
\[
z^2(z^2-1)^2=(z^4+z^3-1)^2.
\]
Thus the critical locus splits into two branches:
\[
p_-(z):=z^4+z-1=0,
\qquad
p_+(z):=z^4+2z^3-z-1=0.
\]
The polynomials $p_-$ and $p_+$ are irreducible and square-free over
$\mathbb{Q}$, are coprime to each other, and are coprime to
$z(z^2+1)(z^4+z^3-1)$. Hence all eight roots give torus points. The critical
points are
\[
P^-_\alpha=\left(
-\frac{\alpha^2+1}{\alpha},
\alpha^2,
\alpha
\right),
\qquad
p_-(\alpha)=0,
\]
and
\[
P^+_\beta=\left(
\frac{\beta^2+1}{\beta},
\beta^2,
\beta
\right),
\qquad
p_+(\beta)=0.
\]
Equivalently, after clearing denominators and saturating by the torus
variables, the two branch lexicographic bases are
\[
x+z^3+z+1,
\qquad
y-z^2,
\qquad
p_-(z),
\]
and
\[
x-z^3-2z^2-z+1,
\qquad
y-z^2,
\qquad
p_+(z).
\]

On the negative branch, the critical values are the roots of
\[
w^4+w^3-8w^2-36w-11=0,
\]
and, modulo $p_-(\alpha)$,
\[
W_{MM_{5-2}}(P^-_\alpha)=-\alpha^3+2\alpha^2-1.
\]
On the positive branch, the critical values are the roots of
\[
w^4+5w^3-44w^2-360w-691=0,
\]
and
\[
W_{MM_{5-2}}(P^+_\beta)=\beta^3+4\beta^2+6\beta-1.
\]
The logarithmic Hessian is nondegenerate at every critical point. More
precisely,
\[
\det(\theta_i\theta_jW)(P^-_\alpha)
=-2(\alpha^3+4\alpha^2-3\alpha+5),
\]
and
\[
\det(\theta_i\theta_jW)(P^+_\beta)
=-2(3\beta^3-2\beta^2-9\beta-5),
\]
with the displayed factors coprime to $p_-$ and $p_+$, respectively.
Therefore all eight critical points are isolated Morse singularities of type
$A_1$. Consequently,
\[
\overline{MF((\mathbb{C}^*)^3,W_{MM_{5-2}})}
\cong \left(\bigoplus_{\alpha:\,p_-(\alpha)=0}\MorseCat{3}\right) \oplus \left(\bigoplus_{\beta:\,p_+(\beta)=0}\MorseCat{3}\right).
\]
\end{example}

\begin{example}
Consider the Mori--Mukai model $MM_{5-3}$ on the algebraic torus
$(\mathbb{C}^*)^3$, with Laurent potential
\[
W_{MM_{5-3}}=x+y+z+\frac{1}{y}+\frac{3}{x}
+\frac{2}{xz}+\frac{3}{x^2z}+\frac{1}{x^3z^2}.
\]
The $y$-coordinate is decoupled from the $(x,z)$-subsystem. Its derivative
gives
\[
1-\frac{1}{y^2}=0,
\qquad\text{hence}\qquad
y=\pm1.
\]
The $z$-derivative gives
\[
1-\frac{2}{xz^2}-\frac{3}{x^2z^2}-\frac{2}{x^3z^3}=0,
\]
or, after clearing denominators,
\[
x^3z^3-2x^2z-3xz-2=0.
\]
The $x$-derivative gives
\[
1-\frac{3}{x^2}-\frac{2}{x^2z}-\frac{6}{x^3z}
-\frac{3}{x^4z^2}=0,
\]
or equivalently
\[
x^4z^2-3x^2z^2-2x^2z-6xz-3=0.
\]
Set
\[
v=xz.
\]
The $z$-equation becomes
\[
v^3-2xv-3v-2=0,
\qquad\text{so}\qquad
2xv=(v+1)^2(v-2).
\]
The $x$-equation becomes
\[
(xv)^2-2xv-3(v+1)^2=0.
\]
Substituting $xv=(v+1)^2(v-2)/2$ gives
\[
\frac{v^2(v-3)(v+1)^3}{4}=0.
\]
The torus condition excludes $v=0$, and $v=-1$ would force $x=0$ in the
equation $2xv=(v+1)^2(v-2)$. Hence
\[
v=3.
\]
It follows that
\[
x=\frac{8}{3},
\qquad
z=\frac{9}{8}.
\]
Together with $y=\pm1$, this gives the two critical points
\[
P_{\varepsilon}=\left(\frac{8}{3},\varepsilon,\frac{9}{8}\right),
\qquad
\varepsilon\in\{\pm1\}.
\]
Equivalently, clearing denominators in the logarithmic critical equations and
saturating by $xyz$ gives the lexicographic basis
\[
3x-8,
\qquad
(y-1)(y+1),
\qquad
8z-9.
\]

The critical values are
\[
W_{MM_{5-3}}(P_{\varepsilon})=6+2\varepsilon.
\]
Thus
\[
W_{MM_{5-3}}(P_1)=8,
\qquad
W_{MM_{5-3}}(P_{-1})=4.
\]
The logarithmic Hessian is nondegenerate at both critical points. More
precisely,
\[
(\theta_i\theta_jW)(P_{\varepsilon})=
\begin{pmatrix}
\frac{19}{3} & 0 & \frac{5}{3} \\
0 & 2\varepsilon & 0 \\
\frac{5}{3} & 0 & \frac{7}{3}
\end{pmatrix},
\qquad
\det(\theta_i\theta_jW)(P_{\varepsilon})=24\varepsilon.
\]
Therefore both critical points are isolated Morse singularities of type
$A_1$. Consequently,
\[
\overline{MF((\mathbb{C}^*)^3,W_{MM_{5-3}})}
\cong \bigoplus_{\varepsilon\in\{\pm1\}}\MorseCat{3}.
\]
\end{example}

\begin{example}
Consider the Mori--Mukai model $MM_{6-1}$ on the algebraic torus
$(\mathbb{C}^*)^3$, with Laurent potential
\[
W_{MM_{6-1}}=x+y+z+\frac{1}{z}+\frac{2}{y}
+\frac{y}{x}+\frac{3}{x}+\frac{3}{xy}+\frac{1}{xy^2}.
\]
The $z$-coordinate is decoupled and gives
\[
1-\frac{1}{z^2}=0,
\qquad\text{hence}\qquad
z=\pm1.
\]
The $x$-derivative gives
\[
1-\frac{y}{x^2}-\frac{3}{x^2}-\frac{3}{x^2y}
-\frac{1}{x^2y^2}=0,
\]
or equivalently
\[
x^2y^2=(y+1)^3.
\]
The $y$-derivative gives
\[
1-\frac{2}{y^2}+\frac{1}{x}-\frac{3}{xy^2}-\frac{2}{xy^3}=0.
\]
Using
\[
y^3-3y-2=(y+1)^2(y-2),
\]
this equation may be written as
\[
\frac{y^2-2}{y^2}+\frac{(y+1)^2(y-2)}{xy^3}=0.
\]
Thus
\[
\frac{1}{x}=\frac{y(2-y^2)}{(y+1)^2(y-2)}.
\]
Squaring this relation and comparing it with
$1/x^2=y^2/(y+1)^3$ gives
\[
(y+1)(y-2)^2=(y^2-2)^2.
\]
After expansion this reduces to
\[
y^2(y^2-y-1)=0.
\]
Since $y$ is invertible on the torus, we obtain
\[
q(y):=y^2-y-1=0.
\]
Modulo $q(y)$, the preceding relation for $1/x$ gives $x=y+1$. Hence the
critical points are
\[
P_{\alpha,\varepsilon}=(\alpha+1,\alpha,\varepsilon),
\qquad
q(\alpha)=0,
\qquad
\varepsilon\in\{\pm1\}.
\]
Equivalently, clearing denominators in the logarithmic critical equations and
saturating by $xyz$ gives the lexicographic basis
\[
x-y-1,
\qquad
y^2-y-1,
\qquad
(z-1)(z+1).
\]

The critical values are
\[
W_{MM_{6-1}}(P_{\alpha,\varepsilon})=5\alpha+2\varepsilon.
\]
Equivalently, on the sheet $z=1$ they are the roots of
\[
w^2-9w-11=0,
\]
while on the sheet $z=-1$ they are the roots of
\[
w^2-w-31=0.
\]
The logarithmic Hessian is nondegenerate at every critical point. More
precisely,
\[
(\theta_i\theta_jW)(P_{\alpha,\varepsilon})=
\begin{pmatrix}
2(\alpha+1) & 2-\alpha & 0 \\
2-\alpha & 8-2\alpha & 0 \\
0 & 0 & 2\varepsilon
\end{pmatrix},
\]
and
\[
\det(\theta_i\theta_jW)(P_{\alpha,\varepsilon})
=2\varepsilon(11\alpha+7).
\]
The displayed factor is coprime to $q$. Therefore all four critical points are
isolated Morse singularities of type $A_1$. Consequently,
\[
\overline{MF((\mathbb{C}^*)^3,W_{MM_{6-1}})}
\cong \bigoplus_{\varepsilon\in\{\pm1\}} \bigoplus_{\alpha:\,q(\alpha)=0}\MorseCat{3}.
\]
\end{example}


\begin{example}
Consider the Mori--Mukai model $MM_{7-1}$ on the algebraic torus
$(\mathbb{C}^*)^3$, with Laurent potential
\[
W_{MM_{7-1}}=xz^2+2xz+x+y+4z+\frac{1}{y}
+\frac{6}{x}+\frac{2}{xz}+\frac{4}{x^2z}+\frac{1}{x^3z^2}.
\]
The $y$-coordinate is decoupled from the $(x,z)$-subsystem. Its derivative
therefore gives
\[
1-\frac{1}{y^2}=0,
\qquad\text{hence}\qquad
 y=\pm1.
\]
The $z$-derivative gives
\[
2xz+2x+4-\frac{2}{xz^2}-\frac{4}{x^2z^2}-\frac{2}{x^3z^3}=0.
\]
After multiplying by $x^3z^3/2$, this factors as
\[
(xz-1)(xz+1)(x^2z^2+x^2z+2xz+1)=0.
\]
The $x$-derivative gives
\[
z^2+2z+1-\frac{6}{x^2}-\frac{2}{x^2z}-\frac{8}{x^3z}
-\frac{3}{x^4z^2}=0,
\]
or equivalently
\[
(x^2z^2+x^2z-2xz-3)(x^2z^2+x^2z+2xz+1)=0.
\]
Thus the common critical locus in the $(x,z)$-variables has two parts.
First, the equations $xz=1$ and
$x^2z^2+x^2z-2xz-3=0$ give
\[
 x=4,
\qquad
 z=\frac{1}{4}.
\]
The branch $xz=-1$ would force $x=0$ and is excluded by the torus condition.
Hence this gives two isolated critical points
\[
Q_{\varepsilon}=\left(4,\varepsilon,\frac{1}{4}\right),
\qquad
\varepsilon\in\{\pm1\}.
\]
At these points
\[
W_{MM_{7-1}}(Q_1)=14,
\qquad
W_{MM_{7-1}}(Q_{-1})=10.
\]

Second, the common factor
\[
C(x,z):=x^2z^2+x^2z+2xz+1
\]
defines a one-dimensional critical component. For each
$\varepsilon\in\{\pm1\}$ set
\[
\Gamma_{\varepsilon}=\{y=\varepsilon,\ C(x,z)=0\}
\subset (\mathbb{C}^*)^3.
\]
Equivalently, writing $t=xz$, this curve is parametrized by
\[
\gamma_{\varepsilon}(t)=\left(
-\frac{(t+1)^2}{t},
\varepsilon,
-\frac{t^2}{(t+1)^2}
\right),
\qquad
 t\in\mathbb{C}^*,\quad t\neq -1.
\]
Along these curves the potential is constant:
\[
W_{MM_{7-1}}|_{\Gamma_1}=-2,
\qquad
W_{MM_{7-1}}|_{\Gamma_{-1}}=-6.
\]
In particular, the points $(-4,1,-1/4)$ and $(-4,-1,-1/4)$ lie on the two
curve components, rather than forming additional isolated components.

The transverse models are global. Indeed, the potential satisfies the exact
identity
\[
W_{MM_{7-1}}-y-\frac1y+4
=\frac{C(x,z)^2}{x^3z^2}.
\]
For $\varepsilon\in\{\pm1\}$, put
\[
c_{\varepsilon}=-4+2\varepsilon,
\qquad
D_{\varepsilon}=y-\varepsilon.
\]
Then
\[
W_{MM_{7-1}}-c_{\varepsilon}
=\frac{C(x,z)^2}{x^3z^2}+\frac{D_{\varepsilon}^2}{y}.
\]
The ideal $(C,D_{\varepsilon})$ defines $\Gamma_{\varepsilon}$. With
\[
A_{\Gamma}=\mathbb C[t^{\pm1},(t+1)^{-1}],
\]
the parametrization above gives
\[
\widehat X_{\Gamma_{\varepsilon}}
\simeq
\operatorname{Spf}\!\left(A_{\Gamma}[[C,D_{\varepsilon}]]\right).
\]
Along $\Gamma_{\varepsilon}$ one has
\[
\frac{1}{x^3z^2}=-\frac{1}{t(t+1)^2},
\qquad
\frac1y=\varepsilon.
\]
Thus the relative Hessian determinant is
\[
-\frac{4\varepsilon}{t(t+1)^2},
\]
a unit in $A_{\Gamma}$. Its square class is $[t]$, since constants over
$\mathbb C$ and $(t+1)^2$ are squares, and this class is nontrivial. The
relative formal Morse lemma therefore gives a global nondegenerate quadratic
bundle with this Clifford twist on each curve.

Equivalently, clearing denominators in the logarithmic critical equations and
saturating by $xyz$ gives the reduced critical locus determined by
\[
(x-4)C(x,z),
\qquad
(4z-1)C(x,z),
\qquad
(y-1)(y+1).
\]
The two curve components are Morse--Bott of transverse type $A_1$. Indeed,
under the parametrization above the logarithmic Hessian has determinant zero,
while the minor with rows and columns $(x,y)$ is
\[
\det(\theta_i\theta_jW)_{i,j\in\{x,y\}}
=-\frac{16\varepsilon}{t},
\]
which is a unit on $\Gamma_{\varepsilon}$. At the isolated points,
\[
(\theta_i\theta_jW)(Q_{\varepsilon})=
\begin{pmatrix}
16 & 0 & 8 \\
0 & 2\varepsilon & 0 \\
8 & 0 & 8
\end{pmatrix},
\qquad
\det(\theta_i\theta_jW)(Q_{\varepsilon})=128\varepsilon.
\]
Therefore $Q_1$ and $Q_{-1}$ are isolated Morse singularities of type $A_1$.
By Proposition~\ref{prop:morse_bott_clifford}, the two curves contribute
$\MBtwcat{\Gamma_1}{t}$ and $\MBtwcat{\Gamma_{-1}}{t}$. Consequently,
\[
\overline{MF((\mathbb{C}^*)^3,W_{MM_{7-1}})}
\cong \MBtwcat{\Gamma_1}{t}\oplus \MBtwcat{\Gamma_{-1}}{t}
\oplus \bigoplus_{\varepsilon\in\{\pm1\}}\MorseCat{3}.
\]
\end{example}

\begin{example}
Consider the Mori--Mukai model $MM_{8-1}$ on the algebraic torus
$(\mathbb{C}^*)^3$, with Laurent potential
\[
W_{MM_{8-1}}=x+\frac{1}{x}+y(z+1)^3
+3\left(z+\frac{1}{z}\right)
+\frac{3(z+1)}{yz^2}+\frac{1}{y^2z^3}.
\]
The $x$-coordinate is completely decoupled and gives
\[
\frac{\partial W_{MM_{8-1}}}{\partial x}=1-\frac{1}{x^2}=0,
\qquad\text{hence}\qquad
x=\pm1.
\]
For the remaining variables, set
\[
A=yz(z+1).
\]
The $y$-derivative factors as
\[
\frac{\partial W_{MM_{8-1}}}{\partial y}
=\frac{(A-2)(A+1)^2}{y^3z^3}.
\]
Moreover,
\[
\frac{1}{3}\frac{\partial W_{MM_{8-1}}}{\partial z}
=\frac{(yz^2-1)(A+1)^2}{y^2z^4}.
\]
Thus the reduced critical locus is the union of the branch $A=-1$ and the
solutions of $A=2$ and $yz^2=1$, together with $x=\pm1$.

The second branch gives
\[
yz^2=1,
\qquad
yz(z+1)=2.
\]
Dividing the second equation by the first gives $(z+1)/z=2$, so $z=1$ and
$y=1$. Hence this branch gives the two isolated critical points
\[
P_{\varepsilon}=(\varepsilon,1,1),
\qquad
\varepsilon\in\{\pm1\}.
\]
At these points,
\[
W_{MM_{8-1}}(P_{\varepsilon})=21+2\varepsilon,
\]
that is,
\[
W_{MM_{8-1}}(P_1)=23,
\qquad
W_{MM_{8-1}}(P_{-1})=19.
\]

The branch $A=-1$ gives two curve components
\[
\Gamma_{\varepsilon}
=\{x=\varepsilon,\ yz(z+1)+1=0\}
\subset (\mathbb{C}^*)^3,
\qquad
\varepsilon\in\{\pm1\}.
\]
Equivalently, they are parametrized by
\[
\gamma_{\varepsilon}(t)=\left(
\varepsilon,
-\frac{1}{t(t+1)},
t
\right),
\qquad
t\in\mathbb{C}^*,\quad t\neq -1.
\]
Along these curves the potential is constant. Indeed,
\[
W_{MM_{8-1}}-\left(x+\frac{1}{x}\right)+6
=\frac{(A+1)^3(z+1)^2}{A^2z},
\]
and therefore
\[
W_{MM_{8-1}}|_{\Gamma_{\varepsilon}}=2\varepsilon-6.
\]
Thus
\[
W_{MM_{8-1}}|_{\Gamma_1}=-4,
\qquad
W_{MM_{8-1}}|_{\Gamma_{-1}}=-8.
\]
Consequently, the singular fibers occur at the four critical values
\[
23,
\qquad
19,
\qquad
-4,
\qquad
-8.
\]

The isolated points are Morse singularities. More precisely, the logarithmic
Hessian at $P_{\varepsilon}$ is
\[
(\theta_i\theta_jW)(P_{\varepsilon})=
\begin{pmatrix}
2\varepsilon & 0 & 0 \\
0 & 18 & 27 \\
0 & 27 & 54
\end{pmatrix},
\qquad
\det(\theta_i\theta_jW)(P_{\varepsilon})=486\varepsilon.
\]
Hence both $P_1$ and $P_{-1}$ are isolated singularities of type $A_1$.

The curve components are not Morse--Bott quadratic components. Under the
parametrization above,
\[
(\theta_i\theta_jW)(\gamma_{\varepsilon}(t))=
\begin{pmatrix}
2\varepsilon & 0 & 0 \\
0 & 0 & 0 \\
0 & 0 & 0
\end{pmatrix}.
\]
However, putting $u=A+1$, the formal expression along $\Gamma_{\varepsilon}$
is
\[
W_{MM_{8-1}}-(2\varepsilon-6)
=\left(x+\frac{1}{x}-2\varepsilon\right)
+\frac{u^3(z+1)^2}{(u-1)^2z}.
\]
The first summand is a nondegenerate quadratic term in the $x$-direction, and
the second summand is a unit times $u^3$. Thus the residual transverse type
along each curve component is the mixed type $v^2+u^3$. The global
coordinates $x-\varepsilon$ and $u=A+1$ have Jacobian minor
\[
\det\frac{\partial(x-\varepsilon,u)}{\partial(x,y)}
=z(z+1),
\]
which restricts to the unit $t(t+1)$ on $\Gamma_\varepsilon$. They therefore
give the global splitting
\[
\widehat X_{\Gamma_\varepsilon}
\simeq
\operatorname{Spf}\!\left(
\mathbb C[t^{\pm1},(t+1)^{-1}][[x-\varepsilon,u]]
\right).
\]
Put $a=x-\varepsilon$ and $s=z^2(z+1)^2$, and set $\lambda=s^2$. In the
completed coordinate ring the shifted potential is
\[
G=\frac{a^2}{\varepsilon+a}
+\frac{(z+1)^2}{z(1-u)^2}u^3.
\]
The unit $\lambda/(\varepsilon+a)$ has a formal square root, and
\[
\lambda\frac{(z+1)^2}{z(1-u)^2}
=\bigl(z(z+1)^2\bigr)^3(1-u)^{-2}
\]
has a formal cube root. Formal rescaling therefore gives the contact identity
$\lambda G=v^2+w^3$. Thus both mixed models are globally untwisted and each
curve contributes $\QAper{\Gamma_\varepsilon}{2}$. This notation retains the
quadratic normal direction together with the residual $A_2$ direction; it is
not the residual category $\Aper{\Gamma_\varepsilon}{2}$ alone.
The formal singular locus of the zero fiber is contained in
$V((x-\varepsilon,u))$, verifying the hypothesis of
Theorem~\ref{thm:general_formal}.
Consequently,
\[
\overline{MF((\mathbb{C}^*)^3,W_{MM_{8-1}})}
\cong
\QAper{\Gamma_1}{2}
\oplus
\QAper{\Gamma_{-1}}{2}
\oplus
\bigoplus_{\varepsilon\in\{\pm1\}}\MorseCat{3}.
\]
\end{example}

\begin{example}[The Laurent mirror for $MM_{9-1}$]
Using the product construction for $S_2\times\mathbb P^1$, the polynomial in
Appendix~A can be regrouped as
\[
W_{MM_{9-1}}
=x\left(z+1+\frac{1}{xz}\right)^4+y+\frac{1}{y}-12.
\]
Its expanded form is
\[
\begin{aligned}
W_{MM_{9-1}}={}&xz^4+4xz^3+6xz^2+4xz+x+y
 +4z^2+12z+\frac{4}{z}+\frac{1}{y}\\
&+\frac{6}{x}+\frac{12}{xz}+\frac{6}{xz^2}
 +\frac{4}{x^2z^2}+\frac{4}{x^2z^3}+\frac{1}{x^3z^4}.
\end{aligned}
\]

Put
\[
F=z+1+\frac1{xz},
\qquad
q=\frac1{xz}.
\]
The logarithmic equations are
\[
\begin{aligned}
\theta_xW_{MM_{9-1}}&=xF^3(F-4q),\\
\theta_zW_{MM_{9-1}}&=4xF^3(z-q),\\
\theta_yW_{MM_{9-1}}&=y-\frac1y.
\end{aligned}
\]
Therefore $y=\varepsilon$, with $\varepsilon\in\{\pm1\}$.  The branch
$F=0$ gives two smooth critical curves
\[
\Gamma_\varepsilon=
\left\{
y=\varepsilon,\quad
x=-\frac1{z(z+1)}
\right\},
\qquad
z\in\mathbb C^*,\quad z\neq-1.
\]
Their critical values are
\[
W_{MM_{9-1}}|_{\Gamma_\varepsilon}=2\varepsilon-12.
\]
If $D_\varepsilon=y-\varepsilon$, then the exact identity
\[
W_{MM_{9-1}}-(2\varepsilon-12)
=xF^4+\frac{D_\varepsilon^2}{y}
\]
gives a global quadratic line together with a residual $A_3$ direction. Write
$t=z$ on $\Gamma_\varepsilon$. The functions $D_\varepsilon$ and $F$ are
global formal normal coordinates because
\[
\det\frac{\partial(z,D_\varepsilon,F)}{\partial(x,y,z)}
\bigg|_{\Gamma_\varepsilon}=t(t+1)^2
\]
is a unit. With
\[
x_0=-\frac1{t(t+1)},\qquad
v=D_\varepsilon(\varepsilon+D_\varepsilon)^{-1/2},\qquad
u=F\left(\frac{x}{x_0}\right)^{1/4},
\]
the completed shifted potential becomes $v^2+x_0u^4$. Taking
$h=it(t+1)$, contact multiplication by $h^2$ and the change $v\mapsto hv$
give
\[
v^2+t(t+1)u^4.
\]
The unit $t(t+1)$ is not a square in
$\mathbb C[t^{\pm1},(t+1)^{-1}]$, so its square class must be retained. Hence
each curve contributes $\QAtwist{\Gamma_\varepsilon}{3}{t(t+1)}$.

On the branch $F\neq0$, one obtains $z=q$ and $F=4q$, hence
\[
P_\varepsilon=\left(4,\varepsilon,\frac12\right).
\]
At these points
\[
W_{MM_{9-1}}(P_\varepsilon)=52+2\varepsilon,
\qquad
\det(\theta_i\theta_jW_{MM_{9-1}})(P_\varepsilon)
=4096\varepsilon.
\]
Thus both are Morse, and
\[
\overline{MF((\mathbb C^*)^3,W_{MM_{9-1}})}
\cong
\bigoplus_{\varepsilon\in\{\pm1\}}
\QAtwist{\Gamma_\varepsilon}{3}{t(t+1)}
\oplus
\bigoplus_{\varepsilon\in\{\pm1\}}\MorseCat{3}.
\]
\end{example}

\begin{example}[The Laurent mirror for $MM_{10-1}$]
For the product $S_1\times\mathbb P^1$, the potential is
\[
W_{MM_{10-1}}
=x\left(z+1+\frac{1}{xz^2}\right)^6+y+\frac{1}{y}-60.
\]
Equivalently,
\[
\begin{aligned}
W_{MM_{10-1}}={}&xz^6+6xz^5+15xz^4+20xz^3+15xz^2+6xz+x+y\\
&+6z^3+30z^2+60z+\frac{30}{z}+\frac{6}{z^2}+\frac{1}{y}
 +\frac{15}{x}+\frac{60}{xz}+\frac{90}{xz^2}\\
&+\frac{60}{xz^3}+\frac{15}{xz^4}
 +\frac{20}{x^2z^3}+\frac{60}{x^2z^4}
 +\frac{60}{x^2z^5}+\frac{20}{x^2z^6}\\
&+\frac{15}{x^3z^6}+\frac{30}{x^3z^7}+\frac{15}{x^3z^8}
 +\frac{6}{x^4z^9}+\frac{6}{x^4z^{10}}+\frac{1}{x^5z^{12}}.
\end{aligned}
\]

Set
\[
F=z+1+\frac1{xz^2},
\qquad
q=\frac1{xz^2}.
\]
Now
\[
\begin{aligned}
\theta_xW_{MM_{10-1}}&=xF^5(F-6q),\\
\theta_zW_{MM_{10-1}}&=6xF^5(z-2q),\\
\theta_yW_{MM_{10-1}}&=y-\frac1y.
\end{aligned}
\]
The branch $F=0$ again gives two smooth curves
\[
\Gamma_\varepsilon=
\left\{
y=\varepsilon,\quad
x=-\frac1{z^2(z+1)}
\right\},
\qquad
z\in\mathbb C^*,\quad z\neq-1,
\]
with critical values
\[
W_{MM_{10-1}}|_{\Gamma_\varepsilon}=2\varepsilon-60.
\]
Writing $D_\varepsilon=y-\varepsilon$, one has
\[
W_{MM_{10-1}}-(2\varepsilon-60)
=xF^6+\frac{D_\varepsilon^2}{y}.
\]
Writing $t=z$ on $\Gamma_\varepsilon$, the Jacobian determinant
\[
\det\frac{\partial(z,D_\varepsilon,F)}{\partial(x,y,z)}
\bigg|_{\Gamma_\varepsilon}=t^2(t+1)^2
\]
is a unit, so $D_\varepsilon$ and $F$ are global formal normal coordinates.
Put
\[
x_0=-\frac1{t^2(t+1)},\qquad
v=D_\varepsilon(\varepsilon+D_\varepsilon)^{-1/2},\qquad
u=F\left(\frac{x}{x_0}\right)^{1/6}.
\]
The completed shifted potential is $v^2+x_0u^6$. For $h=it(t+1)$,
contact multiplication by $h^2$ followed by $v\mapsto hv$ gives
\[
v^2+(t+1)u^6.
\]
Since $[t+1]$ is nontrivial in the unit square-class group of the curve, each
curve contributes $\QAtwist{\Gamma_\varepsilon}{5}{t+1}$.

On the complement of $F=0$, the equations give $z=2q$ and $F=6q$.
Consequently,
\[
P_\varepsilon=
\left(\frac{27}{4},\varepsilon,\frac23\right),
\]
and
\[
W_{MM_{10-1}}(P_\varepsilon)=372+2\varepsilon,
\qquad
\det(\theta_i\theta_jW_{MM_{10-1}})(P_\varepsilon)
=373248\varepsilon.
\]
Both isolated points are Morse.  Therefore
\[
\overline{MF((\mathbb C^*)^3,W_{MM_{10-1}})}
\cong
\bigoplus_{\varepsilon\in\{\pm1\}}
\QAtwist{\Gamma_\varepsilon}{5}{t+1}
\oplus
\bigoplus_{\varepsilon\in\{\pm1\}}\MorseCat{3}.
\]
\end{example}


\subsection*{Global formal-splitting and comparison audit}

For each positive-dimensional component used in a categorical identification,
the preceding computations give global generators of the conormal ideal and
an inverse formal substitution, or equivalently a Jacobian minor which is a
unit on the component. The following table records the resulting certificate.
The symbol $\delta(q)$ denotes the discriminant square class in
$\mathcal O_\Gamma^*/\mathcal O_\Gamma^{*2}$.

\begingroup
\scriptsize
\begin{longtblr}{
  width = \textwidth,
  colspec = {Q[l,wd=0.17\textwidth]Q[l,wd=0.31\textwidth]Q[l,wd=0.39\textwidth]},
  rowhead = 1,
  row{1} = {font=\bfseries},
  rowsep = 1.5pt,
  hline{1,2,Z} = {0.4pt},
}
Component & Global conormal/splitting certificate & Formal model and twist \\
$B_4$, $\Gamma$ & normals $(xy-1,z+1)$; minor $x$ & quadratic $t(U-V)^2+t^{-1}V^2$; hyperbolic, untwisted \\
$B_2$, $S$ & $U$; $\partial_xU=-2/(x^3yz)$ & $xU^2$; nontrivial class $[x]$ \\
$B_1$, $S^\circ$ and $K$ & $G$ on $S^\circ$; $(x+y,G)$ on $K$ & $AG^2$ with class $[A]$ off $K$; crossing $uv^2$ on $K$ \\
$V_{16}$, $\Gamma$ & $(y-1,z+1)$; minor $1$ & quadratic; $\delta=-4$, untwisted \\
$V_{14}$, $\Gamma$ & $(A,x-1)$; Jacobian minor $z(z+1)$ & relative $D_4$; discriminant class $[t(t+1)]$ \\
$V_{12}$, $\Gamma$ & $(y+1,z(x-1)^2+1)$; minor $(t-1)^2$ & quadratic; $\delta=-1$, untwisted \\
$V_{10}$, $S^\circ$ and $K$ & $F$ on $S^\circ$; $(F,G)$ on $K$, with Jacobian minor $-2xz^2(z+1)$ & $F^2$ off $K$; untwisted crossing $u^2v$ on $K$ \\
$V_8$, $S^\circ$ and $K$ & $U$ on $S^\circ$; $(A,B)$ on $K$ & $U^2$ off $K$; twisted crossing family on $K$ \\
$V_6$, $S_A^\circ$, $S_U^\circ$, and $K$ & $A$ on $S_A^\circ$; $U$ on $S_U^\circ$; $(A,U)$ on $K$ & $A^2$ and $U^3$ off $K$; crossing $a^2u^3$ on $K$ \\
$V_4$, $S$ & $U$; $\partial_xU=-1/(x^2yz)$ & $U^4$; residual $A_3$ \\
$V_2$, $S$ & $F$; $\partial_xF=-1/(x^2y^3z)$ & $W+120=xF^6$; residual $A_5$ \\
$MM_{2-1}$, $S^\circ$ and $K$ & $F$ on $S^\circ$; $(F,R+x+y)$ on $K$ & residual $A_2$ off $K$; crossing $u^3v$ on $K$ \\
$MM_{2-1}$, $\Gamma$ & $(R+x,R+y)$; exact product & hyperbolic quadratic, untwisted \\
$MM_{2-2}$, $S$ & $F$; $\partial_xF=-1/x$ on $S$ & $F^2Q$ with $Q|_S=1/(xy^2)$ \\
$MM_{2-3}$, $S^\circ$ and $K$ & $F$ on $S^\circ$; $(F,Q)$ on $K$ & quadratic off $K$; crossing $u^2v$ on $K$ \\
$MM_{2-3}$, $\Gamma$ & $(yz^2-1,H)$ & quadratic; square discriminant, untwisted \\
$MM_{2-5}$, $\Gamma$ & $(f_1,f_2)$; minor $t^2$ & quadratic; $\delta=-(t^3/D^3)^2$, untwisted \\
$MM_{2-6}$, $\Gamma_1$ & $(y-1,z+1)$; minor $1$ & quadratic; $\delta=[t]$, twisted \\
$MM_{2-6}$, $\Gamma_2$ & $(R,S)$; minor $-1$ & quadratic; $\delta=[-t^2/4]=1$ \\
$MM_{2-7}$, $\Gamma$ & $(F_1,F_2)$; minor $t^3(t+1)^2$ & quadratic; $\delta=[-D^{-6}]=1$ \\
$MM_{2-8}$, $\Sigma$ & $F$; $dF$ is nowhere zero on $\Sigma$ & $F^2/x$; untwisted \\
$MM_{2-12}$, $\Gamma$ & $(x-1,z+1)$; minor $1$ & quadratic; $\delta=[t]$, twisted \\
$MM_{2-13}$, $\Gamma$ & $(x-1,y+1)$; minor $1$ & quadratic; $\delta=[-5]=1$ \\
$MM_{2-15}$, $\Gamma$ & $(x+y+z,z+1)$; minor $1$ & quadratic; $\delta=[-1]=1$ \\
$MM_{2-19}$, $\Gamma$ & $(x+1,y-1)$; minor $1$ & quadratic; $\delta=[3]=1$ \\
$MM_{2-28}$, $\Gamma$ & $(F,G)$; their $2\times2$ minors generate $(1)$ on $\Gamma$ & exact $uv$; hyperbolic \\
$MM_{3-1}$, $\Gamma_3$ & $(u+1,y+1)$ in torus coordinates & quadratic; $\delta=[v]$, twisted \\
$MM_{3-1}$, $\Gamma_\zeta$ & $(p,q)$; unit triangular conormal change & exact $pq$; hyperbolic \\
$MM_{3-2}$, $\Gamma$ & $(r,s)$; unit triangular conormal change from $(H,F)$ & exact $rs$; hyperbolic \\
$MM_{4-1}$, $\Gamma$ & $(A,B)$; minor $t+1$ & quadratic; $\delta=[t(t+1)]$, twisted \\
$MM_{4-3}$, $\Gamma$ & $(A,B)$; minor $(t-1)/t$ & quadratic; $\delta=[1-t]$, twisted \\
$MM_{7-1}$, $\Gamma_\varepsilon$ & $(C,y-\varepsilon)$ with parameter $t=xz$; minor $-tx$ & quadratic; $\delta=[t]$, twisted \\
$MM_{8-1}$, $\Gamma_\varepsilon$ & $(x-\varepsilon,A+1)$; minor $t(t+1)$ & $v^2+u^3$; $\QAper{\Gamma_\varepsilon}{2}$ \\
$MM_{9-1}$, $\Gamma_\varepsilon$ & $(y-\varepsilon,F)$; minor $t(t+1)^2$ & $v^2+t(t+1)u^4$; $\QAtwist{\Gamma_\varepsilon}{3}{t(t+1)}$ \\
$MM_{10-1}$, $\Gamma_\varepsilon$ & $(y-\varepsilon,F)$; minor $t^2(t+1)^2$ & $v^2+(t+1)u^6$; $\QAtwist{\Gamma_\varepsilon}{5}{t+1}$ \\
\end{longtblr}
\endgroup

For every row with a nondegenerate quadratic model, the normal derivatives
generate the ideal $J$ of the zero section; hence
$\Sing(\operatorname{Spec}B)\subseteq V(J)$. For the residual models $u^n$,
the formal zero fiber is $u^n=0$, so the same inclusion follows from
$\partial_u(u^n)=nu^{n-1}$. For the mixed models $v^2+\delta u^{m+1}$,
the two normal derivatives again cut out the zero section; the notation
$\QAper{Z}{m}$ or $\QAtwist{Z}{m}{\delta}$ retains the quadratic direction
together with the residual one. The singular $V_8$ crossing stratum is retained
inside $\mathcal C_{V_8,-8}$ and is not replaced by an unrestricted split
model. This verifies, component by component, every use of
Theorem~\ref{thm:general_formal} in the table below.

\subsection*{Summary of Fano threefold decompositions}
The following table collects all Fano threefold Landau--Ginzburg models
computed in this section, including the seven exceptional families integrated
above in classification order.
Each entry records the corresponding decomposition of $\overline{MF((\mathbb{C}^*)^3,W)}$. The column
\emph{No.} records the number asserted by the computational certificate and should not be interpreted as the number of ordinary derived-category summands.
The $B_4$ quadratic curve is globally hyperbolic. The $B_2$ surface and the
open $B_1$ surface stratum retain the nontrivial classes $[x]$ and $[A]$,
respectively. The $V_8$ crossing stratum remains explicitly retained. For
$V_6$, the two residual surfaces are retained as one connected gluing
category because they meet along the $a^2u^3$ crossing curve. For $V_4$, the
displayed periodic $A$-type contribution comes from an explicit global formal
normal coordinate. The entries carrying
$\MBtwcat{Z}{\delta}$ have a verified split quadratic neighborhood but a
nontrivial discriminant square class $[\delta]$. The notation
$\QAper{Z}{m}$ records a standard mixed normal form $v^2+u^{m+1}$, while
$\QAtwist{Z}{m}{\delta}$ retains the nontrivial coefficient square class in
$v^2+\delta u^{m+1}$.
For $B_1$, $MM_{2-1}$, and $MM_{2-3}$, the symbols
$\mathcal C_{B_1,0}$, $\mathcal C_{MM_{2-1},-60}$, and
$\mathcal C_{MM_{2-3},-12}$ retain the crossing strata described in the
corresponding examples.  Their supported-localization sequences are not
asserted to split.

\begingroup
\footnotesize
\begin{longtblr}[
  caption = {Matrix factorization categorical decompositions for the Landau--Ginzburg models of Fano threefolds.},
]{
  width = \textwidth,
  colspec = {Q[l,wd=0.15\textwidth]Q[c,wd=0.06\textwidth]Q[l,wd=0.69\textwidth]},
  rowhead = 1,
  row{1} = {font=\bfseries},
  rowsep = 2pt,
  hline{1,2,Z} = {0.4pt},
}
Threefold & No. & Decomposition of $\overline{MF((\mathbb{C}^*)^3,W)}$ \\
\SetCell[c=3]{l}\textbf{Picard-rank-one models} & & \\
$\mathbb{P}^3$ & 4 & $\bigoplus_{n=1}^{4}\MorseCat{3}$ \\
$Q\subset\mathbb{P}^4$ & 3 & $\bigoplus_{n=1}^{3}\MorseCat{3}$ \\
$B_5$ & 4 & $\bigoplus_{n=1}^{4}\MorseCat{3}$ \\
$B_4$ & 3 & $\MBcat{\Gamma}\oplus\MorseCat{3}\oplus\MorseCat{3}$ \\
$B_3$ & 2 & $\bigoplus_{n=1}^{2}\MorseCat{3}$ \\
$B_2$ & 3 & $\left(\bigoplus_{n=1}^{2}\MorseCat{3}\right)\oplus \MBtwcat{S}{x}$ \\
$B_1$ & 3 & $\mathcal C_{B_1,0}\oplus\bigoplus_{a:\,48a^2-1=0}\MorseCat{3}$ \\
$V_{22}$ & 3 & $\bigoplus_{a:\,a^3+2a^2-2=0}\MorseCat{3}$ \\
$V_{18}$ & 2 & $\bigoplus_{a:\,36a^2-12a-11=0}\MorseCat{3}$ \\
$V_{16}$ & 3 & $\MBcat{\Gamma}\oplus\bigoplus_{a^2=2}\MorseCat{3}$ \\
$V_{14}$ & 3 & $\Dfourcat{\Gamma}{t(t+1)}\oplus\MorseCat{3}\oplus\MorseCat{3}$ \\
$V_{12}$ & 3 & $\MBcat{\Gamma}\oplus\bigoplus_{a:\,32a^2-32a-1=0}\MorseCat{3}$ \\
$V_{10}$ & 3 & $\mathcal C_{V_{10},-6}\oplus\bigoplus_{a:\,4a^2-2a-1=0}\MorseCat{3}$ \\
$V_8$ & 2 & $\mathcal C_{V_8,-8}\oplus \MorseCat{3}$ \\
$V_6$ & 2 & $\mathcal C_{V_6,-12}\oplus \MorseCat{3}$ \\
$V_4$ & 2 & $\Aper{S}{3}\oplus \MorseCat{3}$ \\
$V_2$ & 2 & $\Aper{S}{5}\oplus \MorseCat{3}$ \\
$MM_{2-1}$ & 4 & $\mathcal C_{MM_{2-1},-60}\oplus\MBcat{\Gamma}\oplus\bigoplus_{a:\,108a^2-1=0}\MorseCat{3}$ \\
$MM_{2-2}$ & 4 & $\MBcat{S}\oplus\bigoplus_{t:\,4t^3-16t^2+16t-1=0}\MorseCat{3}$ \\
$MM_{2-3}$ & 4 & $\mathcal C_{MM_{2-3},-12}\oplus\MBcat{\Gamma}\oplus\MorseCat{3}\oplus\MorseCat{3}$ \\
$MM_{2-4}$ & 4 & $\bigoplus_{i=1}^{4}\MorseCat{3}$ \\
$MM_{2-5}$ & 3 & $\MBcat{\Gamma} \oplus \MorseCat{3}\oplus \MorseCat{3}$ \\
$MM_{2-6}$ & 3 & $\MorseCat{3}\oplus \MBtwcat{\Gamma_1}{t}\oplus \MBcat{\Gamma_2}$ \\
$MM_{2-7}$ & 4 & $\MBcat{\Gamma}\oplus\bigoplus_{\alpha:\,4\alpha^3-4\alpha^2-4\alpha+3=0}\MorseCat{3}$ \\
$MM_{2-8}$ & 4 & $\MBcat{\Sigma}\oplus \bigoplus_{\rho:\,\rho^3-\rho-1=0}\MorseCat{3}$ \\
$MM_{2-9}$ & 4 & $\bigoplus_{\zeta:\,q(\zeta)=0}\MorseCat{3}$ \\
$MM_{2-10}$ & 2 & $\MorseCat{3}\oplus \MorseCat{3}$ \\
$MM_{2-11}$ & 4 & $\bigoplus_{\tau:\,q(\tau)=0}\MorseCat{3}$ \\
$MM_{2-12}$ & 2 & $\MBtwcat{\Gamma}{t}\oplus \MorseCat{3}$ \\
$MM_{2-13}$ & 4 & $\MBcat{\Gamma}\oplus \bigoplus_{\xi:\,\xi^3=16}\MorseCat{3}$ \\
$MM_{2-14}$ & 4 & $\bigoplus_{\eta:\,r(\eta)=0}\MorseCat{3}$ \\
$MM_{2-15}$ & 5 & $\MBcat{\Gamma}\oplus \bigoplus_{\alpha:\,q(\alpha)=0}\MorseCat{3}$ \\
$MM_{2-16}$ & 3 & $\bigoplus_{\beta:\,\beta^3-4\beta-16=0}\MorseCat{3}$ \\
$MM_{2-17}$ & 4 & $\bigoplus_{\eta:\,r(\eta)=0}\MorseCat{3}$ \\
$MM_{2-18}$ & 4 & $\bigoplus_{\eta:\,s(\eta)=0}\MorseCat{3}$ \\
$MM_{2-19}$ & 5 & $\MBcat{\Gamma} \oplus \bigoplus_{\rho:\,h(\rho)=0}\MorseCat{3}$ \\
$MM_{2-20}$ & 6 & $\bigoplus_{\alpha:\,q(\alpha)=0}\MorseCat{3}$ \\
$MM_{2-21}$ & 4 & $\bigoplus_{i=1}^{4}\MorseCat{3}$ \\
$MM_{2-22}$ & 5 & $\bigoplus_{\lambda:\,p(\lambda)=0}\MorseCat{3}$ \\
$MM_{2-23}$ & 3 & $\bigoplus_{\alpha:\,q(\alpha)=0}\MorseCat{3}$ \\
$MM_{2-24}$ & 6 & $\bigoplus_{\zeta^3=1}\left(\MorseCat{3}\oplus \MorseCat{3}\right)$ \\
$MM_{2-25}$ & 3 & $\bigoplus_{z:\,q(z)=0}\MorseCat{3}$ \\
$MM_{2-26}$ & 6 & $\bigoplus_{u:\,p(u)=0}\MorseCat{3}$ \\
$MM_{2-27}$ & 6 & $\bigoplus_{u:\,p(u)=0}\MorseCat{3}$ \\
$MM_{2-28}$ & 5 & $\MBcat{\Gamma}\oplus\bigoplus_{\lambda:\,r(\lambda)=0}\MorseCat{3}$ \\
$MM_{2-29}$ & 5 & $\MorseCat{3}\oplus\bigoplus_{\alpha:\,q(\alpha)=0}\MorseCat{3}$ \\
$MM_{2-30}$ & 5 & $\bigoplus_{\lambda:\,p(\lambda)=0}\MorseCat{3}$ \\
$MM_{2-31}$ & 6 & $\bigoplus_{\lambda:\,p(\lambda)=0}\MorseCat{3}$ \\
$MM_{2-32}$ & 6 & $\bigoplus_{\epsilon^2=8}\MorseCat{3} \oplus \bigoplus_{\substack{\omega^2+\omega+1=0\\ \delta^2+1=0}}\MorseCat{3}$ \\
$MM_{2-33}$ & 6 & $\bigoplus_{\lambda:\,p_+(\lambda)=0}\MorseCat{3} \oplus \bigoplus_{\mu:\,p_-(\mu)=0}\MorseCat{3}$ \\
$MM_{2-34}$ & 6 & $\bigoplus_{\substack{\epsilon^2=1\\ \rho^3=1}}\MorseCat{3}$ \\
$MM_{2-35}$ & 6 & $\bigoplus_{\lambda:\,\lambda^6+\lambda^4-1=0}\MorseCat{3}$ \\
$MM_{2-36}$ & 6 & $\bigoplus_{t:\,t^6+2t^5-1=0}\MorseCat{3}$ \\
$MM_{3-1}$ & 5 & $\MorseCat{3}\oplus \MorseCat{3} \oplus \MBtwcat{\Gamma_3}{v} \oplus\bigoplus_{\zeta^2+\zeta+1=0}\MBcat{\Gamma_\zeta}$ \\
$MM_{3-2}$ & 7 & $\MBcat{\Gamma} \oplus\bigoplus_{\lambda:\,p_+(\lambda)=0}\MorseCat{3} \oplus\bigoplus_{\mu:\,p_-(\mu)=0}\MorseCat{3}$ \\
$MM_{3-3}$ & 5 & $\bigoplus_{\alpha:\,q(\alpha)=0}\MorseCat{3} \oplus \bigoplus_{\varepsilon=\pm}\MorseCat{3}$ \\
$MM_{3-4}$ & 6 & $\MorseCat{3} \oplus \bigoplus_{\beta:\,p(\beta)=0}\MorseCat{3}$ \\
$MM_{3-5}$ & 8 & $\bigoplus_{\lambda:\,p_+(\lambda)=0}\MorseCat{3} \oplus \bigoplus_{\mu:\,p_-(\mu)=0}\MorseCat{3}$ \\
$MM_{3-6}$ & 6 & $\bigoplus_{\lambda:\,r_A(\lambda)=0}\MorseCat{3} \oplus \bigoplus_{\mu:\,r_B(\mu)=0}\MorseCat{3}$ \\
$MM_{3-7}$ & 2 & $\bigoplus_{\sigma:\,\sigma^2-6\sigma+7=0}\MorseCat{3}$ \\
$MM_{3-8}$ & 8 & $\MorseCat{3} \oplus \bigoplus_{\lambda:\,q(\lambda)=0}\MorseCat{3}$ \\
$MM_{3-9}$ & 5 & $\bigoplus_{\eta:\,\eta^3=16}\MorseCat{3} \oplus \bigoplus_{\alpha:\,\alpha^2+10\alpha+1=0}\MorseCat{3}$ \\
$MM_{3-10}$ & 6 & $\bigoplus_{\lambda:\,2\lambda^3-2\lambda-1=0}\MorseCat{3} \oplus \MorseCat{3} \oplus \bigoplus_{\alpha:\,\alpha^2+9\alpha+16=0}\MorseCat{3}$ \\
$MM_{3-11}$ & 6 & $\bigoplus_{\nu:\,\Phi(\nu)=0}\MorseCat{3}$ \\
$MM_{3-12}$ & 8 & $\bigoplus_{\lambda:\,\Phi_A(\lambda)=0}\MorseCat{3} \oplus \bigoplus_{\mu:\,\Phi_B(\mu)=0}\MorseCat{3}$ \\
$MM_{3-13}$ & 1 & $\MorseCat{3}$ \\
$MM_{3-14}$ & 6 & $\bigoplus_{\lambda:\,\Phi(\lambda)=0}\MorseCat{3}$ \\
$MM_{3-15}$ & 7 & $\bigoplus_{\lambda:\,\Phi(\lambda)=0}\MorseCat{3}$ \\
$MM_{3-16}$ & 8 & $\bigoplus_{\lambda:\,\Phi(\lambda)=0}\MorseCat{3}$ \\
$MM_{3-17}$ & 4 & $\bigoplus_{\lambda:\,\Phi(\lambda)=0}\MorseCat{3}$ \\
$MM_{3-18}$ & 7 & $\bigoplus_{\lambda:\,\Phi(\lambda)=0}\MorseCat{3}$ \\
$MM_{3-19}$ & 3 & $\bigoplus_{\lambda:\,\Phi(\lambda)=0}\MorseCat{3}$ \\
$MM_{3-20}$ & 4 & $\left(\bigoplus_{\epsilon=\pm1}\MorseCat{3}\right) \oplus \bigoplus_{\mu:\,\Psi(\mu)=0}\MorseCat{3}$ \\
$MM_{3-21}$ & 8 & $\bigoplus_{\tau:\,\Phi(\tau)=0}\MorseCat{3}$ \\
$MM_{3-22}$ & 7 & $\bigoplus_{\lambda:\,\Phi(\lambda)=0}\MorseCat{3}$ \\
$MM_{3-23}$ & 8 & $\bigoplus_{\lambda:\,\Psi(\lambda)=0}\MorseCat{3}$ \\
$MM_{3-24}$ & 8 & $\bigoplus_{\tau:\,\Phi(\tau)=0}\MorseCat{3}$ \\
$MM_{3-25}$ & 8 & $\left(\bigoplus_{\alpha:\,q_+(\alpha)=0}\MorseCat{3}\right) \oplus \left(\bigoplus_{\beta:\,q_-(\beta)=0}\MorseCat{3}\right)$ \\
$MM_{3-26}$ & 8 & $\bigoplus_{\lambda:\,\Psi(\lambda)=0}\MorseCat{3}$ \\
$MM_{3-27}$ & 8 & $\bigoplus_{\epsilon\in\{1,-1\}^3}\MorseCat{3}$ \\
$MM_{3-28}$ & 8 & $\bigoplus_{\epsilon\in\{1,-1\}} \left(\bigoplus_{\lambda:\,\Phi(\lambda)=0}\MorseCat{3}\right)$ \\
$MM_{3-29}$ & 8 & $\bigoplus_{\lambda:\,\Phi(\lambda)=0}\MorseCat{3}$ \\
$MM_{3-30}$ & 8 & $\bigoplus_{\lambda:\,\Phi(\lambda)=0}\MorseCat{3}$ \\
$MM_{3-31}$ & 4 & $\bigoplus_{\lambda:\,\Phi(\lambda)=0}\MorseCat{3}$ \\
$MM_{4-1}$ & 2 & $\MBtwcat{\Gamma}{t(t+1)}\oplus \MorseCat{3}$ \\
$MM_{4-2}$ & 6 & $\left(\bigoplus_{\alpha:\,p_+(\alpha)=0}\MorseCat{3}\right) \oplus \left(\bigoplus_{\beta:\,p_-(\beta)=0}\MorseCat{3}\right)$ \\
$MM_{4-3}$ & 2 & $\MBtwcat{\Gamma}{1-t}\oplus \MorseCat{3}$ \\
$MM_{4-4}$ & 5 & $\bigoplus_{\lambda:\,\Phi(\lambda)=0}\MorseCat{3}$ \\
$MM_{4-5}$ & 5 & $\MorseCat{3}\oplus \bigoplus_{\alpha:\,q(\alpha)=0}\MorseCat{3}$ \\
$MM_{4-6}$ & 10 & $\bigoplus_{\tau:\,P(\tau)=0}\MorseCat{3}$ \\
$MM_{4-7}$ & 5 & $\left(\bigoplus_{\alpha:\,p_+(\alpha)=0}\MorseCat{3}\right) \oplus \left(\bigoplus_{\beta:\,p_-(\beta)=0}\MorseCat{3}\right)$ \\
$MM_{4-8}$ & 10 & $\left(\bigoplus_{\sigma^2=2}\MorseCat{3}\right) \oplus \left(\bigoplus_{\tau:\,R(\tau)=0}\MorseCat{3}\right)$ \\
$MM_{4-9}$ & 5 & $\bigoplus_{\lambda:\,\Phi(\lambda)=0}\MorseCat{3}$ \\
$MM_{4-10}$ & 10 & $\bigoplus_{\gamma:\,\Psi(\gamma)=0}\MorseCat{3}$ \\
$MM_{4-11}$ & 6 & $\bigoplus_{\varepsilon\in\{\pm1\}} \bigoplus_{\alpha:\,p(\alpha)=0}\MorseCat{3}$ \\
$MM_{4-12}$ & 10 & $\bigoplus_{\delta:\,\Xi(\delta)=0}\MorseCat{3}$ \\
$MM_{4-13}$ & 10 & $\bigoplus_{p\in\CritLoc(W_{MM_{4-13}})}\MorseCat{3}$ \\
$MM_{5-1}$ & 5 & $\bigoplus_{\alpha:\,p(\alpha)=0}\MorseCat{3}$ \\
$MM_{5-2}$ & 8 & $\left(\bigoplus_{\alpha:\,p_-(\alpha)=0}\MorseCat{3}\right) \oplus \left(\bigoplus_{\beta:\,p_+(\beta)=0}\MorseCat{3}\right)$ \\
$MM_{5-3}$ & 2 & $\bigoplus_{\varepsilon\in\{\pm1\}}\MorseCat{3}$ \\
$MM_{6-1}$ & 4 & $\bigoplus_{\varepsilon\in\{\pm1\}} \bigoplus_{\alpha:\,q(\alpha)=0}\MorseCat{3}$ \\
$MM_{7-1}$ & 4 & $\MBtwcat{\Gamma_1}{t}\oplus \MBtwcat{\Gamma_{-1}}{t} \oplus \bigoplus_{\varepsilon\in\{\pm1\}}\MorseCat{3}$ \\
$MM_{8-1}$ & 4 & $\QAper{\Gamma_1}{2}\oplus\QAper{\Gamma_{-1}}{2}\oplus\bigoplus_{\varepsilon\in\{\pm1\}}\MorseCat{3}$ \\
$MM_{9-1}$ & 4 & $\displaystyle\bigoplus_{\varepsilon\in\{\pm1\}}\QAtwist{\Gamma_\varepsilon}{3}{t(t+1)}\oplus\bigoplus_{\varepsilon\in\{\pm1\}}\MorseCat{3}$ \\
$MM_{10-1}$ & 4 & $\displaystyle\bigoplus_{\varepsilon\in\{\pm1\}}\QAtwist{\Gamma_\varepsilon}{5}{t+1}\oplus\bigoplus_{\varepsilon\in\{\pm1\}}\MorseCat{3}$ \\
\end{longtblr}
\endgroup

\bibliographystyle{amsplain}
\bibliography{ref}

\end{document}
